% GUID v6: the ten primer boxes, one file each in primers/, each a complete primer environment (preamble.tex) with its label.
% Place each at first use, preferably by a direct input at the %PRIMER markers of part1.tex and in the sections named below, e.g.
%   % GUID v6, Primer P1 (label box:dmp): DMP and trust systems, with a worked act. Background only; it states no result of the paper.
% Place it at the start of Section sec:dmp (marker %PRIMER:dmp in part1.tex): % GUID v6, Primer P1 (label box:dmp): DMP and trust systems, with a worked act. Background only; it states no result of the paper.
% Place it at the start of Section sec:dmp (marker %PRIMER:dmp in part1.tex): % GUID v6, Primer P1 (label box:dmp): DMP and trust systems, with a worked act. Background only; it states no result of the paper.
% Place it at the start of Section sec:dmp (marker %PRIMER:dmp in part1.tex): \input{primers/p1-dmp}
\begin{primer}{DMP and trust systems}{box:dmp}
\emph{What a trust system is for.} In a computer system an access-control or trust protocol decides, for every request, who may do what
to which information, and it keeps the evidence of what was done. Its classic design rules are to grant the least privilege an act needs,
to demand two independent keys for what matters (separation of privilege), and to keep a complete audit trail~\cite{saltzer1975}. Its
formal tools include lattices that say where information may flow~\cite{denning}, noninterference, which says what one side may learn of
the other~\cite{gm}, and quorums of witnesses that fail independently~\cite{mr}. DMP~\cite{dmp} states such a system as three axioms and a
ground, G0; every later statement is a numbered theorem (T) or derived constraint (C).

\emph{Words.} A \emph{part} is ``a Duality of Trust'': a trust half and an information half. An \emph{act} is what a part does. A
\emph{layer} is where trust over acts is decided; its two \emph{witnesses} are its halves. The \emph{trust half} decides whether the act
may happen, the \emph{information half} whether it is meant and right. Two witnesses are \emph{independent} when they share no member of the
\emph{failure classes} the layer declares (host, key, author, engine, person, time source, and whatever the minting lists). A
\emph{standing} is trust held in the field of Axiom~2. It has a source (the witness that granted it), a depth (its \emph{grant chain} back
to the first pair) and a freshness that decays with time; \emph{staleness} is that decay. A \emph{record} is evidence that outlives its act.

\emph{The gate.} Each witness returns a verdict in the lattice $\textsf{refuse}<\textsf{hold}<\textsf{pass}$, and the gate of layer $L$
is their meet, $\operatorname{gate}_L(a,t)=v_{w_T}(a,t)\wedge v_{w_I}(a,t)$. The stricter answer stands, and the act is performed only on
$\textsf{pass}\wedge\textsf{pass}$.

\emph{The axioms in one line each.}
\begin{itemize}[nosep,leftmargin=*]
\item \emph{Axiom~1, the Duality.} A part acts only while both halves of its layer trust it for that act; neither overrules the other. Its
\emph{balance}: every part works towards the trust its boundary declares.
\item \emph{Axiom~2, Trust.} A witness trusts only on evidence it verified itself. ``Trust is gravitational. It is a field, never a
position.'' A strong enough boundary ``bends trustspace'' so that only entropy leaves.
\item \emph{Axiom~3, Information.} Every grant, use and loss of trust leaves a record that every witness of the act reads in full. Any other
reader sees the \emph{counts} always and the content only by a \emph{grant}; from trust, nothing about information can be derived but its
existence. Every layer is a tree rooted at its own Duality. The \emph{layout}, where an item sits in that tree, is ``the median of that
demand'', the place of fewest hops in total, recomputed by both halves and never by one.
\end{itemize}

\emph{A worked act} (Figure~\ref{fig:act}). A part $p$ asks to change one entry of a directory. The trust half checks $p$'s standing: it
was granted by a witness, lies two grants from the minted pair and is still fresh, so the verdict is \textsf{pass}. The information half
finds the change well formed but not yet confirmed: \textsf{hold}. The gate gives $\textsf{pass}\wedge\textsf{hold}=\textsf{hold}$, and
nothing is done. Once the change is confirmed, both return \textsf{pass} and the act is performed. A record is written at the act, and both
witnesses can read it in full. A reader outside the layer, say an auditor, sees the counts (here, one more act of this layer) but not the
new entry, unless it is granted to him. From the trust side alone one learns only that the entry exists.

\emph{Containment and G0.} A boundary \emph{contains} information when no part beyond it is trusted for any act on that information.
G0: ``The dual gate cannot witness its own first instance: somebody has to trust the first two witnesses.'' The first act, the
\emph{minting}, makes at least five parts the holders of the \emph{root keys}; it is stated, not derived, and every later standing
descends from it through grant chains.

\emph{The theorems this paper cites.}
\begin{itemize}[nosep,leftmargin=*]
\item \emph{T4}: every parameter is configuration. DMP fixes structure, never a value, and so derives no constant.
\item \emph{T27}: a contained item is destroyed at its scheduled date; with T31 it describes the trap.
\item \emph{T31}: contained information reaches no part by any act. What still leaves is minimal, thermal emission, the \emph{entropy
floor}, which ``carries entropy and no information''. Outside, one learns only the size, the staleness and the rotation interval.
\item \emph{T32}: the VIBE field. Each member has a depth, a mass (one plus the members deriving through it) and a shortfall; its
potential $\Phi_{\mathrm V}\ge0$ never rises under the system's own acts and vanishes in the \emph{balanced state}. Readers other than the
witnesses see one number, the boundary's temperature, its mass-weighted staleness.
\item \emph{T37}: information that no admission names is contained at the item itself.
\item \emph{T40}: a minting originates only from G0 together with a state in which every axiom and theorem provably holds.
\end{itemize}
\end{primer}

\begin{primer}{DMP and trust systems}{box:dmp}
\emph{What a trust system is for.} In a computer system an access-control or trust protocol decides, for every request, who may do what
to which information, and it keeps the evidence of what was done. Its classic design rules are to grant the least privilege an act needs,
to demand two independent keys for what matters (separation of privilege), and to keep a complete audit trail~\cite{saltzer1975}. Its
formal tools include lattices that say where information may flow~\cite{denning}, noninterference, which says what one side may learn of
the other~\cite{gm}, and quorums of witnesses that fail independently~\cite{mr}. DMP~\cite{dmp} states such a system as three axioms and a
ground, G0; every later statement is a numbered theorem (T) or derived constraint (C).

\emph{Words.} A \emph{part} is ``a Duality of Trust'': a trust half and an information half. An \emph{act} is what a part does. A
\emph{layer} is where trust over acts is decided; its two \emph{witnesses} are its halves. The \emph{trust half} decides whether the act
may happen, the \emph{information half} whether it is meant and right. Two witnesses are \emph{independent} when they share no member of the
\emph{failure classes} the layer declares (host, key, author, engine, person, time source, and whatever the minting lists). A
\emph{standing} is trust held in the field of Axiom~2. It has a source (the witness that granted it), a depth (its \emph{grant chain} back
to the first pair) and a freshness that decays with time; \emph{staleness} is that decay. A \emph{record} is evidence that outlives its act.

\emph{The gate.} Each witness returns a verdict in the lattice $\textsf{refuse}<\textsf{hold}<\textsf{pass}$, and the gate of layer $L$
is their meet, $\operatorname{gate}_L(a,t)=v_{w_T}(a,t)\wedge v_{w_I}(a,t)$. The stricter answer stands, and the act is performed only on
$\textsf{pass}\wedge\textsf{pass}$.

\emph{The axioms in one line each.}
\begin{itemize}[nosep,leftmargin=*]
\item \emph{Axiom~1, the Duality.} A part acts only while both halves of its layer trust it for that act; neither overrules the other. Its
\emph{balance}: every part works towards the trust its boundary declares.
\item \emph{Axiom~2, Trust.} A witness trusts only on evidence it verified itself. ``Trust is gravitational. It is a field, never a
position.'' A strong enough boundary ``bends trustspace'' so that only entropy leaves.
\item \emph{Axiom~3, Information.} Every grant, use and loss of trust leaves a record that every witness of the act reads in full. Any other
reader sees the \emph{counts} always and the content only by a \emph{grant}; from trust, nothing about information can be derived but its
existence. Every layer is a tree rooted at its own Duality. The \emph{layout}, where an item sits in that tree, is ``the median of that
demand'', the place of fewest hops in total, recomputed by both halves and never by one.
\end{itemize}

\emph{A worked act} (Figure~\ref{fig:act}). A part $p$ asks to change one entry of a directory. The trust half checks $p$'s standing: it
was granted by a witness, lies two grants from the minted pair and is still fresh, so the verdict is \textsf{pass}. The information half
finds the change well formed but not yet confirmed: \textsf{hold}. The gate gives $\textsf{pass}\wedge\textsf{hold}=\textsf{hold}$, and
nothing is done. Once the change is confirmed, both return \textsf{pass} and the act is performed. A record is written at the act, and both
witnesses can read it in full. A reader outside the layer, say an auditor, sees the counts (here, one more act of this layer) but not the
new entry, unless it is granted to him. From the trust side alone one learns only that the entry exists.

\emph{Containment and G0.} A boundary \emph{contains} information when no part beyond it is trusted for any act on that information.
G0: ``The dual gate cannot witness its own first instance: somebody has to trust the first two witnesses.'' The first act, the
\emph{minting}, makes at least five parts the holders of the \emph{root keys}; it is stated, not derived, and every later standing
descends from it through grant chains.

\emph{The theorems this paper cites.}
\begin{itemize}[nosep,leftmargin=*]
\item \emph{T4}: every parameter is configuration. DMP fixes structure, never a value, and so derives no constant.
\item \emph{T27}: a contained item is destroyed at its scheduled date; with T31 it describes the trap.
\item \emph{T31}: contained information reaches no part by any act. What still leaves is minimal, thermal emission, the \emph{entropy
floor}, which ``carries entropy and no information''. Outside, one learns only the size, the staleness and the rotation interval.
\item \emph{T32}: the VIBE field. Each member has a depth, a mass (one plus the members deriving through it) and a shortfall; its
potential $\Phi_{\mathrm V}\ge0$ never rises under the system's own acts and vanishes in the \emph{balanced state}. Readers other than the
witnesses see one number, the boundary's temperature, its mass-weighted staleness.
\item \emph{T37}: information that no admission names is contained at the item itself.
\item \emph{T40}: a minting originates only from G0 together with a state in which every axiom and theorem provably holds.
\end{itemize}
\end{primer}

\begin{primer}{DMP and trust systems}{box:dmp}
\emph{What a trust system is for.} In a computer system an access-control or trust protocol decides, for every request, who may do what
to which information, and it keeps the evidence of what was done. Its classic design rules are to grant the least privilege an act needs,
to demand two independent keys for what matters (separation of privilege), and to keep a complete audit trail~\cite{saltzer1975}. Its
formal tools include lattices that say where information may flow~\cite{denning}, noninterference, which says what one side may learn of
the other~\cite{gm}, and quorums of witnesses that fail independently~\cite{mr}. DMP~\cite{dmp} states such a system as three axioms and a
ground, G0; every later statement is a numbered theorem (T) or derived constraint (C).

\emph{Words.} A \emph{part} is ``a Duality of Trust'': a trust half and an information half. An \emph{act} is what a part does. A
\emph{layer} is where trust over acts is decided; its two \emph{witnesses} are its halves. The \emph{trust half} decides whether the act
may happen, the \emph{information half} whether it is meant and right. Two witnesses are \emph{independent} when they share no member of the
\emph{failure classes} the layer declares (host, key, author, engine, person, time source, and whatever the minting lists). A
\emph{standing} is trust held in the field of Axiom~2. It has a source (the witness that granted it), a depth (its \emph{grant chain} back
to the first pair) and a freshness that decays with time; \emph{staleness} is that decay. A \emph{record} is evidence that outlives its act.

\emph{The gate.} Each witness returns a verdict in the lattice $\textsf{refuse}<\textsf{hold}<\textsf{pass}$, and the gate of layer $L$
is their meet, $\operatorname{gate}_L(a,t)=v_{w_T}(a,t)\wedge v_{w_I}(a,t)$. The stricter answer stands, and the act is performed only on
$\textsf{pass}\wedge\textsf{pass}$.

\emph{The axioms in one line each.}
\begin{itemize}[nosep,leftmargin=*]
\item \emph{Axiom~1, the Duality.} A part acts only while both halves of its layer trust it for that act; neither overrules the other. Its
\emph{balance}: every part works towards the trust its boundary declares.
\item \emph{Axiom~2, Trust.} A witness trusts only on evidence it verified itself. ``Trust is gravitational. It is a field, never a
position.'' A strong enough boundary ``bends trustspace'' so that only entropy leaves.
\item \emph{Axiom~3, Information.} Every grant, use and loss of trust leaves a record that every witness of the act reads in full. Any other
reader sees the \emph{counts} always and the content only by a \emph{grant}; from trust, nothing about information can be derived but its
existence. Every layer is a tree rooted at its own Duality. The \emph{layout}, where an item sits in that tree, is ``the median of that
demand'', the place of fewest hops in total, recomputed by both halves and never by one.
\end{itemize}

\emph{A worked act} (Figure~\ref{fig:act}). A part $p$ asks to change one entry of a directory. The trust half checks $p$'s standing: it
was granted by a witness, lies two grants from the minted pair and is still fresh, so the verdict is \textsf{pass}. The information half
finds the change well formed but not yet confirmed: \textsf{hold}. The gate gives $\textsf{pass}\wedge\textsf{hold}=\textsf{hold}$, and
nothing is done. Once the change is confirmed, both return \textsf{pass} and the act is performed. A record is written at the act, and both
witnesses can read it in full. A reader outside the layer, say an auditor, sees the counts (here, one more act of this layer) but not the
new entry, unless it is granted to him. From the trust side alone one learns only that the entry exists.

\emph{Containment and G0.} A boundary \emph{contains} information when no part beyond it is trusted for any act on that information.
G0: ``The dual gate cannot witness its own first instance: somebody has to trust the first two witnesses.'' The first act, the
\emph{minting}, makes at least five parts the holders of the \emph{root keys}; it is stated, not derived, and every later standing
descends from it through grant chains.

\emph{The theorems this paper cites.}
\begin{itemize}[nosep,leftmargin=*]
\item \emph{T4}: every parameter is configuration. DMP fixes structure, never a value, and so derives no constant.
\item \emph{T27}: a contained item is destroyed at its scheduled date; with T31 it describes the trap.
\item \emph{T31}: contained information reaches no part by any act. What still leaves is minimal, thermal emission, the \emph{entropy
floor}, which ``carries entropy and no information''. Outside, one learns only the size, the staleness and the rotation interval.
\item \emph{T32}: the VIBE field. Each member has a depth, a mass (one plus the members deriving through it) and a shortfall; its
potential $\Phi_{\mathrm V}\ge0$ never rises under the system's own acts and vanishes in the \emph{balanced state}. Readers other than the
witnesses see one number, the boundary's temperature, its mass-weighted staleness.
\item \emph{T37}: information that no admission names is contained at the item itself.
\item \emph{T40}: a minting originates only from G0 together with a state in which every axiom and theorem provably holds.
\end{itemize}
\end{primer}

% A direct input keeps build_bib.py's citation order at the place of first use. Alternatively, % GUID v6: the ten primer boxes, one file each in primers/, each a complete primer environment (preamble.tex) with its label.
% Place each at first use, preferably by a direct input at the %PRIMER markers of part1.tex and in the sections named below, e.g.
%   % GUID v6, Primer P1 (label box:dmp): DMP and trust systems, with a worked act. Background only; it states no result of the paper.
% Place it at the start of Section sec:dmp (marker %PRIMER:dmp in part1.tex): % GUID v6, Primer P1 (label box:dmp): DMP and trust systems, with a worked act. Background only; it states no result of the paper.
% Place it at the start of Section sec:dmp (marker %PRIMER:dmp in part1.tex): \input{primers/p1-dmp}
\begin{primer}{DMP and trust systems}{box:dmp}
\emph{What a trust system is for.} In a computer system an access-control or trust protocol decides, for every request, who may do what
to which information, and it keeps the evidence of what was done. Its classic design rules are to grant the least privilege an act needs,
to demand two independent keys for what matters (separation of privilege), and to keep a complete audit trail~\cite{saltzer1975}. Its
formal tools include lattices that say where information may flow~\cite{denning}, noninterference, which says what one side may learn of
the other~\cite{gm}, and quorums of witnesses that fail independently~\cite{mr}. DMP~\cite{dmp} states such a system as three axioms and a
ground, G0; every later statement is a numbered theorem (T) or derived constraint (C).

\emph{Words.} A \emph{part} is ``a Duality of Trust'': a trust half and an information half. An \emph{act} is what a part does. A
\emph{layer} is where trust over acts is decided; its two \emph{witnesses} are its halves. The \emph{trust half} decides whether the act
may happen, the \emph{information half} whether it is meant and right. Two witnesses are \emph{independent} when they share no member of the
\emph{failure classes} the layer declares (host, key, author, engine, person, time source, and whatever the minting lists). A
\emph{standing} is trust held in the field of Axiom~2. It has a source (the witness that granted it), a depth (its \emph{grant chain} back
to the first pair) and a freshness that decays with time; \emph{staleness} is that decay. A \emph{record} is evidence that outlives its act.

\emph{The gate.} Each witness returns a verdict in the lattice $\textsf{refuse}<\textsf{hold}<\textsf{pass}$, and the gate of layer $L$
is their meet, $\operatorname{gate}_L(a,t)=v_{w_T}(a,t)\wedge v_{w_I}(a,t)$. The stricter answer stands, and the act is performed only on
$\textsf{pass}\wedge\textsf{pass}$.

\emph{The axioms in one line each.}
\begin{itemize}[nosep,leftmargin=*]
\item \emph{Axiom~1, the Duality.} A part acts only while both halves of its layer trust it for that act; neither overrules the other. Its
\emph{balance}: every part works towards the trust its boundary declares.
\item \emph{Axiom~2, Trust.} A witness trusts only on evidence it verified itself. ``Trust is gravitational. It is a field, never a
position.'' A strong enough boundary ``bends trustspace'' so that only entropy leaves.
\item \emph{Axiom~3, Information.} Every grant, use and loss of trust leaves a record that every witness of the act reads in full. Any other
reader sees the \emph{counts} always and the content only by a \emph{grant}; from trust, nothing about information can be derived but its
existence. Every layer is a tree rooted at its own Duality. The \emph{layout}, where an item sits in that tree, is ``the median of that
demand'', the place of fewest hops in total, recomputed by both halves and never by one.
\end{itemize}

\emph{A worked act} (Figure~\ref{fig:act}). A part $p$ asks to change one entry of a directory. The trust half checks $p$'s standing: it
was granted by a witness, lies two grants from the minted pair and is still fresh, so the verdict is \textsf{pass}. The information half
finds the change well formed but not yet confirmed: \textsf{hold}. The gate gives $\textsf{pass}\wedge\textsf{hold}=\textsf{hold}$, and
nothing is done. Once the change is confirmed, both return \textsf{pass} and the act is performed. A record is written at the act, and both
witnesses can read it in full. A reader outside the layer, say an auditor, sees the counts (here, one more act of this layer) but not the
new entry, unless it is granted to him. From the trust side alone one learns only that the entry exists.

\emph{Containment and G0.} A boundary \emph{contains} information when no part beyond it is trusted for any act on that information.
G0: ``The dual gate cannot witness its own first instance: somebody has to trust the first two witnesses.'' The first act, the
\emph{minting}, makes at least five parts the holders of the \emph{root keys}; it is stated, not derived, and every later standing
descends from it through grant chains.

\emph{The theorems this paper cites.}
\begin{itemize}[nosep,leftmargin=*]
\item \emph{T4}: every parameter is configuration. DMP fixes structure, never a value, and so derives no constant.
\item \emph{T27}: a contained item is destroyed at its scheduled date; with T31 it describes the trap.
\item \emph{T31}: contained information reaches no part by any act. What still leaves is minimal, thermal emission, the \emph{entropy
floor}, which ``carries entropy and no information''. Outside, one learns only the size, the staleness and the rotation interval.
\item \emph{T32}: the VIBE field. Each member has a depth, a mass (one plus the members deriving through it) and a shortfall; its
potential $\Phi_{\mathrm V}\ge0$ never rises under the system's own acts and vanishes in the \emph{balanced state}. Readers other than the
witnesses see one number, the boundary's temperature, its mass-weighted staleness.
\item \emph{T37}: information that no admission names is contained at the item itself.
\item \emph{T40}: a minting originates only from G0 together with a state in which every axiom and theorem provably holds.
\end{itemize}
\end{primer}

\begin{primer}{DMP and trust systems}{box:dmp}
\emph{What a trust system is for.} In a computer system an access-control or trust protocol decides, for every request, who may do what
to which information, and it keeps the evidence of what was done. Its classic design rules are to grant the least privilege an act needs,
to demand two independent keys for what matters (separation of privilege), and to keep a complete audit trail~\cite{saltzer1975}. Its
formal tools include lattices that say where information may flow~\cite{denning}, noninterference, which says what one side may learn of
the other~\cite{gm}, and quorums of witnesses that fail independently~\cite{mr}. DMP~\cite{dmp} states such a system as three axioms and a
ground, G0; every later statement is a numbered theorem (T) or derived constraint (C).

\emph{Words.} A \emph{part} is ``a Duality of Trust'': a trust half and an information half. An \emph{act} is what a part does. A
\emph{layer} is where trust over acts is decided; its two \emph{witnesses} are its halves. The \emph{trust half} decides whether the act
may happen, the \emph{information half} whether it is meant and right. Two witnesses are \emph{independent} when they share no member of the
\emph{failure classes} the layer declares (host, key, author, engine, person, time source, and whatever the minting lists). A
\emph{standing} is trust held in the field of Axiom~2. It has a source (the witness that granted it), a depth (its \emph{grant chain} back
to the first pair) and a freshness that decays with time; \emph{staleness} is that decay. A \emph{record} is evidence that outlives its act.

\emph{The gate.} Each witness returns a verdict in the lattice $\textsf{refuse}<\textsf{hold}<\textsf{pass}$, and the gate of layer $L$
is their meet, $\operatorname{gate}_L(a,t)=v_{w_T}(a,t)\wedge v_{w_I}(a,t)$. The stricter answer stands, and the act is performed only on
$\textsf{pass}\wedge\textsf{pass}$.

\emph{The axioms in one line each.}
\begin{itemize}[nosep,leftmargin=*]
\item \emph{Axiom~1, the Duality.} A part acts only while both halves of its layer trust it for that act; neither overrules the other. Its
\emph{balance}: every part works towards the trust its boundary declares.
\item \emph{Axiom~2, Trust.} A witness trusts only on evidence it verified itself. ``Trust is gravitational. It is a field, never a
position.'' A strong enough boundary ``bends trustspace'' so that only entropy leaves.
\item \emph{Axiom~3, Information.} Every grant, use and loss of trust leaves a record that every witness of the act reads in full. Any other
reader sees the \emph{counts} always and the content only by a \emph{grant}; from trust, nothing about information can be derived but its
existence. Every layer is a tree rooted at its own Duality. The \emph{layout}, where an item sits in that tree, is ``the median of that
demand'', the place of fewest hops in total, recomputed by both halves and never by one.
\end{itemize}

\emph{A worked act} (Figure~\ref{fig:act}). A part $p$ asks to change one entry of a directory. The trust half checks $p$'s standing: it
was granted by a witness, lies two grants from the minted pair and is still fresh, so the verdict is \textsf{pass}. The information half
finds the change well formed but not yet confirmed: \textsf{hold}. The gate gives $\textsf{pass}\wedge\textsf{hold}=\textsf{hold}$, and
nothing is done. Once the change is confirmed, both return \textsf{pass} and the act is performed. A record is written at the act, and both
witnesses can read it in full. A reader outside the layer, say an auditor, sees the counts (here, one more act of this layer) but not the
new entry, unless it is granted to him. From the trust side alone one learns only that the entry exists.

\emph{Containment and G0.} A boundary \emph{contains} information when no part beyond it is trusted for any act on that information.
G0: ``The dual gate cannot witness its own first instance: somebody has to trust the first two witnesses.'' The first act, the
\emph{minting}, makes at least five parts the holders of the \emph{root keys}; it is stated, not derived, and every later standing
descends from it through grant chains.

\emph{The theorems this paper cites.}
\begin{itemize}[nosep,leftmargin=*]
\item \emph{T4}: every parameter is configuration. DMP fixes structure, never a value, and so derives no constant.
\item \emph{T27}: a contained item is destroyed at its scheduled date; with T31 it describes the trap.
\item \emph{T31}: contained information reaches no part by any act. What still leaves is minimal, thermal emission, the \emph{entropy
floor}, which ``carries entropy and no information''. Outside, one learns only the size, the staleness and the rotation interval.
\item \emph{T32}: the VIBE field. Each member has a depth, a mass (one plus the members deriving through it) and a shortfall; its
potential $\Phi_{\mathrm V}\ge0$ never rises under the system's own acts and vanishes in the \emph{balanced state}. Readers other than the
witnesses see one number, the boundary's temperature, its mass-weighted staleness.
\item \emph{T37}: information that no admission names is contained at the item itself.
\item \emph{T40}: a minting originates only from G0 together with a state in which every axiom and theorem provably holds.
\end{itemize}
\end{primer}

% A direct input keeps build_bib.py's citation order at the place of first use. Alternatively, % GUID v6: the ten primer boxes, one file each in primers/, each a complete primer environment (preamble.tex) with its label.
% Place each at first use, preferably by a direct input at the %PRIMER markers of part1.tex and in the sections named below, e.g.
%   % GUID v6, Primer P1 (label box:dmp): DMP and trust systems, with a worked act. Background only; it states no result of the paper.
% Place it at the start of Section sec:dmp (marker %PRIMER:dmp in part1.tex): \input{primers/p1-dmp}
\begin{primer}{DMP and trust systems}{box:dmp}
\emph{What a trust system is for.} In a computer system an access-control or trust protocol decides, for every request, who may do what
to which information, and it keeps the evidence of what was done. Its classic design rules are to grant the least privilege an act needs,
to demand two independent keys for what matters (separation of privilege), and to keep a complete audit trail~\cite{saltzer1975}. Its
formal tools include lattices that say where information may flow~\cite{denning}, noninterference, which says what one side may learn of
the other~\cite{gm}, and quorums of witnesses that fail independently~\cite{mr}. DMP~\cite{dmp} states such a system as three axioms and a
ground, G0; every later statement is a numbered theorem (T) or derived constraint (C).

\emph{Words.} A \emph{part} is ``a Duality of Trust'': a trust half and an information half. An \emph{act} is what a part does. A
\emph{layer} is where trust over acts is decided; its two \emph{witnesses} are its halves. The \emph{trust half} decides whether the act
may happen, the \emph{information half} whether it is meant and right. Two witnesses are \emph{independent} when they share no member of the
\emph{failure classes} the layer declares (host, key, author, engine, person, time source, and whatever the minting lists). A
\emph{standing} is trust held in the field of Axiom~2. It has a source (the witness that granted it), a depth (its \emph{grant chain} back
to the first pair) and a freshness that decays with time; \emph{staleness} is that decay. A \emph{record} is evidence that outlives its act.

\emph{The gate.} Each witness returns a verdict in the lattice $\textsf{refuse}<\textsf{hold}<\textsf{pass}$, and the gate of layer $L$
is their meet, $\operatorname{gate}_L(a,t)=v_{w_T}(a,t)\wedge v_{w_I}(a,t)$. The stricter answer stands, and the act is performed only on
$\textsf{pass}\wedge\textsf{pass}$.

\emph{The axioms in one line each.}
\begin{itemize}[nosep,leftmargin=*]
\item \emph{Axiom~1, the Duality.} A part acts only while both halves of its layer trust it for that act; neither overrules the other. Its
\emph{balance}: every part works towards the trust its boundary declares.
\item \emph{Axiom~2, Trust.} A witness trusts only on evidence it verified itself. ``Trust is gravitational. It is a field, never a
position.'' A strong enough boundary ``bends trustspace'' so that only entropy leaves.
\item \emph{Axiom~3, Information.} Every grant, use and loss of trust leaves a record that every witness of the act reads in full. Any other
reader sees the \emph{counts} always and the content only by a \emph{grant}; from trust, nothing about information can be derived but its
existence. Every layer is a tree rooted at its own Duality. The \emph{layout}, where an item sits in that tree, is ``the median of that
demand'', the place of fewest hops in total, recomputed by both halves and never by one.
\end{itemize}

\emph{A worked act} (Figure~\ref{fig:act}). A part $p$ asks to change one entry of a directory. The trust half checks $p$'s standing: it
was granted by a witness, lies two grants from the minted pair and is still fresh, so the verdict is \textsf{pass}. The information half
finds the change well formed but not yet confirmed: \textsf{hold}. The gate gives $\textsf{pass}\wedge\textsf{hold}=\textsf{hold}$, and
nothing is done. Once the change is confirmed, both return \textsf{pass} and the act is performed. A record is written at the act, and both
witnesses can read it in full. A reader outside the layer, say an auditor, sees the counts (here, one more act of this layer) but not the
new entry, unless it is granted to him. From the trust side alone one learns only that the entry exists.

\emph{Containment and G0.} A boundary \emph{contains} information when no part beyond it is trusted for any act on that information.
G0: ``The dual gate cannot witness its own first instance: somebody has to trust the first two witnesses.'' The first act, the
\emph{minting}, makes at least five parts the holders of the \emph{root keys}; it is stated, not derived, and every later standing
descends from it through grant chains.

\emph{The theorems this paper cites.}
\begin{itemize}[nosep,leftmargin=*]
\item \emph{T4}: every parameter is configuration. DMP fixes structure, never a value, and so derives no constant.
\item \emph{T27}: a contained item is destroyed at its scheduled date; with T31 it describes the trap.
\item \emph{T31}: contained information reaches no part by any act. What still leaves is minimal, thermal emission, the \emph{entropy
floor}, which ``carries entropy and no information''. Outside, one learns only the size, the staleness and the rotation interval.
\item \emph{T32}: the VIBE field. Each member has a depth, a mass (one plus the members deriving through it) and a shortfall; its
potential $\Phi_{\mathrm V}\ge0$ never rises under the system's own acts and vanishes in the \emph{balanced state}. Readers other than the
witnesses see one number, the boundary's temperature, its mass-weighted staleness.
\item \emph{T37}: information that no admission names is contained at the item itself.
\item \emph{T40}: a minting originates only from G0 together with a state in which every axiom and theorem provably holds.
\end{itemize}
\end{primer}

% A direct input keeps build_bib.py's citation order at the place of first use. Alternatively, % GUID v6: the ten primer boxes, one file each in primers/, each a complete primer environment (preamble.tex) with its label.
% Place each at first use, preferably by a direct input at the %PRIMER markers of part1.tex and in the sections named below, e.g.
%   \input{primers/p1-dmp}
% A direct input keeps build_bib.py's citation order at the place of first use. Alternatively, \input{primers} once after
% \begin{document} defines the macros below, and \primerDMP and so on place the boxes (build_bib.py then numbers the primers'
% references from the position of that \input).
%
%   primer  label         file                      first use
%   P1      box:dmp       primers/p1-dmp.tex        Section sec:dmp           (%PRIMER:dmp)
%   P2      box:qi        primers/p2-qi.tex         Section sec:fix           (%PRIMER:qi, before "The frame is classical")
%   P3      box:gr        primers/p3-gr.tex         Section sec:fields        (%PRIMER:gr, before "The frame")
%   P4      box:cq        primers/p4-cq.tex         Section sec:cq            (%PRIMER:cq)
%   P5      box:holo      primers/p5-holo.tex       Section sec:holo          (%PRIMER:holo)
%   P6      box:cosmo     primers/p6-cosmo.tex      Section sec:cosmo         (%PRIMER:cosmo)
%   P7      box:galaxy    primers/p7-galaxy.tex     Section sec:sparc         (or sec:darkmatter)
%   P8      box:stats     primers/p8-stats.tex      Section sec:scoring
%   P9      box:particle  primers/p9-particle.tex   Section sec:matter        (%PRIMER:particle)
%   P10     box:hadron    primers/p10-hadron.tex    Section sec:trees         (%PRIMER:hadron)
%
% Each box is background only: established physics or DMP's own text, with citations; none states a result of this paper.
% Every key cited is defined in bib-core.tex, bib-extra.tex or sections/*-bib.tex.
\newcommand{\primerDMP}{\input{primers/p1-dmp}}
\newcommand{\primerQI}{\input{primers/p2-qi}}
\newcommand{\primerGR}{\input{primers/p3-gr}}
\newcommand{\primerCQ}{\input{primers/p4-cq}}
\newcommand{\primerHolo}{\input{primers/p5-holo}}
\newcommand{\primerCosmo}{\input{primers/p6-cosmo}}
\newcommand{\primerGalaxy}{\input{primers/p7-galaxy}}
\newcommand{\primerStats}{\input{primers/p8-stats}}
\newcommand{\primerParticle}{\input{primers/p9-particle}}
\newcommand{\primerHadron}{\input{primers/p10-hadron}}
 once after
% \begin{document} defines the macros below, and \primerDMP and so on place the boxes (build_bib.py then numbers the primers'
% references from the position of that \input).
%
%   primer  label         file                      first use
%   P1      box:dmp       primers/p1-dmp.tex        Section sec:dmp           (%PRIMER:dmp)
%   P2      box:qi        primers/p2-qi.tex         Section sec:fix           (%PRIMER:qi, before "The frame is classical")
%   P3      box:gr        primers/p3-gr.tex         Section sec:fields        (%PRIMER:gr, before "The frame")
%   P4      box:cq        primers/p4-cq.tex         Section sec:cq            (%PRIMER:cq)
%   P5      box:holo      primers/p5-holo.tex       Section sec:holo          (%PRIMER:holo)
%   P6      box:cosmo     primers/p6-cosmo.tex      Section sec:cosmo         (%PRIMER:cosmo)
%   P7      box:galaxy    primers/p7-galaxy.tex     Section sec:sparc         (or sec:darkmatter)
%   P8      box:stats     primers/p8-stats.tex      Section sec:scoring
%   P9      box:particle  primers/p9-particle.tex   Section sec:matter        (%PRIMER:particle)
%   P10     box:hadron    primers/p10-hadron.tex    Section sec:trees         (%PRIMER:hadron)
%
% Each box is background only: established physics or DMP's own text, with citations; none states a result of this paper.
% Every key cited is defined in bib-core.tex, bib-extra.tex or sections/*-bib.tex.
\newcommand{\primerDMP}{% GUID v6, Primer P1 (label box:dmp): DMP and trust systems, with a worked act. Background only; it states no result of the paper.
% Place it at the start of Section sec:dmp (marker %PRIMER:dmp in part1.tex): \input{primers/p1-dmp}
\begin{primer}{DMP and trust systems}{box:dmp}
\emph{What a trust system is for.} In a computer system an access-control or trust protocol decides, for every request, who may do what
to which information, and it keeps the evidence of what was done. Its classic design rules are to grant the least privilege an act needs,
to demand two independent keys for what matters (separation of privilege), and to keep a complete audit trail~\cite{saltzer1975}. Its
formal tools include lattices that say where information may flow~\cite{denning}, noninterference, which says what one side may learn of
the other~\cite{gm}, and quorums of witnesses that fail independently~\cite{mr}. DMP~\cite{dmp} states such a system as three axioms and a
ground, G0; every later statement is a numbered theorem (T) or derived constraint (C).

\emph{Words.} A \emph{part} is ``a Duality of Trust'': a trust half and an information half. An \emph{act} is what a part does. A
\emph{layer} is where trust over acts is decided; its two \emph{witnesses} are its halves. The \emph{trust half} decides whether the act
may happen, the \emph{information half} whether it is meant and right. Two witnesses are \emph{independent} when they share no member of the
\emph{failure classes} the layer declares (host, key, author, engine, person, time source, and whatever the minting lists). A
\emph{standing} is trust held in the field of Axiom~2. It has a source (the witness that granted it), a depth (its \emph{grant chain} back
to the first pair) and a freshness that decays with time; \emph{staleness} is that decay. A \emph{record} is evidence that outlives its act.

\emph{The gate.} Each witness returns a verdict in the lattice $\textsf{refuse}<\textsf{hold}<\textsf{pass}$, and the gate of layer $L$
is their meet, $\operatorname{gate}_L(a,t)=v_{w_T}(a,t)\wedge v_{w_I}(a,t)$. The stricter answer stands, and the act is performed only on
$\textsf{pass}\wedge\textsf{pass}$.

\emph{The axioms in one line each.}
\begin{itemize}[nosep,leftmargin=*]
\item \emph{Axiom~1, the Duality.} A part acts only while both halves of its layer trust it for that act; neither overrules the other. Its
\emph{balance}: every part works towards the trust its boundary declares.
\item \emph{Axiom~2, Trust.} A witness trusts only on evidence it verified itself. ``Trust is gravitational. It is a field, never a
position.'' A strong enough boundary ``bends trustspace'' so that only entropy leaves.
\item \emph{Axiom~3, Information.} Every grant, use and loss of trust leaves a record that every witness of the act reads in full. Any other
reader sees the \emph{counts} always and the content only by a \emph{grant}; from trust, nothing about information can be derived but its
existence. Every layer is a tree rooted at its own Duality. The \emph{layout}, where an item sits in that tree, is ``the median of that
demand'', the place of fewest hops in total, recomputed by both halves and never by one.
\end{itemize}

\emph{A worked act} (Figure~\ref{fig:act}). A part $p$ asks to change one entry of a directory. The trust half checks $p$'s standing: it
was granted by a witness, lies two grants from the minted pair and is still fresh, so the verdict is \textsf{pass}. The information half
finds the change well formed but not yet confirmed: \textsf{hold}. The gate gives $\textsf{pass}\wedge\textsf{hold}=\textsf{hold}$, and
nothing is done. Once the change is confirmed, both return \textsf{pass} and the act is performed. A record is written at the act, and both
witnesses can read it in full. A reader outside the layer, say an auditor, sees the counts (here, one more act of this layer) but not the
new entry, unless it is granted to him. From the trust side alone one learns only that the entry exists.

\emph{Containment and G0.} A boundary \emph{contains} information when no part beyond it is trusted for any act on that information.
G0: ``The dual gate cannot witness its own first instance: somebody has to trust the first two witnesses.'' The first act, the
\emph{minting}, makes at least five parts the holders of the \emph{root keys}; it is stated, not derived, and every later standing
descends from it through grant chains.

\emph{The theorems this paper cites.}
\begin{itemize}[nosep,leftmargin=*]
\item \emph{T4}: every parameter is configuration. DMP fixes structure, never a value, and so derives no constant.
\item \emph{T27}: a contained item is destroyed at its scheduled date; with T31 it describes the trap.
\item \emph{T31}: contained information reaches no part by any act. What still leaves is minimal, thermal emission, the \emph{entropy
floor}, which ``carries entropy and no information''. Outside, one learns only the size, the staleness and the rotation interval.
\item \emph{T32}: the VIBE field. Each member has a depth, a mass (one plus the members deriving through it) and a shortfall; its
potential $\Phi_{\mathrm V}\ge0$ never rises under the system's own acts and vanishes in the \emph{balanced state}. Readers other than the
witnesses see one number, the boundary's temperature, its mass-weighted staleness.
\item \emph{T37}: information that no admission names is contained at the item itself.
\item \emph{T40}: a minting originates only from G0 together with a state in which every axiom and theorem provably holds.
\end{itemize}
\end{primer}
}
\newcommand{\primerQI}{% GUID v6, Primer P2 (label box:qi): quantum information. Background only; it states no result of the paper.
% Place it before "The frame is classical" (marker %PRIMER:qi in part1.tex): \input{primers/p2-qi}
\begin{primer}{Quantum information}{box:qi}
\emph{States and processes.} A quantum state is a \emph{density matrix} $\rho$: Hermitian, positive ($\rho\succeq0$) and of unit trace.
It is pure when $\rho=\ket\psi\bra\psi$ and mixed otherwise. The state of one part $A$ of a system $AB$ is the partial trace
$\rho_A=\Tr_B\rho_{AB}$~\cite{nielsen2010}. Every physical process is a \emph{completely positive, trace-preserving} (CPTP) map,
$\rho\mapsto\sum_kK_k\rho K_k^\dagger$ with $\sum_kK_k^\dagger K_k=\mathbb 1$; ``completely'' means that the map stays positive when it
acts on one part of a larger system. A measurement with outcomes $k$ returns $k$ with probability $p_k=\Tr(K_k\rho K_k^\dagger)$ and leaves
the state $K_k\rho K_k^\dagger/p_k$.

\emph{Dephasing is reading.} Dephasing in a basis $\{\ket k\}$, $\rho\mapsto\sum_k\ket k\bra k\rho\ket k\bra k$, erases the off-diagonal
elements (the coherences) and keeps the diagonal. It is what a measurement in that basis does to a system when the outcome is written
somewhere else. A system that is dephased in one fixed basis after every interaction behaves as a classical variable.

\emph{Entanglement.} A state of $AB$ is \emph{separable} if it is a mixture of products, $\rho_{AB}=\sum_ip_i\,\rho_A^i\otimes\rho_B^i$,
and \emph{entangled} otherwise~\cite{horodecki2009}. Two numbers measure it for two qubits.
\begin{itemize}[nosep,leftmargin=*]
\item The \emph{negativity} $\mathcal N$ is the magnitude of the sum of the negative eigenvalues of the partial transpose $\rho^{T_B}$
(the transpose taken on $B$ alone)~\cite{vw}. For two qubits $\mathcal N>0$ exactly when the state is entangled~\cite{peres1996,horodecki1996}.
\item The \emph{concurrence} of a pure two-qubit state $a\ket{00}+b\ket{01}+c\ket{10}+d\ket{11}$ is $C=2|ad-bc|$, from 0 for a product to 1
for a maximally entangled state~\cite{wootters1998}. For pure two-qubit states $\mathcal N=C/2$.
\end{itemize}
An \emph{entanglement witness} is an observable $W$ with $\Tr(W\sigma)\ge0$ for every separable $\sigma$; a measured $\Tr(W\rho)<0$
certifies entanglement~\cite{guhne2009}. A given witness misses every entangled state for which $\Tr(W\rho)\ge0$.

\emph{LOCC.} In local operations and classical communication each party acts on its own system and the parties exchange only measurement
outcomes. LOCC maps take separable states to separable states, and no entanglement measure grows under them~\cite{vw,horodecki2009}.
Creating entanglement between two systems therefore needs a carrier that is itself quantum.

\emph{No-cloning and no-broadcasting.} No process copies two pure states that are neither equal nor orthogonal~\cite{wz,dieks}. A family
of states can be \emph{broadcast}, so that each of two readers ends up with the original, only if its members commute pairwise~\cite{barnum}.
A quantity that several readers hold in full is therefore diagonal in one basis, which is the core of quantum Darwinism~\cite{zurek}.

\emph{Worked example: a dephased mediator.} Qubits $A$ and $B$ start in $\ket+=(\ket0+\ket1)/\sqrt2$ and a mediator qubit $M$ in $\ket0$.
They interact only through $M$: a CNOT from $A$ to $M$, a controlled-$Z$ between $M$ and $B$, and the same CNOT again.
\begin{itemize}[nosep,leftmargin=*]
\item \emph{$M$ coherent.} The first CNOT gives $(\ket{0_A0_M}+\ket{1_A1_M})/\sqrt2$. The controlled-$Z$ turns $B$ into
$\ket-=(\ket0-\ket1)/\sqrt2$ in the second branch, and the second CNOT returns $M$ to $\ket0$. What remains is
$(\ket{0}_A\ket{+}_B+\ket{1}_A\ket{-}_B)/\sqrt2$, a maximally entangled state: $C=1$, $\mathcal N=1/2$.
\item \emph{$M$ dephased} in $\{\ket0,\ket1\}$ after the first CNOT. The state becomes the mixture
$\tfrac12(\ket{00}\bra{00}+\ket{11}\bra{11})_{AM}\otimes\ket+\bra+_B$, and the same two steps leave
$\tfrac12(\ket0\bra0_A\otimes\ket+\bra+_B+\ket1\bra1_A\otimes\ket-\bra-_B)$. This is separable: $\mathcal N=0$.
\end{itemize}
The dephased mediator carried one classical bit from $A$ to $B$; that is LOCC, and it correlates the two qubits without entangling them.
\end{primer}
}
\newcommand{\primerGR}{% GUID v6, Primer P3 (label box:gr): general relativity and the tetrad. Background only; it states no result of the paper.
% Place it before "The frame" (marker %PRIMER:gr in part1.tex): \input{primers/p3-gr}
\begin{primer}{General relativity and the tetrad}{box:gr}
\setlength{\parskip}{3pt plus 1pt}%
\emph{The metric.} General relativity describes gravity as the geometry of spacetime. A metric $g_{\mu\nu}$ gives the interval
$\dd s^2=g_{\mu\nu}\dd x^\mu\dd x^\nu$, with signature $(-,+,+,+)$, and Einstein's equation
$G_{\mu\nu}+\Lambda g_{\mu\nu}=(8\pi G/c^4)\,T_{\mu\nu}$ ties its curvature to the energy-momentum tensor $T_{\mu\nu}$ of matter and
radiation~\cite{wald1984}.

\emph{The tetrad.} A tetrad, or frame, is a set of four orthonormal one-forms $e^a=e^a{}_\mu\dd x^\mu$ ($a=0,\dots,3$) with
$g_{\mu\nu}=\eta_{ab}e^a{}_\mu e^b{}_\nu$ and $\eta=\operatorname{diag}(-1,1,1,1)$. At each point it can be rotated and boosted,
${e^a\to\Lambda^a{}_b(x)\,e^b}$, without changing $g_{\mu\nu}$: this is the \emph{local Lorentz freedom}. Of the tetrad's 16 components, 6 are
that freedom and 10 are the metric. General relativity can be written with the tetrad as its only field~\cite{aldrovandi}. A single unit
timelike vector field added to the metric is different: it picks out a rest frame at every point. That is Einstein--aether
theory~\cite{aether}, and its preferred-frame effects are tightly bounded, for example by the equal speeds of gravitational waves and
light~\cite{gw170817}.

\emph{Free fall.} Matter couples \emph{minimally} when it feels gravity only through $g_{\mu\nu}$, with partial derivatives replaced by
covariant ones. Then $\nabla^\mu T_{\mu\nu}=0$, and a small body follows a geodesic,
$\ddot x^\mu+\Gamma^\mu{}_{\alpha\beta}\dot x^\alpha\dot x^\beta=0$ (dots: derivatives by proper time), whatever it is made of. That is the
weak equivalence principle. Experiments compare the accelerations $a_1$, $a_2$ of two materials through the E\"otv\"os ratio
${\eta=2|a_1-a_2|/|a_1+a_2|}$~\cite{microscope}; with indices, $\eta$ is the Minkowski metric.

\emph{PPN parameters.} In weak fields, metric theories are compared through the parametrised post-Newtonian (PPN) expansion~\cite{will}.
With $U>0$ the Newtonian potential,
\[
g_{00}=-1+\frac{2U}{c^2}-2\beta\frac{U^2}{c^4},\qquad g_{ij}=\Big(1+2\gamma\frac{U}{c^2}\Big)\delta_{ij}.
\]
$\gamma$ says how much space curvature a unit mass produces; light bending and the Shapiro delay scale as $(1+\gamma)/2$. $\beta$ measures
the nonlinearity of gravity and enters the perihelion advance with $\gamma$. General relativity has $\gamma=\beta=1$.

\emph{The lapse.} Split into space and time, spacetime has a \emph{lapse} $N$: the proper time that a static clock records per unit of
coordinate time. In a weak potential $\Phi=-U$ it is $N=\sqrt{1+2\Phi/c^2}$, outside a mass $M$ it is $N=\sqrt{1-2GM/rc^2}$, and it
vanishes at a horizon.

\emph{Black holes.} A stationary black hole is fixed by three numbers: its mass $M$, angular momentum $J$ and charge
$Q$~\cite{israel,kerr}. The rotating hole (Kerr) has its horizon at $r_+=(GM/c^2)\big(1+\sqrt{1-a_*^2}\big)$, with spin
$a_*=cJ/GM^2\le1$~\cite{kerr}. Its \emph{surface gravity} $\kappa$ is the acceleration, as measured from far away, needed to hold a body
at the horizon; for a hole without spin $\kappa=c^4/4GM$. Quantum fields make the horizon radiate at $k_BT_H=\hbar\kappa/2\pi c$~\cite{hawking75}.
The \emph{area theorem}: with positive energy, the horizon area of classical general relativity never decreases~\cite{hawking1971}. It is
the second law of black-hole mechanics~\cite{bch}, with entropy $S=k_Bc^3A/4G\hbar$~\cite{bekenstein73}.
\end{primer}
}
\newcommand{\primerCQ}{% GUID v6, Primer P4 (label box:cq): classical-quantum dynamics and collapse. Background only; it states no result of the paper.
% Place it before "Consistent classical-quantum dynamics" (marker %PRIMER:cq in part1.tex): \input{primers/p4-cq}
\begin{primer}{Classical-quantum dynamics and collapse}{box:cq}
\setlength{\parskip}{3pt plus 1pt}%
\emph{The classical-quantum (CQ) master equation.} A consistent coupling keeps the joint description linear, positive and
normalised, as probability theory demands. The joint state is a family $\varrho(z)$ of positive operators, one for each value
$z$ of the classical variables (here the frame), normalised as $\int\dd z\,\Tr\varrho(z)=1$~\cite{oppenheim}. For one classical variable
$z$ driven by one quantum observable $L$ (for gravity, the mass density), continuous Markovian dynamics has the form~\cite{osswd}
\[
\partial_t\varrho=-\frac i\hbar[H,\varrho]+D_0\Big(L\varrho L-\tfrac12\{L^2,\varrho\}\Big)-\frac{D_1}2\,\partial_z\{L,\varrho\}
+\frac{D_2}2\,\partial_z^2\varrho ,
\]
with three coefficients beyond quantum mechanics:
\begin{itemize}[nosep,leftmargin=*]
\item $D_0$, \emph{decoherence}: the quantum system loses coherence between states that source the classical field differently (for
gravity, different mass distributions);
\item $D_1$, \emph{back-reaction}: a drift of $z$ driven by the quantum state, that is, how matter pushes the frame;
\item $D_2$, \emph{diffusion}: random kicks of $z$, that is, noise in the frame.
\end{itemize}
Complete positivity requires $D_1^2\le4D_0D_2$ here, and $D_0>0$ wherever $D_1\ne0$; Oppenheim et al.\ give the general matrix form of
this trade-off~\cite[Eq.~(23)]{osswd}, which is Eq.~\eqref{eq:tradeoff} of this paper. In words: \emph{a push implies noise}. If quantum matter pushes a
classical field at all, the matter decoheres and the field diffuses, and the product of the two is bounded below by the square of the push,
$D_0D_2\ge D_1^2/4$.

\emph{Measurement and feedback.} An older family of consistent models makes gravity a channel. Each mass's position, or its mass density,
is weakly and continuously measured, and a classical force computed from the record acts back on the other masses (Kafri, Taylor and
Milburn, KTM~\cite{ktm}; Tilloy and Di\'osi~\cite{tilloydiosi}). Newton's law then holds on average, the measurement decoheres
superpositions, and the feedback adds noise. The masses talk only through a classical record, so the channel is LOCC
(Primer~\ref{box:qi}) and cannot entangle them. A weaker measurement decoheres less but gives a noisier record and so more heating; KTM find
both effects at least of the order of the gravitational coupling. The least decoherence of two masses over all measurement strengths is
what this paper calls the \emph{KTM minimum}, $\omega/2$, where $\omega$ is the rate at which a coherent field would build up the
entangling phase.

\emph{Gravity-related collapse.} Di\'osi and Penrose proposed that a superposition of two mass distributions $\rho_1$, $\rho_2$ decays at
the rate $\lambda=E_\Delta/\hbar$, where
\[
E_\Delta=\frac G2\int\!\!\int\dd^3x\,\dd^3y\,\frac{[\rho_1(\mathbf x)-\rho_2(\mathbf x)][\rho_1(\mathbf y)-\rho_2(\mathbf y)]}
{|\mathbf x-\mathbf y|}
\]
is the gravitational self-energy of their difference~\cite{diosi87,penrose96,bassi2013}. For a body of mass $M$ and size $R$ displaced by
more than its size, $\lambda\approx GM^2/\hbar R$: the \emph{Di\'osi--Penrose (DP) rate}. For point-like nuclei $E_\Delta$ diverges, so the
mass density must be smeared over a length $R_0$. The same noise makes charged matter radiate. An underground germanium detector sees no
such radiation, which excludes smearing at nuclear size and requires $R_0$ of order $10^{-10}$\,m or more~\cite{donadi,figurato}. In the
Newtonian limit of CQ gravity the DP kernel is one admissible choice of the decoherence term~\cite{osswd}.
\end{primer}
}
\newcommand{\primerHolo}{% GUID v6, Primer P5 (label box:holo): holography and quantum error correction. Background only; it states no result of the paper.
% Place it at the start of Section sec:holo (marker %PRIMER:holo in part1.tex): \input{primers/p5-holo}
\begin{primer}{Holography and quantum error correction}{box:holo}
\emph{Stabiliser codes.} A quantum error-correcting code stores $k$ logical qubits in $n$ physical ones so that some errors can be
undone. In a \emph{stabiliser code} the code space is the joint $+1$ eigenspace of $n-k$ commuting products of Pauli matrices, the
\emph{stabilisers}~\cite{gottesman1997}. Measuring them gives the \emph{syndrome}. An error that anticommutes with a stabiliser flips its
sign and so shows in the syndrome. The stabilisers commute with the logical operators, so the syndrome reveals nothing about the stored
content. A code of distance $d$, written $[[n,k,d]]$, corrects any $\lfloor(d-1)/2\rfloor$ unknown errors, or any $d-1$
\emph{erasures} (losses at known places).

\emph{The five-qubit code.} The smallest code that corrects every single-qubit error is $[[5,1,3]]$, with the stabiliser $XZZXI$ and its
cyclic shifts~\cite{laflamme,bennett1996}. It corrects any two erasures, so its logical qubit can be recovered from any three of the five
qubits and from no two: a $((3,5))$ quantum threshold scheme~\cite{cgl}.

\emph{Perfect tensors and the HaPPY code.} A tensor with $2m$ legs is \emph{perfect} if every split of its legs into $m$ and $m$ is, up to
normalisation, a unitary map. The encoding map of the five-qubit code, with its logical leg as a sixth, is such a tensor. Pastawski,
Yoshida, Harlow and Preskill tile the hyperbolic disc with pentagons, place one perfect tensor on each, contract neighbouring legs, and
leave one leg per tile in the bulk and the outermost legs on the boundary~\cite{happy}. The network is an isometry from the bulk to the
boundary, a toy model of the bulk-boundary correspondence of AdS/CFT~\cite{maldacena}.

\emph{Ryu--Takayanagi as a minimal cut.} In AdS/CFT the entropy of a boundary region $A$ is $S(A)=\mathrm{Area}(\gamma_A)/4G\hbar$ (with
$c=k_B=1$), where $\gamma_A$ is the bulk surface of least area anchored on the edge of $A$~\cite{rt}. In a network of perfect tensors the
entropy of $A$, in bits, is for many regions the number of qubit legs cut by the smallest cut that separates $A$ from the rest of the
boundary~\cite{happy}.

\emph{Reconstruction.} The bulk region between $A$ and $\gamma_A$ is the \emph{entanglement wedge} of $A$. Every bulk operator inside it can
be represented by an operator acting on $A$ alone~\cite{adh,dong2016}: the bulk is encoded redundantly, like a logical qubit. A code has
\emph{complementary recovery} when what cannot be reconstructed on $A$ can be reconstructed on the rest of the boundary.

\emph{The area operator.} Harlow showed that every code with complementary recovery obeys a quantum Ryu--Takayanagi formula
(Eq.~\eqref{eq:harlow} in Section~\ref{sec:holo}). Its area term is the expectation value of an \emph{area operator} that lies in the centre
of the algebra reconstructable on $A$: it commutes with everything that either side can do~\cite{harlow}. A \emph{fixed-area state} is an
eigenstate of this operator, and such states have a flat entanglement spectrum~\cite{dhm}. A general state of quantum gravity is a
superposition of areas.
\end{primer}
}
\newcommand{\primerCosmo}{% GUID v6, Primer P6 (label box:cosmo): cosmology. Background only; it states no result of the paper.
% Place it at the start of Section sec:cosmo (marker %PRIMER:cosmo in part1.tex): \input{primers/p6-cosmo}
\begin{primer}{Cosmology}{box:cosmo}
\emph{The expanding universe.} On large scales the universe is homogeneous and isotropic. Distances grow with a scale factor $a(t)$, light
from the past is redshifted by $1+z=a_0/a$, and the expansion rate $H=\dot a/a$ obeys, for a flat universe,
$H^2=8\pi G\rho/3+\Lambda c^2/3$. Each component's share is $\Omega_i=\rho_i/\rho_c$ with the critical density $\rho_c=3H^2/8\pi G$; the
value today is $H_0$.

\emph{$\Lambda$CDM.} The standard model of cosmology adds to the known particles a cold, collisionless \emph{dark matter} and a
cosmological constant $\Lambda$. Six parameters fix it: the densities of baryons and dark matter ($\Omega_bh^2$, $\Omega_ch^2$), the angular
size of the sound horizon, the optical depth to reionisation, and the amplitude $A_s$ and tilt $n_s$ of the primordial
fluctuations~\cite{planck18}.

\emph{Inflation.} An early epoch of nearly exponential expansion, $a\propto e^{Ht}$, flattens and smooths the universe and stretches
quantum fluctuations into the seeds of galaxies~\cite{guth}. It is driven by a scalar field, the \emph{inflaton}, with potential $V$, and
its length is counted in \emph{e-folds}, $N=\ln(a_{\rm end}/a_{\rm start})$. In \emph{slow roll} the field creeps down a flat potential;
with the reduced Planck mass $M_{\rm pl}$, the parameters $\epsilon=\tfrac12M_{\rm pl}^2(V'/V)^2$ and $\eta_V=M_{\rm pl}^2V''/V$ are small. The
fluctuations then have a nearly scale-invariant spectrum with tilt $n_s\approx1-6\epsilon+2\eta_V$, and gravitational waves with a
tensor-to-scalar ratio $r\approx16\epsilon$~\cite{planck2020x}.

\emph{Reheating and oscillons.} After inflation the inflaton oscillates about the minimum of its potential and decays into ordinary
particles, which \emph{reheats} the universe~\cite{kls}. In a broad class of models it first fragments into \emph{oscillons}, long-lived
localised lumps of the field~\cite{amin}.

\emph{The probes.} The cosmic microwave background (CMB) was released at recombination, $z\approx1100$. Its acoustic peaks are the imprint of
sound waves in the early plasma: their spacing fixes the sound horizon and their heights the densities~\cite{planck18}. The same sound
horizon, about 150\,Mpc today, is imprinted on the clustering of galaxies as \emph{baryon acoustic oscillations} (BAO), a standard ruler
for distance against redshift~\cite{desi2}.

\emph{Dark energy.} A cosmological constant has the equation of state $w=p/\rho c^2=-1$ at all times. DESI fits the
Chevallier--Polarski--Linder (CPL) form $w(a)=w_0+w_a(1-a)$~\cite{chevallier2001,linder2003}. A \emph{thawing} dark energy has $w_0>-1$
and $w_a<0$: it was close to $-1$ early and weakens today.

\emph{Data names and tensions.} P-ACT combines the CMB spectra of Planck and of the Atacama Cosmology Telescope (ACT, data release 6);
P-ACT-LB adds CMB lensing (L) and DESI BAO (B)~\cite{louis2025}. The \emph{$H_0$ tension} is the disagreement between the Hubble constant
inferred from the CMB with $\Lambda$CDM, near 67\,km\,s$^{-1}$\,Mpc$^{-1}$, and the value from a distance ladder of Cepheids and supernovae,
near 73~\cite{planck18,riess2022}. A ladder calibrated on the tip of the red-giant branch lies in between~\cite{freedman2025}.
\end{primer}
}
\newcommand{\primerGalaxy}{% GUID v6, Primer P7 (label box:galaxy): galaxy dynamics. Background only; it states no result of the paper.
% Place it at the start of Section sec:sparc (or at Section sec:darkmatter if the lead prefers): \input{primers/p7-galaxy}
\begin{primer}{Galaxy dynamics}{box:galaxy}
\setlength{\parskip}{3pt plus 1pt}%
\emph{Rotation curves.} Gas and stars in a disc galaxy move on nearly circular orbits. Their speed $V$ against radius $R$ is the
\emph{rotation curve}, and the centripetal acceleration is $g_{\rm obs}=V^2/R$. Far out, most curves stay flat where the visible mass alone
would make them fall as $R^{-1/2}$: the missing-mass problem.

\emph{Baryonic mass models.} The visible, or \emph{baryonic}, mass has three parts, each with its own contribution to $V^2$, computed from
its observed surface density: the gas (atomic hydrogen with its helium), the stellar disc and the central bulge. Starlight becomes stellar
mass through a mass-to-light ratio $\Upsilon$; at a wavelength of 3.6\,\textmu m it is near 0.5 for discs and 0.7 for
bulges~\cite{mcgaugh2016}. The sum of the three, divided by $R$, is the baryonic acceleration $g_{\rm bar}$.

\emph{Two empirical laws.} The \emph{baryonic Tully--Fisher relation} (BTFR): a galaxy's baryonic mass grows as the fourth power of its flat
rotation speed, $M_{\rm bar}\propto V_{\rm flat}^4$, with small scatter~\cite{lelli2016a}. The \emph{radial acceleration relation} (RAR): at
every radius of every galaxy, $g_{\rm obs}$ is one function of $g_{\rm bar}$. It is close to $g_{\rm bar}$ at high acceleration and close to
$\sqrt{g_{\rm bar}\,g_\dagger}$ below a scale $g_\dagger\approx10^{-10}$\,m\,s$^{-2}$~\cite{mcgaugh2016,lelli2017}.

\emph{MOND.} In Milgrom's modified Newtonian dynamics, $g_{\rm obs}=\nu(g_{\rm bar}/a_0)\,g_{\rm bar}$, where the new constant
$a_0$ is about $1.2\times10^{-10}$\,m\,s$^{-2}$~\cite{milgrom}. The \emph{interpolating function} $\nu(y)$ tends to 1 at high acceleration
($y\gg1$) and to $y^{-1/2}$ at low; the deep regime gives $V^4=GMa_0$, the BTFR. Common choices are the ``simple'' function
$\nu(y)=\tfrac12+(\tfrac14+1/y)^{1/2}$ and the ``standard'' one, and the RAR's fitting function is another~\cite{famaey2012}.

\emph{Dark-matter halos.} In cold dark matter each galaxy sits in a halo. Simulated halos follow the NFW profile
\[
\rho(r)=\frac{\rho_s}{(r/r_s)(1+r/r_s)^2}
\]
of Navarro, Frenk and White~\cite{navarro1996}. Its mass $M_{200}$ is counted inside the radius $r_{200}$ that encloses 200 times the
critical density, and its concentration $c=r_{200}/r_s$ falls slowly with mass~\cite{dutton2014}. At the centre the density rises as
$r^{-1}$, a \emph{cusp}. Many dwarf and low-surface-brightness galaxies prefer a \emph{core} of nearly constant density, as in the Burkert
profile $\rho=\rho_0r_0^3/[(r+r_0)(r^2+r_0^2)]$~\cite{burkert1995}; this mismatch is the cusp-core problem~\cite{deblok2010}. The
\emph{stellar-to-halo relation}, found by matching the abundances of galaxies and halos, gives the stellar mass a halo of given mass hosts;
the ratio peaks at a few per cent near $10^{12}$ solar masses~\cite{moster2013}.

\emph{SPARC.} The Spitzer Photometry and Accurate Rotation Curves sample holds 175 disc galaxies with 3.6\,\textmu m photometry and
rotation curves from atomic-hydrogen and H$\alpha$ data, with mass models for gas, disc and bulge~\cite{lelli2016}.
\end{primer}
}
\newcommand{\primerStats}{% GUID v6, Primer P8 (label box:stats): the statistics behind the scoring rules of Section sec:scoring. Background only; it states
% no result of the paper and does not repeat the rules themselves. Place it in Section sec:scoring: \input{primers/p8-stats}
\begin{primer}{Statistics}{box:stats}
\emph{Pulls and 95\%.} A measurement $x\pm\sigma$ usually means a Gaussian error with standard deviation $\sigma$. If a prediction
$x_G\pm\sigma_G$ is right and the two errors are independent, the difference $x-x_G$ is Gaussian with variance $\sigma^2+\sigma_G^2$, so the
\emph{pull} $p=(x-x_G)/\sqrt{\sigma^2+\sigma_G^2}$ is a standard normal number. It falls outside $\pm1.96$ with probability 5\% (two-sided),
so $|p|<1.96$ is a test at 95\%, and a published 95\% half-width $b$ corresponds to $\sigma=b/1.96$. \emph{Worked pull:} a model predicts
$10.0\pm0.3$ and an experiment finds $11.0\pm0.4$. Then $p=1.0/\sqrt{0.4^2+0.3^2}=1.0/0.5=2.0$: a FAIL at 95\%, although a right model
gives $|p|\ge2.0$ in 4.6\% of independent tests. For $n$ independent tests, about $0.05\,n$ FAILs are expected by chance, and the sum of the
squared pulls follows a $\chi^2$ distribution with $n$ degrees of freedom.

\emph{BIC and $\Delta$BIC.} A model with $k$ fitted parameters that reaches the maximum likelihood $\mathcal L_{\max}$ on $N$ data points has
$\mathrm{BIC}=-2\ln\mathcal L_{\max}+k\ln N$~\cite{schwarz1978}. The penalty $k\ln N$ charges each extra parameter. For large $N$ and
independent points, the difference $\Delta$BIC of two models approximates $-2$ times the logarithm of their Bayes factor, the ratio of the
probabilities of the data under the two models. A $\Delta$BIC of 2 to 6 is then positive evidence, 6 to 10 strong and above 10 very
strong~\cite{kassraftery1995}. The approximation assumes $N$ independent points. If the points come in groups that share an error, such
as the distance of one galaxy, the effective number of independent data is closer to the number of groups. A point-wise $\Delta$BIC then
overstates the evidence, roughly in proportion to the number of points per group.

\emph{The bootstrap.} To learn how much a fitted value would vary in a new sample, one draws many new samples of the same size from the
data themselves, with replacement, and refits each~\cite{efron1979}. The spread of the refitted values estimates the uncertainty; the 16th
to 84th percentiles bracket a range comparable to $\pm1\sigma$. The units resampled must be the independent ones, for rotation curves whole
galaxies, not single points.

\emph{Many tests.} With $m$ tests at 5\% each, the chance of at least one false positive, the \emph{family-wise error rate}, grows towards
$1-0.95^m$. Bonferroni's remedy tests each at $0.05/m$. Holm's method orders the $p$-values, $p_{(1)}\le\dots\le p_{(m)}$, and rejects
$p_{(i)}$ while $p_{(i)}\le0.05/(m-i+1)$; it keeps the family-wise error rate at 5\% under any dependence and never rejects less than
Bonferroni~\cite{holm1979}. Fisher's method combines $n$ \emph{independent} $p$-values: under the null each $p_i$ is uniform, so
$-2\ln p_i$ follows $\chi^2$ with 2 degrees of freedom, and $X^2=-2\sum_i\ln p_i$ follows $\chi^2$ with $2n$~\cite{fisher1932}. It responds
strongly to a few very small $p_i$.

\emph{Look-elsewhere and preregistration.} A search over many masses, scales or candidate values finds a small \emph{local} $p$-value
somewhere by chance. The \emph{global} $p$-value, the chance of so small a value anywhere in the search, is larger, roughly the local one
times a trials factor~\cite{grossvitells2010}. A test is confirmatory only if its hypothesis, data, statistic and decision rule were
fixed before the outcome was seen; \emph{preregistration} writes them down in a dated record beforehand, and every choice made afterwards
is then visible as exploratory~\cite{nosek2018}.
\end{primer}
}
\newcommand{\primerParticle}{% GUID v6, Primer P9 (label box:particle): particle physics. Background only; it states no result of the paper.
% Place it before "The spin ladder" (marker %PRIMER:particle in part1.tex): \input{primers/p9-particle}
\begin{primer}{Particle physics}{box:particle}
\setlength{\parskip}{3pt plus 1pt}%
\emph{The Higgs mechanism.} The Standard Model is a gauge theory of $SU(3)\times SU(2)\times U(1)$: colour for the strong force, weak isospin
and hypercharge for the electroweak force~\cite{glashow,weinberg67,salam}. The Higgs field is a doublet whose potential has its minimum
away from zero, at a vacuum value $v=246$\,GeV~\cite{higgs64,be64}. In that vacuum the W and Z bosons acquire masses proportional to $v$,
while one combination of the electroweak fields, the photon, stays massless. Quarks and charged leptons get their masses from
\emph{Yukawa couplings} $y_f$ to the doublet, $m_f=y_fv/\sqrt2$. The Higgs boson therefore couples to each fermion with strength $m_f/v$
and to the W and Z with $2m_V^2/v$.

\emph{The $\kappa$ framework.} Measurements of Higgs production and decay are summarised by \emph{coupling modifiers} $\kappa_i$, which
scale each Standard Model coupling, so that a rate involving particle $i$ scales as $\kappa_i^2$; the Standard Model has every
$\kappa_i=1$~\cite{david2012}.

\emph{Schemes and running.} A coupling or mass in quantum field theory is defined only together with a \emph{renormalisation scheme} and
an energy scale $\mu$. In the $\overline{\text{MS}}$ scheme only the divergent part of each loop is subtracted, and parameters run with
$\mu$. The \emph{on-shell} scheme instead fixes masses at the poles of the propagators and $\alpha$ at zero momentum transfer. The same
quantity can differ between the two by more than its measurement error, so every comparison states its scheme~\cite{erler2017,navas2024}.
The dependence on $\mu$ obeys the renormalisation-group equations (RGEs). At one loop
$\dd\alpha_i^{-1}/\dd\ln\mu=-b_i/2\pi$, with $(b_1,b_2,b_3)=(41/10,-19/6,-7)$ in the Standard Model (hypercharge normalised as in grand
unification). The strong coupling therefore weakens at high energy, and the hypercharge coupling grows. Precision work runs at two or
three loops~\cite{buttazzo2013}.

\emph{Dirac or Majorana.} A \emph{Dirac} mass joins a left-handed neutrino to a right-handed partner and conserves lepton number $L$. A
\emph{Majorana} mass joins the neutrino to its own antiparticle and changes $L$ by two. The lowest-dimension operator of Standard Model
fields that gives neutrinos mass is Weinberg's dimension-five operator, $(c/\Lambda)\,(LH)(LH)$, with $L$ the lepton doublet; in the Higgs
vacuum it gives a Majorana mass of order $cv^2/\Lambda$, small when the new scale $\Lambda$ is large~\cite{weinberg1979}. Majorana masses
can therefore also come from the Higgs vacuum.

\emph{Neutrinoless double-beta decay} ($0\nu\beta\beta$) turns two neutrons of a nucleus into two protons and two electrons with no
neutrinos. It changes $L$ by two and can occur only if neutrinos have a Majorana mass. Through light-neutrino exchange its rate is
proportional to $m_{\beta\beta}^2$, with the effective mass $m_{\beta\beta}=|\sum_iU_{ei}^2m_i|$ set by the mixing matrix $U$ and the
masses $m_i$~\cite{navas2024}. Xenon and germanium experiments search for it~\cite{kamlandzen2025,legend2025}.

\emph{Grand unification.} A grand unified theory (GUT) embeds the three gauge groups in one simple group, the smallest being
$SU(5)$~\cite{georgiglashow}, with a single coupling above a scale $M_X$. With the Standard Model's particles alone, the three running
couplings do not meet at one point~\cite{amaldi}. A GUT has two generic consequences. Its heavy gauge bosons let protons decay, with a
lifetime of order $M_X^4/(\alpha_G^2m_p^5)$; Super-Kamiokande's bound on $p\to e^+\pi^0$ excludes the minimal $SU(5)$
model~\cite{superk}. And its breaking leaves magnetic monopoles of mass near $M_X/\alpha_G$, about one per causal volume, which would
overfill the universe unless inflation dilutes them~\cite{preskill,guth}.
\end{primer}
}
\newcommand{\primerHadron}{% GUID v6, Primer P10 (label box:hadron): hadron spectroscopy. Background only; it states no result of the paper.
% Place it before "Hadrons as trees of dualities" (marker %PRIMER:hadron in part1.tex): \input{primers/p10-hadron}
\begin{primer}{Hadron spectroscopy}{box:hadron}
\emph{Colour.} Quarks carry the colour charge of $SU(3)$: a quark is a triplet $\mathbf3$, an antiquark an antitriplet $\mathbf{\bar3}$.
Only colour \emph{singlets} exist as free particles, the hadrons~\cite{navas2024}. The products
$\mathbf3\otimes\mathbf{\bar3}=\mathbf1\oplus\mathbf8$ and $\mathbf3\otimes\mathbf3=\mathbf{\bar3}\oplus\mathbf6$ give the ordinary
hadrons: a meson $q\bar q$ in the singlet, and a baryon $qqq$, whose first two quarks couple to $\mathbf{\bar3}$ and then with the third to
a singlet. One-gluon exchange attracts in the singlet $q\bar q$ and, half as strongly, in the $\mathbf{\bar3}$ $qq$, and repels in the
$\mathbf8$ and the $\mathbf6$.

\emph{Diquarks.} A \emph{diquark} is a pair of quarks in the attractive $\mathbf{\bar3}$. The ``good'' diquark, with spin 0 and a flavour-
antisymmetric pair such as $[ud]$, binds most; the ``bad'' one has spin 1~\cite{jaffe2005}.

\emph{Compact or molecular.} A hadron with more than three quarks, an \emph{exotic}, can be \emph{compact}, for example a diquark and an
antidiquark bound by colour forces, $[qq][\bar q\bar q]$~\cite{maiani2005}. Or it can be a \emph{hadronic molecule}: two colour-singlet
hadrons held together by the residual strong force, like the deuteron~\cite{guo2018}. Molecules sit just below the threshold
$m_1+m_2$ of their two constituents.

\emph{The colour-spin interaction.} In the constituent-quark model a hadron's mass is the sum of effective quark masses plus a hyperfine
term proportional to $(\boldsymbol\lambda_i\cdot\boldsymbol\lambda_j)(\boldsymbol\sigma_i\cdot\boldsymbol\sigma_j)/m_im_j$ for each pair,
where the $\boldsymbol\lambda$ are colour and the $\boldsymbol\sigma$ spin matrices. This chromomagnetic interaction of De R\'ujula, Georgi
and Glashow explains, for example, why the $\Delta$ is heavier than the nucleon~\cite{derujula1975}. Karliner and Rosner (KR) refined it
with separate effective quark masses in mesons and in baryons, the baryon masses for every constituent of a compact multiquark, and binding
terms for heavy pairs with $B(QQ)=B(Q\bar Q)/2$; with it they predicted the doubly charmed baryon and doubly heavy
tetraquarks~\cite{karliner2014,karliner2017}.

\emph{Thresholds and cusps.} Where a new two-hadron channel opens, at $M=m_1+m_2$, the scattering amplitude has a branch point. A
\emph{cusp} can then appear in a mass spectrum without any state, and a state that couples to the channel is pulled towards the threshold.
Many exotic candidates lie within a few MeV of a threshold~\cite{guo2018}.

\emph{Names.} The Particle Data Group names hadrons by their quantum numbers, not by their structure~\cite{navas2024}; the
$X(3872)$, for example, is listed as $\chi_{c1}(3872)$. LHCb's convention for exotics~\cite{gershon2022} uses $T$ for four-quark and $P$
for five-quark states, with subscripts for heavy flavour, $\psi$ standing for hidden charm $c\bar c$. So $T_{cc}^+$ is $cc\bar u\bar d$,
seen just below the $D^{*+}D^0$ threshold~\cite{lhcb2022h}, and the hidden-charm pentaquark $P_c(4312)^+$~\cite{lhcb2019h} becomes
$P_\psi^N(4312)^+$.

\emph{String junctions.} In a baryon the three colour flux tubes meet at a point, a \emph{string junction}~\cite{rossiveneziano1977}. A
compact tetraquark has two junctions and a compact pentaquark three, and some models add a constant energy per
junction~\cite{karliner2017a}.
\end{primer}
}
 once after
% \begin{document} defines the macros below, and \primerDMP and so on place the boxes (build_bib.py then numbers the primers'
% references from the position of that \input).
%
%   primer  label         file                      first use
%   P1      box:dmp       primers/p1-dmp.tex        Section sec:dmp           (%PRIMER:dmp)
%   P2      box:qi        primers/p2-qi.tex         Section sec:fix           (%PRIMER:qi, before "The frame is classical")
%   P3      box:gr        primers/p3-gr.tex         Section sec:fields        (%PRIMER:gr, before "The frame")
%   P4      box:cq        primers/p4-cq.tex         Section sec:cq            (%PRIMER:cq)
%   P5      box:holo      primers/p5-holo.tex       Section sec:holo          (%PRIMER:holo)
%   P6      box:cosmo     primers/p6-cosmo.tex      Section sec:cosmo         (%PRIMER:cosmo)
%   P7      box:galaxy    primers/p7-galaxy.tex     Section sec:sparc         (or sec:darkmatter)
%   P8      box:stats     primers/p8-stats.tex      Section sec:scoring
%   P9      box:particle  primers/p9-particle.tex   Section sec:matter        (%PRIMER:particle)
%   P10     box:hadron    primers/p10-hadron.tex    Section sec:trees         (%PRIMER:hadron)
%
% Each box is background only: established physics or DMP's own text, with citations; none states a result of this paper.
% Every key cited is defined in bib-core.tex, bib-extra.tex or sections/*-bib.tex.
\newcommand{\primerDMP}{% GUID v6, Primer P1 (label box:dmp): DMP and trust systems, with a worked act. Background only; it states no result of the paper.
% Place it at the start of Section sec:dmp (marker %PRIMER:dmp in part1.tex): % GUID v6, Primer P1 (label box:dmp): DMP and trust systems, with a worked act. Background only; it states no result of the paper.
% Place it at the start of Section sec:dmp (marker %PRIMER:dmp in part1.tex): \input{primers/p1-dmp}
\begin{primer}{DMP and trust systems}{box:dmp}
\emph{What a trust system is for.} In a computer system an access-control or trust protocol decides, for every request, who may do what
to which information, and it keeps the evidence of what was done. Its classic design rules are to grant the least privilege an act needs,
to demand two independent keys for what matters (separation of privilege), and to keep a complete audit trail~\cite{saltzer1975}. Its
formal tools include lattices that say where information may flow~\cite{denning}, noninterference, which says what one side may learn of
the other~\cite{gm}, and quorums of witnesses that fail independently~\cite{mr}. DMP~\cite{dmp} states such a system as three axioms and a
ground, G0; every later statement is a numbered theorem (T) or derived constraint (C).

\emph{Words.} A \emph{part} is ``a Duality of Trust'': a trust half and an information half. An \emph{act} is what a part does. A
\emph{layer} is where trust over acts is decided; its two \emph{witnesses} are its halves. The \emph{trust half} decides whether the act
may happen, the \emph{information half} whether it is meant and right. Two witnesses are \emph{independent} when they share no member of the
\emph{failure classes} the layer declares (host, key, author, engine, person, time source, and whatever the minting lists). A
\emph{standing} is trust held in the field of Axiom~2. It has a source (the witness that granted it), a depth (its \emph{grant chain} back
to the first pair) and a freshness that decays with time; \emph{staleness} is that decay. A \emph{record} is evidence that outlives its act.

\emph{The gate.} Each witness returns a verdict in the lattice $\textsf{refuse}<\textsf{hold}<\textsf{pass}$, and the gate of layer $L$
is their meet, $\operatorname{gate}_L(a,t)=v_{w_T}(a,t)\wedge v_{w_I}(a,t)$. The stricter answer stands, and the act is performed only on
$\textsf{pass}\wedge\textsf{pass}$.

\emph{The axioms in one line each.}
\begin{itemize}[nosep,leftmargin=*]
\item \emph{Axiom~1, the Duality.} A part acts only while both halves of its layer trust it for that act; neither overrules the other. Its
\emph{balance}: every part works towards the trust its boundary declares.
\item \emph{Axiom~2, Trust.} A witness trusts only on evidence it verified itself. ``Trust is gravitational. It is a field, never a
position.'' A strong enough boundary ``bends trustspace'' so that only entropy leaves.
\item \emph{Axiom~3, Information.} Every grant, use and loss of trust leaves a record that every witness of the act reads in full. Any other
reader sees the \emph{counts} always and the content only by a \emph{grant}; from trust, nothing about information can be derived but its
existence. Every layer is a tree rooted at its own Duality. The \emph{layout}, where an item sits in that tree, is ``the median of that
demand'', the place of fewest hops in total, recomputed by both halves and never by one.
\end{itemize}

\emph{A worked act} (Figure~\ref{fig:act}). A part $p$ asks to change one entry of a directory. The trust half checks $p$'s standing: it
was granted by a witness, lies two grants from the minted pair and is still fresh, so the verdict is \textsf{pass}. The information half
finds the change well formed but not yet confirmed: \textsf{hold}. The gate gives $\textsf{pass}\wedge\textsf{hold}=\textsf{hold}$, and
nothing is done. Once the change is confirmed, both return \textsf{pass} and the act is performed. A record is written at the act, and both
witnesses can read it in full. A reader outside the layer, say an auditor, sees the counts (here, one more act of this layer) but not the
new entry, unless it is granted to him. From the trust side alone one learns only that the entry exists.

\emph{Containment and G0.} A boundary \emph{contains} information when no part beyond it is trusted for any act on that information.
G0: ``The dual gate cannot witness its own first instance: somebody has to trust the first two witnesses.'' The first act, the
\emph{minting}, makes at least five parts the holders of the \emph{root keys}; it is stated, not derived, and every later standing
descends from it through grant chains.

\emph{The theorems this paper cites.}
\begin{itemize}[nosep,leftmargin=*]
\item \emph{T4}: every parameter is configuration. DMP fixes structure, never a value, and so derives no constant.
\item \emph{T27}: a contained item is destroyed at its scheduled date; with T31 it describes the trap.
\item \emph{T31}: contained information reaches no part by any act. What still leaves is minimal, thermal emission, the \emph{entropy
floor}, which ``carries entropy and no information''. Outside, one learns only the size, the staleness and the rotation interval.
\item \emph{T32}: the VIBE field. Each member has a depth, a mass (one plus the members deriving through it) and a shortfall; its
potential $\Phi_{\mathrm V}\ge0$ never rises under the system's own acts and vanishes in the \emph{balanced state}. Readers other than the
witnesses see one number, the boundary's temperature, its mass-weighted staleness.
\item \emph{T37}: information that no admission names is contained at the item itself.
\item \emph{T40}: a minting originates only from G0 together with a state in which every axiom and theorem provably holds.
\end{itemize}
\end{primer}

\begin{primer}{DMP and trust systems}{box:dmp}
\emph{What a trust system is for.} In a computer system an access-control or trust protocol decides, for every request, who may do what
to which information, and it keeps the evidence of what was done. Its classic design rules are to grant the least privilege an act needs,
to demand two independent keys for what matters (separation of privilege), and to keep a complete audit trail~\cite{saltzer1975}. Its
formal tools include lattices that say where information may flow~\cite{denning}, noninterference, which says what one side may learn of
the other~\cite{gm}, and quorums of witnesses that fail independently~\cite{mr}. DMP~\cite{dmp} states such a system as three axioms and a
ground, G0; every later statement is a numbered theorem (T) or derived constraint (C).

\emph{Words.} A \emph{part} is ``a Duality of Trust'': a trust half and an information half. An \emph{act} is what a part does. A
\emph{layer} is where trust over acts is decided; its two \emph{witnesses} are its halves. The \emph{trust half} decides whether the act
may happen, the \emph{information half} whether it is meant and right. Two witnesses are \emph{independent} when they share no member of the
\emph{failure classes} the layer declares (host, key, author, engine, person, time source, and whatever the minting lists). A
\emph{standing} is trust held in the field of Axiom~2. It has a source (the witness that granted it), a depth (its \emph{grant chain} back
to the first pair) and a freshness that decays with time; \emph{staleness} is that decay. A \emph{record} is evidence that outlives its act.

\emph{The gate.} Each witness returns a verdict in the lattice $\textsf{refuse}<\textsf{hold}<\textsf{pass}$, and the gate of layer $L$
is their meet, $\operatorname{gate}_L(a,t)=v_{w_T}(a,t)\wedge v_{w_I}(a,t)$. The stricter answer stands, and the act is performed only on
$\textsf{pass}\wedge\textsf{pass}$.

\emph{The axioms in one line each.}
\begin{itemize}[nosep,leftmargin=*]
\item \emph{Axiom~1, the Duality.} A part acts only while both halves of its layer trust it for that act; neither overrules the other. Its
\emph{balance}: every part works towards the trust its boundary declares.
\item \emph{Axiom~2, Trust.} A witness trusts only on evidence it verified itself. ``Trust is gravitational. It is a field, never a
position.'' A strong enough boundary ``bends trustspace'' so that only entropy leaves.
\item \emph{Axiom~3, Information.} Every grant, use and loss of trust leaves a record that every witness of the act reads in full. Any other
reader sees the \emph{counts} always and the content only by a \emph{grant}; from trust, nothing about information can be derived but its
existence. Every layer is a tree rooted at its own Duality. The \emph{layout}, where an item sits in that tree, is ``the median of that
demand'', the place of fewest hops in total, recomputed by both halves and never by one.
\end{itemize}

\emph{A worked act} (Figure~\ref{fig:act}). A part $p$ asks to change one entry of a directory. The trust half checks $p$'s standing: it
was granted by a witness, lies two grants from the minted pair and is still fresh, so the verdict is \textsf{pass}. The information half
finds the change well formed but not yet confirmed: \textsf{hold}. The gate gives $\textsf{pass}\wedge\textsf{hold}=\textsf{hold}$, and
nothing is done. Once the change is confirmed, both return \textsf{pass} and the act is performed. A record is written at the act, and both
witnesses can read it in full. A reader outside the layer, say an auditor, sees the counts (here, one more act of this layer) but not the
new entry, unless it is granted to him. From the trust side alone one learns only that the entry exists.

\emph{Containment and G0.} A boundary \emph{contains} information when no part beyond it is trusted for any act on that information.
G0: ``The dual gate cannot witness its own first instance: somebody has to trust the first two witnesses.'' The first act, the
\emph{minting}, makes at least five parts the holders of the \emph{root keys}; it is stated, not derived, and every later standing
descends from it through grant chains.

\emph{The theorems this paper cites.}
\begin{itemize}[nosep,leftmargin=*]
\item \emph{T4}: every parameter is configuration. DMP fixes structure, never a value, and so derives no constant.
\item \emph{T27}: a contained item is destroyed at its scheduled date; with T31 it describes the trap.
\item \emph{T31}: contained information reaches no part by any act. What still leaves is minimal, thermal emission, the \emph{entropy
floor}, which ``carries entropy and no information''. Outside, one learns only the size, the staleness and the rotation interval.
\item \emph{T32}: the VIBE field. Each member has a depth, a mass (one plus the members deriving through it) and a shortfall; its
potential $\Phi_{\mathrm V}\ge0$ never rises under the system's own acts and vanishes in the \emph{balanced state}. Readers other than the
witnesses see one number, the boundary's temperature, its mass-weighted staleness.
\item \emph{T37}: information that no admission names is contained at the item itself.
\item \emph{T40}: a minting originates only from G0 together with a state in which every axiom and theorem provably holds.
\end{itemize}
\end{primer}
}
\newcommand{\primerQI}{% GUID v6, Primer P2 (label box:qi): quantum information. Background only; it states no result of the paper.
% Place it before "The frame is classical" (marker %PRIMER:qi in part1.tex): % GUID v6, Primer P2 (label box:qi): quantum information. Background only; it states no result of the paper.
% Place it before "The frame is classical" (marker %PRIMER:qi in part1.tex): \input{primers/p2-qi}
\begin{primer}{Quantum information}{box:qi}
\emph{States and processes.} A quantum state is a \emph{density matrix} $\rho$: Hermitian, positive ($\rho\succeq0$) and of unit trace.
It is pure when $\rho=\ket\psi\bra\psi$ and mixed otherwise. The state of one part $A$ of a system $AB$ is the partial trace
$\rho_A=\Tr_B\rho_{AB}$~\cite{nielsen2010}. Every physical process is a \emph{completely positive, trace-preserving} (CPTP) map,
$\rho\mapsto\sum_kK_k\rho K_k^\dagger$ with $\sum_kK_k^\dagger K_k=\mathbb 1$; ``completely'' means that the map stays positive when it
acts on one part of a larger system. A measurement with outcomes $k$ returns $k$ with probability $p_k=\Tr(K_k\rho K_k^\dagger)$ and leaves
the state $K_k\rho K_k^\dagger/p_k$.

\emph{Dephasing is reading.} Dephasing in a basis $\{\ket k\}$, $\rho\mapsto\sum_k\ket k\bra k\rho\ket k\bra k$, erases the off-diagonal
elements (the coherences) and keeps the diagonal. It is what a measurement in that basis does to a system when the outcome is written
somewhere else. A system that is dephased in one fixed basis after every interaction behaves as a classical variable.

\emph{Entanglement.} A state of $AB$ is \emph{separable} if it is a mixture of products, $\rho_{AB}=\sum_ip_i\,\rho_A^i\otimes\rho_B^i$,
and \emph{entangled} otherwise~\cite{horodecki2009}. Two numbers measure it for two qubits.
\begin{itemize}[nosep,leftmargin=*]
\item The \emph{negativity} $\mathcal N$ is the magnitude of the sum of the negative eigenvalues of the partial transpose $\rho^{T_B}$
(the transpose taken on $B$ alone)~\cite{vw}. For two qubits $\mathcal N>0$ exactly when the state is entangled~\cite{peres1996,horodecki1996}.
\item The \emph{concurrence} of a pure two-qubit state $a\ket{00}+b\ket{01}+c\ket{10}+d\ket{11}$ is $C=2|ad-bc|$, from 0 for a product to 1
for a maximally entangled state~\cite{wootters1998}. For pure two-qubit states $\mathcal N=C/2$.
\end{itemize}
An \emph{entanglement witness} is an observable $W$ with $\Tr(W\sigma)\ge0$ for every separable $\sigma$; a measured $\Tr(W\rho)<0$
certifies entanglement~\cite{guhne2009}. A given witness misses every entangled state for which $\Tr(W\rho)\ge0$.

\emph{LOCC.} In local operations and classical communication each party acts on its own system and the parties exchange only measurement
outcomes. LOCC maps take separable states to separable states, and no entanglement measure grows under them~\cite{vw,horodecki2009}.
Creating entanglement between two systems therefore needs a carrier that is itself quantum.

\emph{No-cloning and no-broadcasting.} No process copies two pure states that are neither equal nor orthogonal~\cite{wz,dieks}. A family
of states can be \emph{broadcast}, so that each of two readers ends up with the original, only if its members commute pairwise~\cite{barnum}.
A quantity that several readers hold in full is therefore diagonal in one basis, which is the core of quantum Darwinism~\cite{zurek}.

\emph{Worked example: a dephased mediator.} Qubits $A$ and $B$ start in $\ket+=(\ket0+\ket1)/\sqrt2$ and a mediator qubit $M$ in $\ket0$.
They interact only through $M$: a CNOT from $A$ to $M$, a controlled-$Z$ between $M$ and $B$, and the same CNOT again.
\begin{itemize}[nosep,leftmargin=*]
\item \emph{$M$ coherent.} The first CNOT gives $(\ket{0_A0_M}+\ket{1_A1_M})/\sqrt2$. The controlled-$Z$ turns $B$ into
$\ket-=(\ket0-\ket1)/\sqrt2$ in the second branch, and the second CNOT returns $M$ to $\ket0$. What remains is
$(\ket{0}_A\ket{+}_B+\ket{1}_A\ket{-}_B)/\sqrt2$, a maximally entangled state: $C=1$, $\mathcal N=1/2$.
\item \emph{$M$ dephased} in $\{\ket0,\ket1\}$ after the first CNOT. The state becomes the mixture
$\tfrac12(\ket{00}\bra{00}+\ket{11}\bra{11})_{AM}\otimes\ket+\bra+_B$, and the same two steps leave
$\tfrac12(\ket0\bra0_A\otimes\ket+\bra+_B+\ket1\bra1_A\otimes\ket-\bra-_B)$. This is separable: $\mathcal N=0$.
\end{itemize}
The dephased mediator carried one classical bit from $A$ to $B$; that is LOCC, and it correlates the two qubits without entangling them.
\end{primer}

\begin{primer}{Quantum information}{box:qi}
\emph{States and processes.} A quantum state is a \emph{density matrix} $\rho$: Hermitian, positive ($\rho\succeq0$) and of unit trace.
It is pure when $\rho=\ket\psi\bra\psi$ and mixed otherwise. The state of one part $A$ of a system $AB$ is the partial trace
$\rho_A=\Tr_B\rho_{AB}$~\cite{nielsen2010}. Every physical process is a \emph{completely positive, trace-preserving} (CPTP) map,
$\rho\mapsto\sum_kK_k\rho K_k^\dagger$ with $\sum_kK_k^\dagger K_k=\mathbb 1$; ``completely'' means that the map stays positive when it
acts on one part of a larger system. A measurement with outcomes $k$ returns $k$ with probability $p_k=\Tr(K_k\rho K_k^\dagger)$ and leaves
the state $K_k\rho K_k^\dagger/p_k$.

\emph{Dephasing is reading.} Dephasing in a basis $\{\ket k\}$, $\rho\mapsto\sum_k\ket k\bra k\rho\ket k\bra k$, erases the off-diagonal
elements (the coherences) and keeps the diagonal. It is what a measurement in that basis does to a system when the outcome is written
somewhere else. A system that is dephased in one fixed basis after every interaction behaves as a classical variable.

\emph{Entanglement.} A state of $AB$ is \emph{separable} if it is a mixture of products, $\rho_{AB}=\sum_ip_i\,\rho_A^i\otimes\rho_B^i$,
and \emph{entangled} otherwise~\cite{horodecki2009}. Two numbers measure it for two qubits.
\begin{itemize}[nosep,leftmargin=*]
\item The \emph{negativity} $\mathcal N$ is the magnitude of the sum of the negative eigenvalues of the partial transpose $\rho^{T_B}$
(the transpose taken on $B$ alone)~\cite{vw}. For two qubits $\mathcal N>0$ exactly when the state is entangled~\cite{peres1996,horodecki1996}.
\item The \emph{concurrence} of a pure two-qubit state $a\ket{00}+b\ket{01}+c\ket{10}+d\ket{11}$ is $C=2|ad-bc|$, from 0 for a product to 1
for a maximally entangled state~\cite{wootters1998}. For pure two-qubit states $\mathcal N=C/2$.
\end{itemize}
An \emph{entanglement witness} is an observable $W$ with $\Tr(W\sigma)\ge0$ for every separable $\sigma$; a measured $\Tr(W\rho)<0$
certifies entanglement~\cite{guhne2009}. A given witness misses every entangled state for which $\Tr(W\rho)\ge0$.

\emph{LOCC.} In local operations and classical communication each party acts on its own system and the parties exchange only measurement
outcomes. LOCC maps take separable states to separable states, and no entanglement measure grows under them~\cite{vw,horodecki2009}.
Creating entanglement between two systems therefore needs a carrier that is itself quantum.

\emph{No-cloning and no-broadcasting.} No process copies two pure states that are neither equal nor orthogonal~\cite{wz,dieks}. A family
of states can be \emph{broadcast}, so that each of two readers ends up with the original, only if its members commute pairwise~\cite{barnum}.
A quantity that several readers hold in full is therefore diagonal in one basis, which is the core of quantum Darwinism~\cite{zurek}.

\emph{Worked example: a dephased mediator.} Qubits $A$ and $B$ start in $\ket+=(\ket0+\ket1)/\sqrt2$ and a mediator qubit $M$ in $\ket0$.
They interact only through $M$: a CNOT from $A$ to $M$, a controlled-$Z$ between $M$ and $B$, and the same CNOT again.
\begin{itemize}[nosep,leftmargin=*]
\item \emph{$M$ coherent.} The first CNOT gives $(\ket{0_A0_M}+\ket{1_A1_M})/\sqrt2$. The controlled-$Z$ turns $B$ into
$\ket-=(\ket0-\ket1)/\sqrt2$ in the second branch, and the second CNOT returns $M$ to $\ket0$. What remains is
$(\ket{0}_A\ket{+}_B+\ket{1}_A\ket{-}_B)/\sqrt2$, a maximally entangled state: $C=1$, $\mathcal N=1/2$.
\item \emph{$M$ dephased} in $\{\ket0,\ket1\}$ after the first CNOT. The state becomes the mixture
$\tfrac12(\ket{00}\bra{00}+\ket{11}\bra{11})_{AM}\otimes\ket+\bra+_B$, and the same two steps leave
$\tfrac12(\ket0\bra0_A\otimes\ket+\bra+_B+\ket1\bra1_A\otimes\ket-\bra-_B)$. This is separable: $\mathcal N=0$.
\end{itemize}
The dephased mediator carried one classical bit from $A$ to $B$; that is LOCC, and it correlates the two qubits without entangling them.
\end{primer}
}
\newcommand{\primerGR}{% GUID v6, Primer P3 (label box:gr): general relativity and the tetrad. Background only; it states no result of the paper.
% Place it before "The frame" (marker %PRIMER:gr in part1.tex): % GUID v6, Primer P3 (label box:gr): general relativity and the tetrad. Background only; it states no result of the paper.
% Place it before "The frame" (marker %PRIMER:gr in part1.tex): \input{primers/p3-gr}
\begin{primer}{General relativity and the tetrad}{box:gr}
\setlength{\parskip}{3pt plus 1pt}%
\emph{The metric.} General relativity describes gravity as the geometry of spacetime. A metric $g_{\mu\nu}$ gives the interval
$\dd s^2=g_{\mu\nu}\dd x^\mu\dd x^\nu$, with signature $(-,+,+,+)$, and Einstein's equation
$G_{\mu\nu}+\Lambda g_{\mu\nu}=(8\pi G/c^4)\,T_{\mu\nu}$ ties its curvature to the energy-momentum tensor $T_{\mu\nu}$ of matter and
radiation~\cite{wald1984}.

\emph{The tetrad.} A tetrad, or frame, is a set of four orthonormal one-forms $e^a=e^a{}_\mu\dd x^\mu$ ($a=0,\dots,3$) with
$g_{\mu\nu}=\eta_{ab}e^a{}_\mu e^b{}_\nu$ and $\eta=\operatorname{diag}(-1,1,1,1)$. At each point it can be rotated and boosted,
${e^a\to\Lambda^a{}_b(x)\,e^b}$, without changing $g_{\mu\nu}$: this is the \emph{local Lorentz freedom}. Of the tetrad's 16 components, 6 are
that freedom and 10 are the metric. General relativity can be written with the tetrad as its only field~\cite{aldrovandi}. A single unit
timelike vector field added to the metric is different: it picks out a rest frame at every point. That is Einstein--aether
theory~\cite{aether}, and its preferred-frame effects are tightly bounded, for example by the equal speeds of gravitational waves and
light~\cite{gw170817}.

\emph{Free fall.} Matter couples \emph{minimally} when it feels gravity only through $g_{\mu\nu}$, with partial derivatives replaced by
covariant ones. Then $\nabla^\mu T_{\mu\nu}=0$, and a small body follows a geodesic,
$\ddot x^\mu+\Gamma^\mu{}_{\alpha\beta}\dot x^\alpha\dot x^\beta=0$ (dots: derivatives by proper time), whatever it is made of. That is the
weak equivalence principle. Experiments compare the accelerations $a_1$, $a_2$ of two materials through the E\"otv\"os ratio
${\eta=2|a_1-a_2|/|a_1+a_2|}$~\cite{microscope}; with indices, $\eta$ is the Minkowski metric.

\emph{PPN parameters.} In weak fields, metric theories are compared through the parametrised post-Newtonian (PPN) expansion~\cite{will}.
With $U>0$ the Newtonian potential,
\[
g_{00}=-1+\frac{2U}{c^2}-2\beta\frac{U^2}{c^4},\qquad g_{ij}=\Big(1+2\gamma\frac{U}{c^2}\Big)\delta_{ij}.
\]
$\gamma$ says how much space curvature a unit mass produces; light bending and the Shapiro delay scale as $(1+\gamma)/2$. $\beta$ measures
the nonlinearity of gravity and enters the perihelion advance with $\gamma$. General relativity has $\gamma=\beta=1$.

\emph{The lapse.} Split into space and time, spacetime has a \emph{lapse} $N$: the proper time that a static clock records per unit of
coordinate time. In a weak potential $\Phi=-U$ it is $N=\sqrt{1+2\Phi/c^2}$, outside a mass $M$ it is $N=\sqrt{1-2GM/rc^2}$, and it
vanishes at a horizon.

\emph{Black holes.} A stationary black hole is fixed by three numbers: its mass $M$, angular momentum $J$ and charge
$Q$~\cite{israel,kerr}. The rotating hole (Kerr) has its horizon at $r_+=(GM/c^2)\big(1+\sqrt{1-a_*^2}\big)$, with spin
$a_*=cJ/GM^2\le1$~\cite{kerr}. Its \emph{surface gravity} $\kappa$ is the acceleration, as measured from far away, needed to hold a body
at the horizon; for a hole without spin $\kappa=c^4/4GM$. Quantum fields make the horizon radiate at $k_BT_H=\hbar\kappa/2\pi c$~\cite{hawking75}.
The \emph{area theorem}: with positive energy, the horizon area of classical general relativity never decreases~\cite{hawking1971}. It is
the second law of black-hole mechanics~\cite{bch}, with entropy $S=k_Bc^3A/4G\hbar$~\cite{bekenstein73}.
\end{primer}

\begin{primer}{General relativity and the tetrad}{box:gr}
\setlength{\parskip}{3pt plus 1pt}%
\emph{The metric.} General relativity describes gravity as the geometry of spacetime. A metric $g_{\mu\nu}$ gives the interval
$\dd s^2=g_{\mu\nu}\dd x^\mu\dd x^\nu$, with signature $(-,+,+,+)$, and Einstein's equation
$G_{\mu\nu}+\Lambda g_{\mu\nu}=(8\pi G/c^4)\,T_{\mu\nu}$ ties its curvature to the energy-momentum tensor $T_{\mu\nu}$ of matter and
radiation~\cite{wald1984}.

\emph{The tetrad.} A tetrad, or frame, is a set of four orthonormal one-forms $e^a=e^a{}_\mu\dd x^\mu$ ($a=0,\dots,3$) with
$g_{\mu\nu}=\eta_{ab}e^a{}_\mu e^b{}_\nu$ and $\eta=\operatorname{diag}(-1,1,1,1)$. At each point it can be rotated and boosted,
${e^a\to\Lambda^a{}_b(x)\,e^b}$, without changing $g_{\mu\nu}$: this is the \emph{local Lorentz freedom}. Of the tetrad's 16 components, 6 are
that freedom and 10 are the metric. General relativity can be written with the tetrad as its only field~\cite{aldrovandi}. A single unit
timelike vector field added to the metric is different: it picks out a rest frame at every point. That is Einstein--aether
theory~\cite{aether}, and its preferred-frame effects are tightly bounded, for example by the equal speeds of gravitational waves and
light~\cite{gw170817}.

\emph{Free fall.} Matter couples \emph{minimally} when it feels gravity only through $g_{\mu\nu}$, with partial derivatives replaced by
covariant ones. Then $\nabla^\mu T_{\mu\nu}=0$, and a small body follows a geodesic,
$\ddot x^\mu+\Gamma^\mu{}_{\alpha\beta}\dot x^\alpha\dot x^\beta=0$ (dots: derivatives by proper time), whatever it is made of. That is the
weak equivalence principle. Experiments compare the accelerations $a_1$, $a_2$ of two materials through the E\"otv\"os ratio
${\eta=2|a_1-a_2|/|a_1+a_2|}$~\cite{microscope}; with indices, $\eta$ is the Minkowski metric.

\emph{PPN parameters.} In weak fields, metric theories are compared through the parametrised post-Newtonian (PPN) expansion~\cite{will}.
With $U>0$ the Newtonian potential,
\[
g_{00}=-1+\frac{2U}{c^2}-2\beta\frac{U^2}{c^4},\qquad g_{ij}=\Big(1+2\gamma\frac{U}{c^2}\Big)\delta_{ij}.
\]
$\gamma$ says how much space curvature a unit mass produces; light bending and the Shapiro delay scale as $(1+\gamma)/2$. $\beta$ measures
the nonlinearity of gravity and enters the perihelion advance with $\gamma$. General relativity has $\gamma=\beta=1$.

\emph{The lapse.} Split into space and time, spacetime has a \emph{lapse} $N$: the proper time that a static clock records per unit of
coordinate time. In a weak potential $\Phi=-U$ it is $N=\sqrt{1+2\Phi/c^2}$, outside a mass $M$ it is $N=\sqrt{1-2GM/rc^2}$, and it
vanishes at a horizon.

\emph{Black holes.} A stationary black hole is fixed by three numbers: its mass $M$, angular momentum $J$ and charge
$Q$~\cite{israel,kerr}. The rotating hole (Kerr) has its horizon at $r_+=(GM/c^2)\big(1+\sqrt{1-a_*^2}\big)$, with spin
$a_*=cJ/GM^2\le1$~\cite{kerr}. Its \emph{surface gravity} $\kappa$ is the acceleration, as measured from far away, needed to hold a body
at the horizon; for a hole without spin $\kappa=c^4/4GM$. Quantum fields make the horizon radiate at $k_BT_H=\hbar\kappa/2\pi c$~\cite{hawking75}.
The \emph{area theorem}: with positive energy, the horizon area of classical general relativity never decreases~\cite{hawking1971}. It is
the second law of black-hole mechanics~\cite{bch}, with entropy $S=k_Bc^3A/4G\hbar$~\cite{bekenstein73}.
\end{primer}
}
\newcommand{\primerCQ}{% GUID v6, Primer P4 (label box:cq): classical-quantum dynamics and collapse. Background only; it states no result of the paper.
% Place it before "Consistent classical-quantum dynamics" (marker %PRIMER:cq in part1.tex): % GUID v6, Primer P4 (label box:cq): classical-quantum dynamics and collapse. Background only; it states no result of the paper.
% Place it before "Consistent classical-quantum dynamics" (marker %PRIMER:cq in part1.tex): \input{primers/p4-cq}
\begin{primer}{Classical-quantum dynamics and collapse}{box:cq}
\setlength{\parskip}{3pt plus 1pt}%
\emph{The classical-quantum (CQ) master equation.} A consistent coupling keeps the joint description linear, positive and
normalised, as probability theory demands. The joint state is a family $\varrho(z)$ of positive operators, one for each value
$z$ of the classical variables (here the frame), normalised as $\int\dd z\,\Tr\varrho(z)=1$~\cite{oppenheim}. For one classical variable
$z$ driven by one quantum observable $L$ (for gravity, the mass density), continuous Markovian dynamics has the form~\cite{osswd}
\[
\partial_t\varrho=-\frac i\hbar[H,\varrho]+D_0\Big(L\varrho L-\tfrac12\{L^2,\varrho\}\Big)-\frac{D_1}2\,\partial_z\{L,\varrho\}
+\frac{D_2}2\,\partial_z^2\varrho ,
\]
with three coefficients beyond quantum mechanics:
\begin{itemize}[nosep,leftmargin=*]
\item $D_0$, \emph{decoherence}: the quantum system loses coherence between states that source the classical field differently (for
gravity, different mass distributions);
\item $D_1$, \emph{back-reaction}: a drift of $z$ driven by the quantum state, that is, how matter pushes the frame;
\item $D_2$, \emph{diffusion}: random kicks of $z$, that is, noise in the frame.
\end{itemize}
Complete positivity requires $D_1^2\le4D_0D_2$ here, and $D_0>0$ wherever $D_1\ne0$; Oppenheim et al.\ give the general matrix form of
this trade-off~\cite[Eq.~(23)]{osswd}, which is Eq.~\eqref{eq:tradeoff} of this paper. In words: \emph{a push implies noise}. If quantum matter pushes a
classical field at all, the matter decoheres and the field diffuses, and the product of the two is bounded below by the square of the push,
$D_0D_2\ge D_1^2/4$.

\emph{Measurement and feedback.} An older family of consistent models makes gravity a channel. Each mass's position, or its mass density,
is weakly and continuously measured, and a classical force computed from the record acts back on the other masses (Kafri, Taylor and
Milburn, KTM~\cite{ktm}; Tilloy and Di\'osi~\cite{tilloydiosi}). Newton's law then holds on average, the measurement decoheres
superpositions, and the feedback adds noise. The masses talk only through a classical record, so the channel is LOCC
(Primer~\ref{box:qi}) and cannot entangle them. A weaker measurement decoheres less but gives a noisier record and so more heating; KTM find
both effects at least of the order of the gravitational coupling. The least decoherence of two masses over all measurement strengths is
what this paper calls the \emph{KTM minimum}, $\omega/2$, where $\omega$ is the rate at which a coherent field would build up the
entangling phase.

\emph{Gravity-related collapse.} Di\'osi and Penrose proposed that a superposition of two mass distributions $\rho_1$, $\rho_2$ decays at
the rate $\lambda=E_\Delta/\hbar$, where
\[
E_\Delta=\frac G2\int\!\!\int\dd^3x\,\dd^3y\,\frac{[\rho_1(\mathbf x)-\rho_2(\mathbf x)][\rho_1(\mathbf y)-\rho_2(\mathbf y)]}
{|\mathbf x-\mathbf y|}
\]
is the gravitational self-energy of their difference~\cite{diosi87,penrose96,bassi2013}. For a body of mass $M$ and size $R$ displaced by
more than its size, $\lambda\approx GM^2/\hbar R$: the \emph{Di\'osi--Penrose (DP) rate}. For point-like nuclei $E_\Delta$ diverges, so the
mass density must be smeared over a length $R_0$. The same noise makes charged matter radiate. An underground germanium detector sees no
such radiation, which excludes smearing at nuclear size and requires $R_0$ of order $10^{-10}$\,m or more~\cite{donadi,figurato}. In the
Newtonian limit of CQ gravity the DP kernel is one admissible choice of the decoherence term~\cite{osswd}.
\end{primer}

\begin{primer}{Classical-quantum dynamics and collapse}{box:cq}
\setlength{\parskip}{3pt plus 1pt}%
\emph{The classical-quantum (CQ) master equation.} A consistent coupling keeps the joint description linear, positive and
normalised, as probability theory demands. The joint state is a family $\varrho(z)$ of positive operators, one for each value
$z$ of the classical variables (here the frame), normalised as $\int\dd z\,\Tr\varrho(z)=1$~\cite{oppenheim}. For one classical variable
$z$ driven by one quantum observable $L$ (for gravity, the mass density), continuous Markovian dynamics has the form~\cite{osswd}
\[
\partial_t\varrho=-\frac i\hbar[H,\varrho]+D_0\Big(L\varrho L-\tfrac12\{L^2,\varrho\}\Big)-\frac{D_1}2\,\partial_z\{L,\varrho\}
+\frac{D_2}2\,\partial_z^2\varrho ,
\]
with three coefficients beyond quantum mechanics:
\begin{itemize}[nosep,leftmargin=*]
\item $D_0$, \emph{decoherence}: the quantum system loses coherence between states that source the classical field differently (for
gravity, different mass distributions);
\item $D_1$, \emph{back-reaction}: a drift of $z$ driven by the quantum state, that is, how matter pushes the frame;
\item $D_2$, \emph{diffusion}: random kicks of $z$, that is, noise in the frame.
\end{itemize}
Complete positivity requires $D_1^2\le4D_0D_2$ here, and $D_0>0$ wherever $D_1\ne0$; Oppenheim et al.\ give the general matrix form of
this trade-off~\cite[Eq.~(23)]{osswd}, which is Eq.~\eqref{eq:tradeoff} of this paper. In words: \emph{a push implies noise}. If quantum matter pushes a
classical field at all, the matter decoheres and the field diffuses, and the product of the two is bounded below by the square of the push,
$D_0D_2\ge D_1^2/4$.

\emph{Measurement and feedback.} An older family of consistent models makes gravity a channel. Each mass's position, or its mass density,
is weakly and continuously measured, and a classical force computed from the record acts back on the other masses (Kafri, Taylor and
Milburn, KTM~\cite{ktm}; Tilloy and Di\'osi~\cite{tilloydiosi}). Newton's law then holds on average, the measurement decoheres
superpositions, and the feedback adds noise. The masses talk only through a classical record, so the channel is LOCC
(Primer~\ref{box:qi}) and cannot entangle them. A weaker measurement decoheres less but gives a noisier record and so more heating; KTM find
both effects at least of the order of the gravitational coupling. The least decoherence of two masses over all measurement strengths is
what this paper calls the \emph{KTM minimum}, $\omega/2$, where $\omega$ is the rate at which a coherent field would build up the
entangling phase.

\emph{Gravity-related collapse.} Di\'osi and Penrose proposed that a superposition of two mass distributions $\rho_1$, $\rho_2$ decays at
the rate $\lambda=E_\Delta/\hbar$, where
\[
E_\Delta=\frac G2\int\!\!\int\dd^3x\,\dd^3y\,\frac{[\rho_1(\mathbf x)-\rho_2(\mathbf x)][\rho_1(\mathbf y)-\rho_2(\mathbf y)]}
{|\mathbf x-\mathbf y|}
\]
is the gravitational self-energy of their difference~\cite{diosi87,penrose96,bassi2013}. For a body of mass $M$ and size $R$ displaced by
more than its size, $\lambda\approx GM^2/\hbar R$: the \emph{Di\'osi--Penrose (DP) rate}. For point-like nuclei $E_\Delta$ diverges, so the
mass density must be smeared over a length $R_0$. The same noise makes charged matter radiate. An underground germanium detector sees no
such radiation, which excludes smearing at nuclear size and requires $R_0$ of order $10^{-10}$\,m or more~\cite{donadi,figurato}. In the
Newtonian limit of CQ gravity the DP kernel is one admissible choice of the decoherence term~\cite{osswd}.
\end{primer}
}
\newcommand{\primerHolo}{% GUID v6, Primer P5 (label box:holo): holography and quantum error correction. Background only; it states no result of the paper.
% Place it at the start of Section sec:holo (marker %PRIMER:holo in part1.tex): % GUID v6, Primer P5 (label box:holo): holography and quantum error correction. Background only; it states no result of the paper.
% Place it at the start of Section sec:holo (marker %PRIMER:holo in part1.tex): \input{primers/p5-holo}
\begin{primer}{Holography and quantum error correction}{box:holo}
\emph{Stabiliser codes.} A quantum error-correcting code stores $k$ logical qubits in $n$ physical ones so that some errors can be
undone. In a \emph{stabiliser code} the code space is the joint $+1$ eigenspace of $n-k$ commuting products of Pauli matrices, the
\emph{stabilisers}~\cite{gottesman1997}. Measuring them gives the \emph{syndrome}. An error that anticommutes with a stabiliser flips its
sign and so shows in the syndrome. The stabilisers commute with the logical operators, so the syndrome reveals nothing about the stored
content. A code of distance $d$, written $[[n,k,d]]$, corrects any $\lfloor(d-1)/2\rfloor$ unknown errors, or any $d-1$
\emph{erasures} (losses at known places).

\emph{The five-qubit code.} The smallest code that corrects every single-qubit error is $[[5,1,3]]$, with the stabiliser $XZZXI$ and its
cyclic shifts~\cite{laflamme,bennett1996}. It corrects any two erasures, so its logical qubit can be recovered from any three of the five
qubits and from no two: a $((3,5))$ quantum threshold scheme~\cite{cgl}.

\emph{Perfect tensors and the HaPPY code.} A tensor with $2m$ legs is \emph{perfect} if every split of its legs into $m$ and $m$ is, up to
normalisation, a unitary map. The encoding map of the five-qubit code, with its logical leg as a sixth, is such a tensor. Pastawski,
Yoshida, Harlow and Preskill tile the hyperbolic disc with pentagons, place one perfect tensor on each, contract neighbouring legs, and
leave one leg per tile in the bulk and the outermost legs on the boundary~\cite{happy}. The network is an isometry from the bulk to the
boundary, a toy model of the bulk-boundary correspondence of AdS/CFT~\cite{maldacena}.

\emph{Ryu--Takayanagi as a minimal cut.} In AdS/CFT the entropy of a boundary region $A$ is $S(A)=\mathrm{Area}(\gamma_A)/4G\hbar$ (with
$c=k_B=1$), where $\gamma_A$ is the bulk surface of least area anchored on the edge of $A$~\cite{rt}. In a network of perfect tensors the
entropy of $A$, in bits, is for many regions the number of qubit legs cut by the smallest cut that separates $A$ from the rest of the
boundary~\cite{happy}.

\emph{Reconstruction.} The bulk region between $A$ and $\gamma_A$ is the \emph{entanglement wedge} of $A$. Every bulk operator inside it can
be represented by an operator acting on $A$ alone~\cite{adh,dong2016}: the bulk is encoded redundantly, like a logical qubit. A code has
\emph{complementary recovery} when what cannot be reconstructed on $A$ can be reconstructed on the rest of the boundary.

\emph{The area operator.} Harlow showed that every code with complementary recovery obeys a quantum Ryu--Takayanagi formula
(Eq.~\eqref{eq:harlow} in Section~\ref{sec:holo}). Its area term is the expectation value of an \emph{area operator} that lies in the centre
of the algebra reconstructable on $A$: it commutes with everything that either side can do~\cite{harlow}. A \emph{fixed-area state} is an
eigenstate of this operator, and such states have a flat entanglement spectrum~\cite{dhm}. A general state of quantum gravity is a
superposition of areas.
\end{primer}

\begin{primer}{Holography and quantum error correction}{box:holo}
\emph{Stabiliser codes.} A quantum error-correcting code stores $k$ logical qubits in $n$ physical ones so that some errors can be
undone. In a \emph{stabiliser code} the code space is the joint $+1$ eigenspace of $n-k$ commuting products of Pauli matrices, the
\emph{stabilisers}~\cite{gottesman1997}. Measuring them gives the \emph{syndrome}. An error that anticommutes with a stabiliser flips its
sign and so shows in the syndrome. The stabilisers commute with the logical operators, so the syndrome reveals nothing about the stored
content. A code of distance $d$, written $[[n,k,d]]$, corrects any $\lfloor(d-1)/2\rfloor$ unknown errors, or any $d-1$
\emph{erasures} (losses at known places).

\emph{The five-qubit code.} The smallest code that corrects every single-qubit error is $[[5,1,3]]$, with the stabiliser $XZZXI$ and its
cyclic shifts~\cite{laflamme,bennett1996}. It corrects any two erasures, so its logical qubit can be recovered from any three of the five
qubits and from no two: a $((3,5))$ quantum threshold scheme~\cite{cgl}.

\emph{Perfect tensors and the HaPPY code.} A tensor with $2m$ legs is \emph{perfect} if every split of its legs into $m$ and $m$ is, up to
normalisation, a unitary map. The encoding map of the five-qubit code, with its logical leg as a sixth, is such a tensor. Pastawski,
Yoshida, Harlow and Preskill tile the hyperbolic disc with pentagons, place one perfect tensor on each, contract neighbouring legs, and
leave one leg per tile in the bulk and the outermost legs on the boundary~\cite{happy}. The network is an isometry from the bulk to the
boundary, a toy model of the bulk-boundary correspondence of AdS/CFT~\cite{maldacena}.

\emph{Ryu--Takayanagi as a minimal cut.} In AdS/CFT the entropy of a boundary region $A$ is $S(A)=\mathrm{Area}(\gamma_A)/4G\hbar$ (with
$c=k_B=1$), where $\gamma_A$ is the bulk surface of least area anchored on the edge of $A$~\cite{rt}. In a network of perfect tensors the
entropy of $A$, in bits, is for many regions the number of qubit legs cut by the smallest cut that separates $A$ from the rest of the
boundary~\cite{happy}.

\emph{Reconstruction.} The bulk region between $A$ and $\gamma_A$ is the \emph{entanglement wedge} of $A$. Every bulk operator inside it can
be represented by an operator acting on $A$ alone~\cite{adh,dong2016}: the bulk is encoded redundantly, like a logical qubit. A code has
\emph{complementary recovery} when what cannot be reconstructed on $A$ can be reconstructed on the rest of the boundary.

\emph{The area operator.} Harlow showed that every code with complementary recovery obeys a quantum Ryu--Takayanagi formula
(Eq.~\eqref{eq:harlow} in Section~\ref{sec:holo}). Its area term is the expectation value of an \emph{area operator} that lies in the centre
of the algebra reconstructable on $A$: it commutes with everything that either side can do~\cite{harlow}. A \emph{fixed-area state} is an
eigenstate of this operator, and such states have a flat entanglement spectrum~\cite{dhm}. A general state of quantum gravity is a
superposition of areas.
\end{primer}
}
\newcommand{\primerCosmo}{% GUID v6, Primer P6 (label box:cosmo): cosmology. Background only; it states no result of the paper.
% Place it at the start of Section sec:cosmo (marker %PRIMER:cosmo in part1.tex): % GUID v6, Primer P6 (label box:cosmo): cosmology. Background only; it states no result of the paper.
% Place it at the start of Section sec:cosmo (marker %PRIMER:cosmo in part1.tex): \input{primers/p6-cosmo}
\begin{primer}{Cosmology}{box:cosmo}
\emph{The expanding universe.} On large scales the universe is homogeneous and isotropic. Distances grow with a scale factor $a(t)$, light
from the past is redshifted by $1+z=a_0/a$, and the expansion rate $H=\dot a/a$ obeys, for a flat universe,
$H^2=8\pi G\rho/3+\Lambda c^2/3$. Each component's share is $\Omega_i=\rho_i/\rho_c$ with the critical density $\rho_c=3H^2/8\pi G$; the
value today is $H_0$.

\emph{$\Lambda$CDM.} The standard model of cosmology adds to the known particles a cold, collisionless \emph{dark matter} and a
cosmological constant $\Lambda$. Six parameters fix it: the densities of baryons and dark matter ($\Omega_bh^2$, $\Omega_ch^2$), the angular
size of the sound horizon, the optical depth to reionisation, and the amplitude $A_s$ and tilt $n_s$ of the primordial
fluctuations~\cite{planck18}.

\emph{Inflation.} An early epoch of nearly exponential expansion, $a\propto e^{Ht}$, flattens and smooths the universe and stretches
quantum fluctuations into the seeds of galaxies~\cite{guth}. It is driven by a scalar field, the \emph{inflaton}, with potential $V$, and
its length is counted in \emph{e-folds}, $N=\ln(a_{\rm end}/a_{\rm start})$. In \emph{slow roll} the field creeps down a flat potential;
with the reduced Planck mass $M_{\rm pl}$, the parameters $\epsilon=\tfrac12M_{\rm pl}^2(V'/V)^2$ and $\eta_V=M_{\rm pl}^2V''/V$ are small. The
fluctuations then have a nearly scale-invariant spectrum with tilt $n_s\approx1-6\epsilon+2\eta_V$, and gravitational waves with a
tensor-to-scalar ratio $r\approx16\epsilon$~\cite{planck2020x}.

\emph{Reheating and oscillons.} After inflation the inflaton oscillates about the minimum of its potential and decays into ordinary
particles, which \emph{reheats} the universe~\cite{kls}. In a broad class of models it first fragments into \emph{oscillons}, long-lived
localised lumps of the field~\cite{amin}.

\emph{The probes.} The cosmic microwave background (CMB) was released at recombination, $z\approx1100$. Its acoustic peaks are the imprint of
sound waves in the early plasma: their spacing fixes the sound horizon and their heights the densities~\cite{planck18}. The same sound
horizon, about 150\,Mpc today, is imprinted on the clustering of galaxies as \emph{baryon acoustic oscillations} (BAO), a standard ruler
for distance against redshift~\cite{desi2}.

\emph{Dark energy.} A cosmological constant has the equation of state $w=p/\rho c^2=-1$ at all times. DESI fits the
Chevallier--Polarski--Linder (CPL) form $w(a)=w_0+w_a(1-a)$~\cite{chevallier2001,linder2003}. A \emph{thawing} dark energy has $w_0>-1$
and $w_a<0$: it was close to $-1$ early and weakens today.

\emph{Data names and tensions.} P-ACT combines the CMB spectra of Planck and of the Atacama Cosmology Telescope (ACT, data release 6);
P-ACT-LB adds CMB lensing (L) and DESI BAO (B)~\cite{louis2025}. The \emph{$H_0$ tension} is the disagreement between the Hubble constant
inferred from the CMB with $\Lambda$CDM, near 67\,km\,s$^{-1}$\,Mpc$^{-1}$, and the value from a distance ladder of Cepheids and supernovae,
near 73~\cite{planck18,riess2022}. A ladder calibrated on the tip of the red-giant branch lies in between~\cite{freedman2025}.
\end{primer}

\begin{primer}{Cosmology}{box:cosmo}
\emph{The expanding universe.} On large scales the universe is homogeneous and isotropic. Distances grow with a scale factor $a(t)$, light
from the past is redshifted by $1+z=a_0/a$, and the expansion rate $H=\dot a/a$ obeys, for a flat universe,
$H^2=8\pi G\rho/3+\Lambda c^2/3$. Each component's share is $\Omega_i=\rho_i/\rho_c$ with the critical density $\rho_c=3H^2/8\pi G$; the
value today is $H_0$.

\emph{$\Lambda$CDM.} The standard model of cosmology adds to the known particles a cold, collisionless \emph{dark matter} and a
cosmological constant $\Lambda$. Six parameters fix it: the densities of baryons and dark matter ($\Omega_bh^2$, $\Omega_ch^2$), the angular
size of the sound horizon, the optical depth to reionisation, and the amplitude $A_s$ and tilt $n_s$ of the primordial
fluctuations~\cite{planck18}.

\emph{Inflation.} An early epoch of nearly exponential expansion, $a\propto e^{Ht}$, flattens and smooths the universe and stretches
quantum fluctuations into the seeds of galaxies~\cite{guth}. It is driven by a scalar field, the \emph{inflaton}, with potential $V$, and
its length is counted in \emph{e-folds}, $N=\ln(a_{\rm end}/a_{\rm start})$. In \emph{slow roll} the field creeps down a flat potential;
with the reduced Planck mass $M_{\rm pl}$, the parameters $\epsilon=\tfrac12M_{\rm pl}^2(V'/V)^2$ and $\eta_V=M_{\rm pl}^2V''/V$ are small. The
fluctuations then have a nearly scale-invariant spectrum with tilt $n_s\approx1-6\epsilon+2\eta_V$, and gravitational waves with a
tensor-to-scalar ratio $r\approx16\epsilon$~\cite{planck2020x}.

\emph{Reheating and oscillons.} After inflation the inflaton oscillates about the minimum of its potential and decays into ordinary
particles, which \emph{reheats} the universe~\cite{kls}. In a broad class of models it first fragments into \emph{oscillons}, long-lived
localised lumps of the field~\cite{amin}.

\emph{The probes.} The cosmic microwave background (CMB) was released at recombination, $z\approx1100$. Its acoustic peaks are the imprint of
sound waves in the early plasma: their spacing fixes the sound horizon and their heights the densities~\cite{planck18}. The same sound
horizon, about 150\,Mpc today, is imprinted on the clustering of galaxies as \emph{baryon acoustic oscillations} (BAO), a standard ruler
for distance against redshift~\cite{desi2}.

\emph{Dark energy.} A cosmological constant has the equation of state $w=p/\rho c^2=-1$ at all times. DESI fits the
Chevallier--Polarski--Linder (CPL) form $w(a)=w_0+w_a(1-a)$~\cite{chevallier2001,linder2003}. A \emph{thawing} dark energy has $w_0>-1$
and $w_a<0$: it was close to $-1$ early and weakens today.

\emph{Data names and tensions.} P-ACT combines the CMB spectra of Planck and of the Atacama Cosmology Telescope (ACT, data release 6);
P-ACT-LB adds CMB lensing (L) and DESI BAO (B)~\cite{louis2025}. The \emph{$H_0$ tension} is the disagreement between the Hubble constant
inferred from the CMB with $\Lambda$CDM, near 67\,km\,s$^{-1}$\,Mpc$^{-1}$, and the value from a distance ladder of Cepheids and supernovae,
near 73~\cite{planck18,riess2022}. A ladder calibrated on the tip of the red-giant branch lies in between~\cite{freedman2025}.
\end{primer}
}
\newcommand{\primerGalaxy}{% GUID v6, Primer P7 (label box:galaxy): galaxy dynamics. Background only; it states no result of the paper.
% Place it at the start of Section sec:sparc (or at Section sec:darkmatter if the lead prefers): % GUID v6, Primer P7 (label box:galaxy): galaxy dynamics. Background only; it states no result of the paper.
% Place it at the start of Section sec:sparc (or at Section sec:darkmatter if the lead prefers): \input{primers/p7-galaxy}
\begin{primer}{Galaxy dynamics}{box:galaxy}
\setlength{\parskip}{3pt plus 1pt}%
\emph{Rotation curves.} Gas and stars in a disc galaxy move on nearly circular orbits. Their speed $V$ against radius $R$ is the
\emph{rotation curve}, and the centripetal acceleration is $g_{\rm obs}=V^2/R$. Far out, most curves stay flat where the visible mass alone
would make them fall as $R^{-1/2}$: the missing-mass problem.

\emph{Baryonic mass models.} The visible, or \emph{baryonic}, mass has three parts, each with its own contribution to $V^2$, computed from
its observed surface density: the gas (atomic hydrogen with its helium), the stellar disc and the central bulge. Starlight becomes stellar
mass through a mass-to-light ratio $\Upsilon$; at a wavelength of 3.6\,\textmu m it is near 0.5 for discs and 0.7 for
bulges~\cite{mcgaugh2016}. The sum of the three, divided by $R$, is the baryonic acceleration $g_{\rm bar}$.

\emph{Two empirical laws.} The \emph{baryonic Tully--Fisher relation} (BTFR): a galaxy's baryonic mass grows as the fourth power of its flat
rotation speed, $M_{\rm bar}\propto V_{\rm flat}^4$, with small scatter~\cite{lelli2016a}. The \emph{radial acceleration relation} (RAR): at
every radius of every galaxy, $g_{\rm obs}$ is one function of $g_{\rm bar}$. It is close to $g_{\rm bar}$ at high acceleration and close to
$\sqrt{g_{\rm bar}\,g_\dagger}$ below a scale $g_\dagger\approx10^{-10}$\,m\,s$^{-2}$~\cite{mcgaugh2016,lelli2017}.

\emph{MOND.} In Milgrom's modified Newtonian dynamics, $g_{\rm obs}=\nu(g_{\rm bar}/a_0)\,g_{\rm bar}$, where the new constant
$a_0$ is about $1.2\times10^{-10}$\,m\,s$^{-2}$~\cite{milgrom}. The \emph{interpolating function} $\nu(y)$ tends to 1 at high acceleration
($y\gg1$) and to $y^{-1/2}$ at low; the deep regime gives $V^4=GMa_0$, the BTFR. Common choices are the ``simple'' function
$\nu(y)=\tfrac12+(\tfrac14+1/y)^{1/2}$ and the ``standard'' one, and the RAR's fitting function is another~\cite{famaey2012}.

\emph{Dark-matter halos.} In cold dark matter each galaxy sits in a halo. Simulated halos follow the NFW profile
\[
\rho(r)=\frac{\rho_s}{(r/r_s)(1+r/r_s)^2}
\]
of Navarro, Frenk and White~\cite{navarro1996}. Its mass $M_{200}$ is counted inside the radius $r_{200}$ that encloses 200 times the
critical density, and its concentration $c=r_{200}/r_s$ falls slowly with mass~\cite{dutton2014}. At the centre the density rises as
$r^{-1}$, a \emph{cusp}. Many dwarf and low-surface-brightness galaxies prefer a \emph{core} of nearly constant density, as in the Burkert
profile $\rho=\rho_0r_0^3/[(r+r_0)(r^2+r_0^2)]$~\cite{burkert1995}; this mismatch is the cusp-core problem~\cite{deblok2010}. The
\emph{stellar-to-halo relation}, found by matching the abundances of galaxies and halos, gives the stellar mass a halo of given mass hosts;
the ratio peaks at a few per cent near $10^{12}$ solar masses~\cite{moster2013}.

\emph{SPARC.} The Spitzer Photometry and Accurate Rotation Curves sample holds 175 disc galaxies with 3.6\,\textmu m photometry and
rotation curves from atomic-hydrogen and H$\alpha$ data, with mass models for gas, disc and bulge~\cite{lelli2016}.
\end{primer}

\begin{primer}{Galaxy dynamics}{box:galaxy}
\setlength{\parskip}{3pt plus 1pt}%
\emph{Rotation curves.} Gas and stars in a disc galaxy move on nearly circular orbits. Their speed $V$ against radius $R$ is the
\emph{rotation curve}, and the centripetal acceleration is $g_{\rm obs}=V^2/R$. Far out, most curves stay flat where the visible mass alone
would make them fall as $R^{-1/2}$: the missing-mass problem.

\emph{Baryonic mass models.} The visible, or \emph{baryonic}, mass has three parts, each with its own contribution to $V^2$, computed from
its observed surface density: the gas (atomic hydrogen with its helium), the stellar disc and the central bulge. Starlight becomes stellar
mass through a mass-to-light ratio $\Upsilon$; at a wavelength of 3.6\,\textmu m it is near 0.5 for discs and 0.7 for
bulges~\cite{mcgaugh2016}. The sum of the three, divided by $R$, is the baryonic acceleration $g_{\rm bar}$.

\emph{Two empirical laws.} The \emph{baryonic Tully--Fisher relation} (BTFR): a galaxy's baryonic mass grows as the fourth power of its flat
rotation speed, $M_{\rm bar}\propto V_{\rm flat}^4$, with small scatter~\cite{lelli2016a}. The \emph{radial acceleration relation} (RAR): at
every radius of every galaxy, $g_{\rm obs}$ is one function of $g_{\rm bar}$. It is close to $g_{\rm bar}$ at high acceleration and close to
$\sqrt{g_{\rm bar}\,g_\dagger}$ below a scale $g_\dagger\approx10^{-10}$\,m\,s$^{-2}$~\cite{mcgaugh2016,lelli2017}.

\emph{MOND.} In Milgrom's modified Newtonian dynamics, $g_{\rm obs}=\nu(g_{\rm bar}/a_0)\,g_{\rm bar}$, where the new constant
$a_0$ is about $1.2\times10^{-10}$\,m\,s$^{-2}$~\cite{milgrom}. The \emph{interpolating function} $\nu(y)$ tends to 1 at high acceleration
($y\gg1$) and to $y^{-1/2}$ at low; the deep regime gives $V^4=GMa_0$, the BTFR. Common choices are the ``simple'' function
$\nu(y)=\tfrac12+(\tfrac14+1/y)^{1/2}$ and the ``standard'' one, and the RAR's fitting function is another~\cite{famaey2012}.

\emph{Dark-matter halos.} In cold dark matter each galaxy sits in a halo. Simulated halos follow the NFW profile
\[
\rho(r)=\frac{\rho_s}{(r/r_s)(1+r/r_s)^2}
\]
of Navarro, Frenk and White~\cite{navarro1996}. Its mass $M_{200}$ is counted inside the radius $r_{200}$ that encloses 200 times the
critical density, and its concentration $c=r_{200}/r_s$ falls slowly with mass~\cite{dutton2014}. At the centre the density rises as
$r^{-1}$, a \emph{cusp}. Many dwarf and low-surface-brightness galaxies prefer a \emph{core} of nearly constant density, as in the Burkert
profile $\rho=\rho_0r_0^3/[(r+r_0)(r^2+r_0^2)]$~\cite{burkert1995}; this mismatch is the cusp-core problem~\cite{deblok2010}. The
\emph{stellar-to-halo relation}, found by matching the abundances of galaxies and halos, gives the stellar mass a halo of given mass hosts;
the ratio peaks at a few per cent near $10^{12}$ solar masses~\cite{moster2013}.

\emph{SPARC.} The Spitzer Photometry and Accurate Rotation Curves sample holds 175 disc galaxies with 3.6\,\textmu m photometry and
rotation curves from atomic-hydrogen and H$\alpha$ data, with mass models for gas, disc and bulge~\cite{lelli2016}.
\end{primer}
}
\newcommand{\primerStats}{% GUID v6, Primer P8 (label box:stats): the statistics behind the scoring rules of Section sec:scoring. Background only; it states
% no result of the paper and does not repeat the rules themselves. Place it in Section sec:scoring: % GUID v6, Primer P8 (label box:stats): the statistics behind the scoring rules of Section sec:scoring. Background only; it states
% no result of the paper and does not repeat the rules themselves. Place it in Section sec:scoring: \input{primers/p8-stats}
\begin{primer}{Statistics}{box:stats}
\emph{Pulls and 95\%.} A measurement $x\pm\sigma$ usually means a Gaussian error with standard deviation $\sigma$. If a prediction
$x_G\pm\sigma_G$ is right and the two errors are independent, the difference $x-x_G$ is Gaussian with variance $\sigma^2+\sigma_G^2$, so the
\emph{pull} $p=(x-x_G)/\sqrt{\sigma^2+\sigma_G^2}$ is a standard normal number. It falls outside $\pm1.96$ with probability 5\% (two-sided),
so $|p|<1.96$ is a test at 95\%, and a published 95\% half-width $b$ corresponds to $\sigma=b/1.96$. \emph{Worked pull:} a model predicts
$10.0\pm0.3$ and an experiment finds $11.0\pm0.4$. Then $p=1.0/\sqrt{0.4^2+0.3^2}=1.0/0.5=2.0$: a FAIL at 95\%, although a right model
gives $|p|\ge2.0$ in 4.6\% of independent tests. For $n$ independent tests, about $0.05\,n$ FAILs are expected by chance, and the sum of the
squared pulls follows a $\chi^2$ distribution with $n$ degrees of freedom.

\emph{BIC and $\Delta$BIC.} A model with $k$ fitted parameters that reaches the maximum likelihood $\mathcal L_{\max}$ on $N$ data points has
$\mathrm{BIC}=-2\ln\mathcal L_{\max}+k\ln N$~\cite{schwarz1978}. The penalty $k\ln N$ charges each extra parameter. For large $N$ and
independent points, the difference $\Delta$BIC of two models approximates $-2$ times the logarithm of their Bayes factor, the ratio of the
probabilities of the data under the two models. A $\Delta$BIC of 2 to 6 is then positive evidence, 6 to 10 strong and above 10 very
strong~\cite{kassraftery1995}. The approximation assumes $N$ independent points. If the points come in groups that share an error, such
as the distance of one galaxy, the effective number of independent data is closer to the number of groups. A point-wise $\Delta$BIC then
overstates the evidence, roughly in proportion to the number of points per group.

\emph{The bootstrap.} To learn how much a fitted value would vary in a new sample, one draws many new samples of the same size from the
data themselves, with replacement, and refits each~\cite{efron1979}. The spread of the refitted values estimates the uncertainty; the 16th
to 84th percentiles bracket a range comparable to $\pm1\sigma$. The units resampled must be the independent ones, for rotation curves whole
galaxies, not single points.

\emph{Many tests.} With $m$ tests at 5\% each, the chance of at least one false positive, the \emph{family-wise error rate}, grows towards
$1-0.95^m$. Bonferroni's remedy tests each at $0.05/m$. Holm's method orders the $p$-values, $p_{(1)}\le\dots\le p_{(m)}$, and rejects
$p_{(i)}$ while $p_{(i)}\le0.05/(m-i+1)$; it keeps the family-wise error rate at 5\% under any dependence and never rejects less than
Bonferroni~\cite{holm1979}. Fisher's method combines $n$ \emph{independent} $p$-values: under the null each $p_i$ is uniform, so
$-2\ln p_i$ follows $\chi^2$ with 2 degrees of freedom, and $X^2=-2\sum_i\ln p_i$ follows $\chi^2$ with $2n$~\cite{fisher1932}. It responds
strongly to a few very small $p_i$.

\emph{Look-elsewhere and preregistration.} A search over many masses, scales or candidate values finds a small \emph{local} $p$-value
somewhere by chance. The \emph{global} $p$-value, the chance of so small a value anywhere in the search, is larger, roughly the local one
times a trials factor~\cite{grossvitells2010}. A test is confirmatory only if its hypothesis, data, statistic and decision rule were
fixed before the outcome was seen; \emph{preregistration} writes them down in a dated record beforehand, and every choice made afterwards
is then visible as exploratory~\cite{nosek2018}.
\end{primer}

\begin{primer}{Statistics}{box:stats}
\emph{Pulls and 95\%.} A measurement $x\pm\sigma$ usually means a Gaussian error with standard deviation $\sigma$. If a prediction
$x_G\pm\sigma_G$ is right and the two errors are independent, the difference $x-x_G$ is Gaussian with variance $\sigma^2+\sigma_G^2$, so the
\emph{pull} $p=(x-x_G)/\sqrt{\sigma^2+\sigma_G^2}$ is a standard normal number. It falls outside $\pm1.96$ with probability 5\% (two-sided),
so $|p|<1.96$ is a test at 95\%, and a published 95\% half-width $b$ corresponds to $\sigma=b/1.96$. \emph{Worked pull:} a model predicts
$10.0\pm0.3$ and an experiment finds $11.0\pm0.4$. Then $p=1.0/\sqrt{0.4^2+0.3^2}=1.0/0.5=2.0$: a FAIL at 95\%, although a right model
gives $|p|\ge2.0$ in 4.6\% of independent tests. For $n$ independent tests, about $0.05\,n$ FAILs are expected by chance, and the sum of the
squared pulls follows a $\chi^2$ distribution with $n$ degrees of freedom.

\emph{BIC and $\Delta$BIC.} A model with $k$ fitted parameters that reaches the maximum likelihood $\mathcal L_{\max}$ on $N$ data points has
$\mathrm{BIC}=-2\ln\mathcal L_{\max}+k\ln N$~\cite{schwarz1978}. The penalty $k\ln N$ charges each extra parameter. For large $N$ and
independent points, the difference $\Delta$BIC of two models approximates $-2$ times the logarithm of their Bayes factor, the ratio of the
probabilities of the data under the two models. A $\Delta$BIC of 2 to 6 is then positive evidence, 6 to 10 strong and above 10 very
strong~\cite{kassraftery1995}. The approximation assumes $N$ independent points. If the points come in groups that share an error, such
as the distance of one galaxy, the effective number of independent data is closer to the number of groups. A point-wise $\Delta$BIC then
overstates the evidence, roughly in proportion to the number of points per group.

\emph{The bootstrap.} To learn how much a fitted value would vary in a new sample, one draws many new samples of the same size from the
data themselves, with replacement, and refits each~\cite{efron1979}. The spread of the refitted values estimates the uncertainty; the 16th
to 84th percentiles bracket a range comparable to $\pm1\sigma$. The units resampled must be the independent ones, for rotation curves whole
galaxies, not single points.

\emph{Many tests.} With $m$ tests at 5\% each, the chance of at least one false positive, the \emph{family-wise error rate}, grows towards
$1-0.95^m$. Bonferroni's remedy tests each at $0.05/m$. Holm's method orders the $p$-values, $p_{(1)}\le\dots\le p_{(m)}$, and rejects
$p_{(i)}$ while $p_{(i)}\le0.05/(m-i+1)$; it keeps the family-wise error rate at 5\% under any dependence and never rejects less than
Bonferroni~\cite{holm1979}. Fisher's method combines $n$ \emph{independent} $p$-values: under the null each $p_i$ is uniform, so
$-2\ln p_i$ follows $\chi^2$ with 2 degrees of freedom, and $X^2=-2\sum_i\ln p_i$ follows $\chi^2$ with $2n$~\cite{fisher1932}. It responds
strongly to a few very small $p_i$.

\emph{Look-elsewhere and preregistration.} A search over many masses, scales or candidate values finds a small \emph{local} $p$-value
somewhere by chance. The \emph{global} $p$-value, the chance of so small a value anywhere in the search, is larger, roughly the local one
times a trials factor~\cite{grossvitells2010}. A test is confirmatory only if its hypothesis, data, statistic and decision rule were
fixed before the outcome was seen; \emph{preregistration} writes them down in a dated record beforehand, and every choice made afterwards
is then visible as exploratory~\cite{nosek2018}.
\end{primer}
}
\newcommand{\primerParticle}{% GUID v6, Primer P9 (label box:particle): particle physics. Background only; it states no result of the paper.
% Place it before "The spin ladder" (marker %PRIMER:particle in part1.tex): % GUID v6, Primer P9 (label box:particle): particle physics. Background only; it states no result of the paper.
% Place it before "The spin ladder" (marker %PRIMER:particle in part1.tex): \input{primers/p9-particle}
\begin{primer}{Particle physics}{box:particle}
\setlength{\parskip}{3pt plus 1pt}%
\emph{The Higgs mechanism.} The Standard Model is a gauge theory of $SU(3)\times SU(2)\times U(1)$: colour for the strong force, weak isospin
and hypercharge for the electroweak force~\cite{glashow,weinberg67,salam}. The Higgs field is a doublet whose potential has its minimum
away from zero, at a vacuum value $v=246$\,GeV~\cite{higgs64,be64}. In that vacuum the W and Z bosons acquire masses proportional to $v$,
while one combination of the electroweak fields, the photon, stays massless. Quarks and charged leptons get their masses from
\emph{Yukawa couplings} $y_f$ to the doublet, $m_f=y_fv/\sqrt2$. The Higgs boson therefore couples to each fermion with strength $m_f/v$
and to the W and Z with $2m_V^2/v$.

\emph{The $\kappa$ framework.} Measurements of Higgs production and decay are summarised by \emph{coupling modifiers} $\kappa_i$, which
scale each Standard Model coupling, so that a rate involving particle $i$ scales as $\kappa_i^2$; the Standard Model has every
$\kappa_i=1$~\cite{david2012}.

\emph{Schemes and running.} A coupling or mass in quantum field theory is defined only together with a \emph{renormalisation scheme} and
an energy scale $\mu$. In the $\overline{\text{MS}}$ scheme only the divergent part of each loop is subtracted, and parameters run with
$\mu$. The \emph{on-shell} scheme instead fixes masses at the poles of the propagators and $\alpha$ at zero momentum transfer. The same
quantity can differ between the two by more than its measurement error, so every comparison states its scheme~\cite{erler2017,navas2024}.
The dependence on $\mu$ obeys the renormalisation-group equations (RGEs). At one loop
$\dd\alpha_i^{-1}/\dd\ln\mu=-b_i/2\pi$, with $(b_1,b_2,b_3)=(41/10,-19/6,-7)$ in the Standard Model (hypercharge normalised as in grand
unification). The strong coupling therefore weakens at high energy, and the hypercharge coupling grows. Precision work runs at two or
three loops~\cite{buttazzo2013}.

\emph{Dirac or Majorana.} A \emph{Dirac} mass joins a left-handed neutrino to a right-handed partner and conserves lepton number $L$. A
\emph{Majorana} mass joins the neutrino to its own antiparticle and changes $L$ by two. The lowest-dimension operator of Standard Model
fields that gives neutrinos mass is Weinberg's dimension-five operator, $(c/\Lambda)\,(LH)(LH)$, with $L$ the lepton doublet; in the Higgs
vacuum it gives a Majorana mass of order $cv^2/\Lambda$, small when the new scale $\Lambda$ is large~\cite{weinberg1979}. Majorana masses
can therefore also come from the Higgs vacuum.

\emph{Neutrinoless double-beta decay} ($0\nu\beta\beta$) turns two neutrons of a nucleus into two protons and two electrons with no
neutrinos. It changes $L$ by two and can occur only if neutrinos have a Majorana mass. Through light-neutrino exchange its rate is
proportional to $m_{\beta\beta}^2$, with the effective mass $m_{\beta\beta}=|\sum_iU_{ei}^2m_i|$ set by the mixing matrix $U$ and the
masses $m_i$~\cite{navas2024}. Xenon and germanium experiments search for it~\cite{kamlandzen2025,legend2025}.

\emph{Grand unification.} A grand unified theory (GUT) embeds the three gauge groups in one simple group, the smallest being
$SU(5)$~\cite{georgiglashow}, with a single coupling above a scale $M_X$. With the Standard Model's particles alone, the three running
couplings do not meet at one point~\cite{amaldi}. A GUT has two generic consequences. Its heavy gauge bosons let protons decay, with a
lifetime of order $M_X^4/(\alpha_G^2m_p^5)$; Super-Kamiokande's bound on $p\to e^+\pi^0$ excludes the minimal $SU(5)$
model~\cite{superk}. And its breaking leaves magnetic monopoles of mass near $M_X/\alpha_G$, about one per causal volume, which would
overfill the universe unless inflation dilutes them~\cite{preskill,guth}.
\end{primer}

\begin{primer}{Particle physics}{box:particle}
\setlength{\parskip}{3pt plus 1pt}%
\emph{The Higgs mechanism.} The Standard Model is a gauge theory of $SU(3)\times SU(2)\times U(1)$: colour for the strong force, weak isospin
and hypercharge for the electroweak force~\cite{glashow,weinberg67,salam}. The Higgs field is a doublet whose potential has its minimum
away from zero, at a vacuum value $v=246$\,GeV~\cite{higgs64,be64}. In that vacuum the W and Z bosons acquire masses proportional to $v$,
while one combination of the electroweak fields, the photon, stays massless. Quarks and charged leptons get their masses from
\emph{Yukawa couplings} $y_f$ to the doublet, $m_f=y_fv/\sqrt2$. The Higgs boson therefore couples to each fermion with strength $m_f/v$
and to the W and Z with $2m_V^2/v$.

\emph{The $\kappa$ framework.} Measurements of Higgs production and decay are summarised by \emph{coupling modifiers} $\kappa_i$, which
scale each Standard Model coupling, so that a rate involving particle $i$ scales as $\kappa_i^2$; the Standard Model has every
$\kappa_i=1$~\cite{david2012}.

\emph{Schemes and running.} A coupling or mass in quantum field theory is defined only together with a \emph{renormalisation scheme} and
an energy scale $\mu$. In the $\overline{\text{MS}}$ scheme only the divergent part of each loop is subtracted, and parameters run with
$\mu$. The \emph{on-shell} scheme instead fixes masses at the poles of the propagators and $\alpha$ at zero momentum transfer. The same
quantity can differ between the two by more than its measurement error, so every comparison states its scheme~\cite{erler2017,navas2024}.
The dependence on $\mu$ obeys the renormalisation-group equations (RGEs). At one loop
$\dd\alpha_i^{-1}/\dd\ln\mu=-b_i/2\pi$, with $(b_1,b_2,b_3)=(41/10,-19/6,-7)$ in the Standard Model (hypercharge normalised as in grand
unification). The strong coupling therefore weakens at high energy, and the hypercharge coupling grows. Precision work runs at two or
three loops~\cite{buttazzo2013}.

\emph{Dirac or Majorana.} A \emph{Dirac} mass joins a left-handed neutrino to a right-handed partner and conserves lepton number $L$. A
\emph{Majorana} mass joins the neutrino to its own antiparticle and changes $L$ by two. The lowest-dimension operator of Standard Model
fields that gives neutrinos mass is Weinberg's dimension-five operator, $(c/\Lambda)\,(LH)(LH)$, with $L$ the lepton doublet; in the Higgs
vacuum it gives a Majorana mass of order $cv^2/\Lambda$, small when the new scale $\Lambda$ is large~\cite{weinberg1979}. Majorana masses
can therefore also come from the Higgs vacuum.

\emph{Neutrinoless double-beta decay} ($0\nu\beta\beta$) turns two neutrons of a nucleus into two protons and two electrons with no
neutrinos. It changes $L$ by two and can occur only if neutrinos have a Majorana mass. Through light-neutrino exchange its rate is
proportional to $m_{\beta\beta}^2$, with the effective mass $m_{\beta\beta}=|\sum_iU_{ei}^2m_i|$ set by the mixing matrix $U$ and the
masses $m_i$~\cite{navas2024}. Xenon and germanium experiments search for it~\cite{kamlandzen2025,legend2025}.

\emph{Grand unification.} A grand unified theory (GUT) embeds the three gauge groups in one simple group, the smallest being
$SU(5)$~\cite{georgiglashow}, with a single coupling above a scale $M_X$. With the Standard Model's particles alone, the three running
couplings do not meet at one point~\cite{amaldi}. A GUT has two generic consequences. Its heavy gauge bosons let protons decay, with a
lifetime of order $M_X^4/(\alpha_G^2m_p^5)$; Super-Kamiokande's bound on $p\to e^+\pi^0$ excludes the minimal $SU(5)$
model~\cite{superk}. And its breaking leaves magnetic monopoles of mass near $M_X/\alpha_G$, about one per causal volume, which would
overfill the universe unless inflation dilutes them~\cite{preskill,guth}.
\end{primer}
}
\newcommand{\primerHadron}{% GUID v6, Primer P10 (label box:hadron): hadron spectroscopy. Background only; it states no result of the paper.
% Place it before "Hadrons as trees of dualities" (marker %PRIMER:hadron in part1.tex): % GUID v6, Primer P10 (label box:hadron): hadron spectroscopy. Background only; it states no result of the paper.
% Place it before "Hadrons as trees of dualities" (marker %PRIMER:hadron in part1.tex): \input{primers/p10-hadron}
\begin{primer}{Hadron spectroscopy}{box:hadron}
\emph{Colour.} Quarks carry the colour charge of $SU(3)$: a quark is a triplet $\mathbf3$, an antiquark an antitriplet $\mathbf{\bar3}$.
Only colour \emph{singlets} exist as free particles, the hadrons~\cite{navas2024}. The products
$\mathbf3\otimes\mathbf{\bar3}=\mathbf1\oplus\mathbf8$ and $\mathbf3\otimes\mathbf3=\mathbf{\bar3}\oplus\mathbf6$ give the ordinary
hadrons: a meson $q\bar q$ in the singlet, and a baryon $qqq$, whose first two quarks couple to $\mathbf{\bar3}$ and then with the third to
a singlet. One-gluon exchange attracts in the singlet $q\bar q$ and, half as strongly, in the $\mathbf{\bar3}$ $qq$, and repels in the
$\mathbf8$ and the $\mathbf6$.

\emph{Diquarks.} A \emph{diquark} is a pair of quarks in the attractive $\mathbf{\bar3}$. The ``good'' diquark, with spin 0 and a flavour-
antisymmetric pair such as $[ud]$, binds most; the ``bad'' one has spin 1~\cite{jaffe2005}.

\emph{Compact or molecular.} A hadron with more than three quarks, an \emph{exotic}, can be \emph{compact}, for example a diquark and an
antidiquark bound by colour forces, $[qq][\bar q\bar q]$~\cite{maiani2005}. Or it can be a \emph{hadronic molecule}: two colour-singlet
hadrons held together by the residual strong force, like the deuteron~\cite{guo2018}. Molecules sit just below the threshold
$m_1+m_2$ of their two constituents.

\emph{The colour-spin interaction.} In the constituent-quark model a hadron's mass is the sum of effective quark masses plus a hyperfine
term proportional to $(\boldsymbol\lambda_i\cdot\boldsymbol\lambda_j)(\boldsymbol\sigma_i\cdot\boldsymbol\sigma_j)/m_im_j$ for each pair,
where the $\boldsymbol\lambda$ are colour and the $\boldsymbol\sigma$ spin matrices. This chromomagnetic interaction of De R\'ujula, Georgi
and Glashow explains, for example, why the $\Delta$ is heavier than the nucleon~\cite{derujula1975}. Karliner and Rosner (KR) refined it
with separate effective quark masses in mesons and in baryons, the baryon masses for every constituent of a compact multiquark, and binding
terms for heavy pairs with $B(QQ)=B(Q\bar Q)/2$; with it they predicted the doubly charmed baryon and doubly heavy
tetraquarks~\cite{karliner2014,karliner2017}.

\emph{Thresholds and cusps.} Where a new two-hadron channel opens, at $M=m_1+m_2$, the scattering amplitude has a branch point. A
\emph{cusp} can then appear in a mass spectrum without any state, and a state that couples to the channel is pulled towards the threshold.
Many exotic candidates lie within a few MeV of a threshold~\cite{guo2018}.

\emph{Names.} The Particle Data Group names hadrons by their quantum numbers, not by their structure~\cite{navas2024}; the
$X(3872)$, for example, is listed as $\chi_{c1}(3872)$. LHCb's convention for exotics~\cite{gershon2022} uses $T$ for four-quark and $P$
for five-quark states, with subscripts for heavy flavour, $\psi$ standing for hidden charm $c\bar c$. So $T_{cc}^+$ is $cc\bar u\bar d$,
seen just below the $D^{*+}D^0$ threshold~\cite{lhcb2022h}, and the hidden-charm pentaquark $P_c(4312)^+$~\cite{lhcb2019h} becomes
$P_\psi^N(4312)^+$.

\emph{String junctions.} In a baryon the three colour flux tubes meet at a point, a \emph{string junction}~\cite{rossiveneziano1977}. A
compact tetraquark has two junctions and a compact pentaquark three, and some models add a constant energy per
junction~\cite{karliner2017a}.
\end{primer}

\begin{primer}{Hadron spectroscopy}{box:hadron}
\emph{Colour.} Quarks carry the colour charge of $SU(3)$: a quark is a triplet $\mathbf3$, an antiquark an antitriplet $\mathbf{\bar3}$.
Only colour \emph{singlets} exist as free particles, the hadrons~\cite{navas2024}. The products
$\mathbf3\otimes\mathbf{\bar3}=\mathbf1\oplus\mathbf8$ and $\mathbf3\otimes\mathbf3=\mathbf{\bar3}\oplus\mathbf6$ give the ordinary
hadrons: a meson $q\bar q$ in the singlet, and a baryon $qqq$, whose first two quarks couple to $\mathbf{\bar3}$ and then with the third to
a singlet. One-gluon exchange attracts in the singlet $q\bar q$ and, half as strongly, in the $\mathbf{\bar3}$ $qq$, and repels in the
$\mathbf8$ and the $\mathbf6$.

\emph{Diquarks.} A \emph{diquark} is a pair of quarks in the attractive $\mathbf{\bar3}$. The ``good'' diquark, with spin 0 and a flavour-
antisymmetric pair such as $[ud]$, binds most; the ``bad'' one has spin 1~\cite{jaffe2005}.

\emph{Compact or molecular.} A hadron with more than three quarks, an \emph{exotic}, can be \emph{compact}, for example a diquark and an
antidiquark bound by colour forces, $[qq][\bar q\bar q]$~\cite{maiani2005}. Or it can be a \emph{hadronic molecule}: two colour-singlet
hadrons held together by the residual strong force, like the deuteron~\cite{guo2018}. Molecules sit just below the threshold
$m_1+m_2$ of their two constituents.

\emph{The colour-spin interaction.} In the constituent-quark model a hadron's mass is the sum of effective quark masses plus a hyperfine
term proportional to $(\boldsymbol\lambda_i\cdot\boldsymbol\lambda_j)(\boldsymbol\sigma_i\cdot\boldsymbol\sigma_j)/m_im_j$ for each pair,
where the $\boldsymbol\lambda$ are colour and the $\boldsymbol\sigma$ spin matrices. This chromomagnetic interaction of De R\'ujula, Georgi
and Glashow explains, for example, why the $\Delta$ is heavier than the nucleon~\cite{derujula1975}. Karliner and Rosner (KR) refined it
with separate effective quark masses in mesons and in baryons, the baryon masses for every constituent of a compact multiquark, and binding
terms for heavy pairs with $B(QQ)=B(Q\bar Q)/2$; with it they predicted the doubly charmed baryon and doubly heavy
tetraquarks~\cite{karliner2014,karliner2017}.

\emph{Thresholds and cusps.} Where a new two-hadron channel opens, at $M=m_1+m_2$, the scattering amplitude has a branch point. A
\emph{cusp} can then appear in a mass spectrum without any state, and a state that couples to the channel is pulled towards the threshold.
Many exotic candidates lie within a few MeV of a threshold~\cite{guo2018}.

\emph{Names.} The Particle Data Group names hadrons by their quantum numbers, not by their structure~\cite{navas2024}; the
$X(3872)$, for example, is listed as $\chi_{c1}(3872)$. LHCb's convention for exotics~\cite{gershon2022} uses $T$ for four-quark and $P$
for five-quark states, with subscripts for heavy flavour, $\psi$ standing for hidden charm $c\bar c$. So $T_{cc}^+$ is $cc\bar u\bar d$,
seen just below the $D^{*+}D^0$ threshold~\cite{lhcb2022h}, and the hidden-charm pentaquark $P_c(4312)^+$~\cite{lhcb2019h} becomes
$P_\psi^N(4312)^+$.

\emph{String junctions.} In a baryon the three colour flux tubes meet at a point, a \emph{string junction}~\cite{rossiveneziano1977}. A
compact tetraquark has two junctions and a compact pentaquark three, and some models add a constant energy per
junction~\cite{karliner2017a}.
\end{primer}
}
 once after
% \begin{document} defines the macros below, and \primerDMP and so on place the boxes (build_bib.py then numbers the primers'
% references from the position of that \input).
%
%   primer  label         file                      first use
%   P1      box:dmp       primers/p1-dmp.tex        Section sec:dmp           (%PRIMER:dmp)
%   P2      box:qi        primers/p2-qi.tex         Section sec:fix           (%PRIMER:qi, before "The frame is classical")
%   P3      box:gr        primers/p3-gr.tex         Section sec:fields        (%PRIMER:gr, before "The frame")
%   P4      box:cq        primers/p4-cq.tex         Section sec:cq            (%PRIMER:cq)
%   P5      box:holo      primers/p5-holo.tex       Section sec:holo          (%PRIMER:holo)
%   P6      box:cosmo     primers/p6-cosmo.tex      Section sec:cosmo         (%PRIMER:cosmo)
%   P7      box:galaxy    primers/p7-galaxy.tex     Section sec:sparc         (or sec:darkmatter)
%   P8      box:stats     primers/p8-stats.tex      Section sec:scoring
%   P9      box:particle  primers/p9-particle.tex   Section sec:matter        (%PRIMER:particle)
%   P10     box:hadron    primers/p10-hadron.tex    Section sec:trees         (%PRIMER:hadron)
%
% Each box is background only: established physics or DMP's own text, with citations; none states a result of this paper.
% Every key cited is defined in bib-core.tex, bib-extra.tex or sections/*-bib.tex.
\newcommand{\primerDMP}{% GUID v6, Primer P1 (label box:dmp): DMP and trust systems, with a worked act. Background only; it states no result of the paper.
% Place it at the start of Section sec:dmp (marker %PRIMER:dmp in part1.tex): % GUID v6, Primer P1 (label box:dmp): DMP and trust systems, with a worked act. Background only; it states no result of the paper.
% Place it at the start of Section sec:dmp (marker %PRIMER:dmp in part1.tex): % GUID v6, Primer P1 (label box:dmp): DMP and trust systems, with a worked act. Background only; it states no result of the paper.
% Place it at the start of Section sec:dmp (marker %PRIMER:dmp in part1.tex): \input{primers/p1-dmp}
\begin{primer}{DMP and trust systems}{box:dmp}
\emph{What a trust system is for.} In a computer system an access-control or trust protocol decides, for every request, who may do what
to which information, and it keeps the evidence of what was done. Its classic design rules are to grant the least privilege an act needs,
to demand two independent keys for what matters (separation of privilege), and to keep a complete audit trail~\cite{saltzer1975}. Its
formal tools include lattices that say where information may flow~\cite{denning}, noninterference, which says what one side may learn of
the other~\cite{gm}, and quorums of witnesses that fail independently~\cite{mr}. DMP~\cite{dmp} states such a system as three axioms and a
ground, G0; every later statement is a numbered theorem (T) or derived constraint (C).

\emph{Words.} A \emph{part} is ``a Duality of Trust'': a trust half and an information half. An \emph{act} is what a part does. A
\emph{layer} is where trust over acts is decided; its two \emph{witnesses} are its halves. The \emph{trust half} decides whether the act
may happen, the \emph{information half} whether it is meant and right. Two witnesses are \emph{independent} when they share no member of the
\emph{failure classes} the layer declares (host, key, author, engine, person, time source, and whatever the minting lists). A
\emph{standing} is trust held in the field of Axiom~2. It has a source (the witness that granted it), a depth (its \emph{grant chain} back
to the first pair) and a freshness that decays with time; \emph{staleness} is that decay. A \emph{record} is evidence that outlives its act.

\emph{The gate.} Each witness returns a verdict in the lattice $\textsf{refuse}<\textsf{hold}<\textsf{pass}$, and the gate of layer $L$
is their meet, $\operatorname{gate}_L(a,t)=v_{w_T}(a,t)\wedge v_{w_I}(a,t)$. The stricter answer stands, and the act is performed only on
$\textsf{pass}\wedge\textsf{pass}$.

\emph{The axioms in one line each.}
\begin{itemize}[nosep,leftmargin=*]
\item \emph{Axiom~1, the Duality.} A part acts only while both halves of its layer trust it for that act; neither overrules the other. Its
\emph{balance}: every part works towards the trust its boundary declares.
\item \emph{Axiom~2, Trust.} A witness trusts only on evidence it verified itself. ``Trust is gravitational. It is a field, never a
position.'' A strong enough boundary ``bends trustspace'' so that only entropy leaves.
\item \emph{Axiom~3, Information.} Every grant, use and loss of trust leaves a record that every witness of the act reads in full. Any other
reader sees the \emph{counts} always and the content only by a \emph{grant}; from trust, nothing about information can be derived but its
existence. Every layer is a tree rooted at its own Duality. The \emph{layout}, where an item sits in that tree, is ``the median of that
demand'', the place of fewest hops in total, recomputed by both halves and never by one.
\end{itemize}

\emph{A worked act} (Figure~\ref{fig:act}). A part $p$ asks to change one entry of a directory. The trust half checks $p$'s standing: it
was granted by a witness, lies two grants from the minted pair and is still fresh, so the verdict is \textsf{pass}. The information half
finds the change well formed but not yet confirmed: \textsf{hold}. The gate gives $\textsf{pass}\wedge\textsf{hold}=\textsf{hold}$, and
nothing is done. Once the change is confirmed, both return \textsf{pass} and the act is performed. A record is written at the act, and both
witnesses can read it in full. A reader outside the layer, say an auditor, sees the counts (here, one more act of this layer) but not the
new entry, unless it is granted to him. From the trust side alone one learns only that the entry exists.

\emph{Containment and G0.} A boundary \emph{contains} information when no part beyond it is trusted for any act on that information.
G0: ``The dual gate cannot witness its own first instance: somebody has to trust the first two witnesses.'' The first act, the
\emph{minting}, makes at least five parts the holders of the \emph{root keys}; it is stated, not derived, and every later standing
descends from it through grant chains.

\emph{The theorems this paper cites.}
\begin{itemize}[nosep,leftmargin=*]
\item \emph{T4}: every parameter is configuration. DMP fixes structure, never a value, and so derives no constant.
\item \emph{T27}: a contained item is destroyed at its scheduled date; with T31 it describes the trap.
\item \emph{T31}: contained information reaches no part by any act. What still leaves is minimal, thermal emission, the \emph{entropy
floor}, which ``carries entropy and no information''. Outside, one learns only the size, the staleness and the rotation interval.
\item \emph{T32}: the VIBE field. Each member has a depth, a mass (one plus the members deriving through it) and a shortfall; its
potential $\Phi_{\mathrm V}\ge0$ never rises under the system's own acts and vanishes in the \emph{balanced state}. Readers other than the
witnesses see one number, the boundary's temperature, its mass-weighted staleness.
\item \emph{T37}: information that no admission names is contained at the item itself.
\item \emph{T40}: a minting originates only from G0 together with a state in which every axiom and theorem provably holds.
\end{itemize}
\end{primer}

\begin{primer}{DMP and trust systems}{box:dmp}
\emph{What a trust system is for.} In a computer system an access-control or trust protocol decides, for every request, who may do what
to which information, and it keeps the evidence of what was done. Its classic design rules are to grant the least privilege an act needs,
to demand two independent keys for what matters (separation of privilege), and to keep a complete audit trail~\cite{saltzer1975}. Its
formal tools include lattices that say where information may flow~\cite{denning}, noninterference, which says what one side may learn of
the other~\cite{gm}, and quorums of witnesses that fail independently~\cite{mr}. DMP~\cite{dmp} states such a system as three axioms and a
ground, G0; every later statement is a numbered theorem (T) or derived constraint (C).

\emph{Words.} A \emph{part} is ``a Duality of Trust'': a trust half and an information half. An \emph{act} is what a part does. A
\emph{layer} is where trust over acts is decided; its two \emph{witnesses} are its halves. The \emph{trust half} decides whether the act
may happen, the \emph{information half} whether it is meant and right. Two witnesses are \emph{independent} when they share no member of the
\emph{failure classes} the layer declares (host, key, author, engine, person, time source, and whatever the minting lists). A
\emph{standing} is trust held in the field of Axiom~2. It has a source (the witness that granted it), a depth (its \emph{grant chain} back
to the first pair) and a freshness that decays with time; \emph{staleness} is that decay. A \emph{record} is evidence that outlives its act.

\emph{The gate.} Each witness returns a verdict in the lattice $\textsf{refuse}<\textsf{hold}<\textsf{pass}$, and the gate of layer $L$
is their meet, $\operatorname{gate}_L(a,t)=v_{w_T}(a,t)\wedge v_{w_I}(a,t)$. The stricter answer stands, and the act is performed only on
$\textsf{pass}\wedge\textsf{pass}$.

\emph{The axioms in one line each.}
\begin{itemize}[nosep,leftmargin=*]
\item \emph{Axiom~1, the Duality.} A part acts only while both halves of its layer trust it for that act; neither overrules the other. Its
\emph{balance}: every part works towards the trust its boundary declares.
\item \emph{Axiom~2, Trust.} A witness trusts only on evidence it verified itself. ``Trust is gravitational. It is a field, never a
position.'' A strong enough boundary ``bends trustspace'' so that only entropy leaves.
\item \emph{Axiom~3, Information.} Every grant, use and loss of trust leaves a record that every witness of the act reads in full. Any other
reader sees the \emph{counts} always and the content only by a \emph{grant}; from trust, nothing about information can be derived but its
existence. Every layer is a tree rooted at its own Duality. The \emph{layout}, where an item sits in that tree, is ``the median of that
demand'', the place of fewest hops in total, recomputed by both halves and never by one.
\end{itemize}

\emph{A worked act} (Figure~\ref{fig:act}). A part $p$ asks to change one entry of a directory. The trust half checks $p$'s standing: it
was granted by a witness, lies two grants from the minted pair and is still fresh, so the verdict is \textsf{pass}. The information half
finds the change well formed but not yet confirmed: \textsf{hold}. The gate gives $\textsf{pass}\wedge\textsf{hold}=\textsf{hold}$, and
nothing is done. Once the change is confirmed, both return \textsf{pass} and the act is performed. A record is written at the act, and both
witnesses can read it in full. A reader outside the layer, say an auditor, sees the counts (here, one more act of this layer) but not the
new entry, unless it is granted to him. From the trust side alone one learns only that the entry exists.

\emph{Containment and G0.} A boundary \emph{contains} information when no part beyond it is trusted for any act on that information.
G0: ``The dual gate cannot witness its own first instance: somebody has to trust the first two witnesses.'' The first act, the
\emph{minting}, makes at least five parts the holders of the \emph{root keys}; it is stated, not derived, and every later standing
descends from it through grant chains.

\emph{The theorems this paper cites.}
\begin{itemize}[nosep,leftmargin=*]
\item \emph{T4}: every parameter is configuration. DMP fixes structure, never a value, and so derives no constant.
\item \emph{T27}: a contained item is destroyed at its scheduled date; with T31 it describes the trap.
\item \emph{T31}: contained information reaches no part by any act. What still leaves is minimal, thermal emission, the \emph{entropy
floor}, which ``carries entropy and no information''. Outside, one learns only the size, the staleness and the rotation interval.
\item \emph{T32}: the VIBE field. Each member has a depth, a mass (one plus the members deriving through it) and a shortfall; its
potential $\Phi_{\mathrm V}\ge0$ never rises under the system's own acts and vanishes in the \emph{balanced state}. Readers other than the
witnesses see one number, the boundary's temperature, its mass-weighted staleness.
\item \emph{T37}: information that no admission names is contained at the item itself.
\item \emph{T40}: a minting originates only from G0 together with a state in which every axiom and theorem provably holds.
\end{itemize}
\end{primer}

\begin{primer}{DMP and trust systems}{box:dmp}
\emph{What a trust system is for.} In a computer system an access-control or trust protocol decides, for every request, who may do what
to which information, and it keeps the evidence of what was done. Its classic design rules are to grant the least privilege an act needs,
to demand two independent keys for what matters (separation of privilege), and to keep a complete audit trail~\cite{saltzer1975}. Its
formal tools include lattices that say where information may flow~\cite{denning}, noninterference, which says what one side may learn of
the other~\cite{gm}, and quorums of witnesses that fail independently~\cite{mr}. DMP~\cite{dmp} states such a system as three axioms and a
ground, G0; every later statement is a numbered theorem (T) or derived constraint (C).

\emph{Words.} A \emph{part} is ``a Duality of Trust'': a trust half and an information half. An \emph{act} is what a part does. A
\emph{layer} is where trust over acts is decided; its two \emph{witnesses} are its halves. The \emph{trust half} decides whether the act
may happen, the \emph{information half} whether it is meant and right. Two witnesses are \emph{independent} when they share no member of the
\emph{failure classes} the layer declares (host, key, author, engine, person, time source, and whatever the minting lists). A
\emph{standing} is trust held in the field of Axiom~2. It has a source (the witness that granted it), a depth (its \emph{grant chain} back
to the first pair) and a freshness that decays with time; \emph{staleness} is that decay. A \emph{record} is evidence that outlives its act.

\emph{The gate.} Each witness returns a verdict in the lattice $\textsf{refuse}<\textsf{hold}<\textsf{pass}$, and the gate of layer $L$
is their meet, $\operatorname{gate}_L(a,t)=v_{w_T}(a,t)\wedge v_{w_I}(a,t)$. The stricter answer stands, and the act is performed only on
$\textsf{pass}\wedge\textsf{pass}$.

\emph{The axioms in one line each.}
\begin{itemize}[nosep,leftmargin=*]
\item \emph{Axiom~1, the Duality.} A part acts only while both halves of its layer trust it for that act; neither overrules the other. Its
\emph{balance}: every part works towards the trust its boundary declares.
\item \emph{Axiom~2, Trust.} A witness trusts only on evidence it verified itself. ``Trust is gravitational. It is a field, never a
position.'' A strong enough boundary ``bends trustspace'' so that only entropy leaves.
\item \emph{Axiom~3, Information.} Every grant, use and loss of trust leaves a record that every witness of the act reads in full. Any other
reader sees the \emph{counts} always and the content only by a \emph{grant}; from trust, nothing about information can be derived but its
existence. Every layer is a tree rooted at its own Duality. The \emph{layout}, where an item sits in that tree, is ``the median of that
demand'', the place of fewest hops in total, recomputed by both halves and never by one.
\end{itemize}

\emph{A worked act} (Figure~\ref{fig:act}). A part $p$ asks to change one entry of a directory. The trust half checks $p$'s standing: it
was granted by a witness, lies two grants from the minted pair and is still fresh, so the verdict is \textsf{pass}. The information half
finds the change well formed but not yet confirmed: \textsf{hold}. The gate gives $\textsf{pass}\wedge\textsf{hold}=\textsf{hold}$, and
nothing is done. Once the change is confirmed, both return \textsf{pass} and the act is performed. A record is written at the act, and both
witnesses can read it in full. A reader outside the layer, say an auditor, sees the counts (here, one more act of this layer) but not the
new entry, unless it is granted to him. From the trust side alone one learns only that the entry exists.

\emph{Containment and G0.} A boundary \emph{contains} information when no part beyond it is trusted for any act on that information.
G0: ``The dual gate cannot witness its own first instance: somebody has to trust the first two witnesses.'' The first act, the
\emph{minting}, makes at least five parts the holders of the \emph{root keys}; it is stated, not derived, and every later standing
descends from it through grant chains.

\emph{The theorems this paper cites.}
\begin{itemize}[nosep,leftmargin=*]
\item \emph{T4}: every parameter is configuration. DMP fixes structure, never a value, and so derives no constant.
\item \emph{T27}: a contained item is destroyed at its scheduled date; with T31 it describes the trap.
\item \emph{T31}: contained information reaches no part by any act. What still leaves is minimal, thermal emission, the \emph{entropy
floor}, which ``carries entropy and no information''. Outside, one learns only the size, the staleness and the rotation interval.
\item \emph{T32}: the VIBE field. Each member has a depth, a mass (one plus the members deriving through it) and a shortfall; its
potential $\Phi_{\mathrm V}\ge0$ never rises under the system's own acts and vanishes in the \emph{balanced state}. Readers other than the
witnesses see one number, the boundary's temperature, its mass-weighted staleness.
\item \emph{T37}: information that no admission names is contained at the item itself.
\item \emph{T40}: a minting originates only from G0 together with a state in which every axiom and theorem provably holds.
\end{itemize}
\end{primer}
}
\newcommand{\primerQI}{% GUID v6, Primer P2 (label box:qi): quantum information. Background only; it states no result of the paper.
% Place it before "The frame is classical" (marker %PRIMER:qi in part1.tex): % GUID v6, Primer P2 (label box:qi): quantum information. Background only; it states no result of the paper.
% Place it before "The frame is classical" (marker %PRIMER:qi in part1.tex): % GUID v6, Primer P2 (label box:qi): quantum information. Background only; it states no result of the paper.
% Place it before "The frame is classical" (marker %PRIMER:qi in part1.tex): \input{primers/p2-qi}
\begin{primer}{Quantum information}{box:qi}
\emph{States and processes.} A quantum state is a \emph{density matrix} $\rho$: Hermitian, positive ($\rho\succeq0$) and of unit trace.
It is pure when $\rho=\ket\psi\bra\psi$ and mixed otherwise. The state of one part $A$ of a system $AB$ is the partial trace
$\rho_A=\Tr_B\rho_{AB}$~\cite{nielsen2010}. Every physical process is a \emph{completely positive, trace-preserving} (CPTP) map,
$\rho\mapsto\sum_kK_k\rho K_k^\dagger$ with $\sum_kK_k^\dagger K_k=\mathbb 1$; ``completely'' means that the map stays positive when it
acts on one part of a larger system. A measurement with outcomes $k$ returns $k$ with probability $p_k=\Tr(K_k\rho K_k^\dagger)$ and leaves
the state $K_k\rho K_k^\dagger/p_k$.

\emph{Dephasing is reading.} Dephasing in a basis $\{\ket k\}$, $\rho\mapsto\sum_k\ket k\bra k\rho\ket k\bra k$, erases the off-diagonal
elements (the coherences) and keeps the diagonal. It is what a measurement in that basis does to a system when the outcome is written
somewhere else. A system that is dephased in one fixed basis after every interaction behaves as a classical variable.

\emph{Entanglement.} A state of $AB$ is \emph{separable} if it is a mixture of products, $\rho_{AB}=\sum_ip_i\,\rho_A^i\otimes\rho_B^i$,
and \emph{entangled} otherwise~\cite{horodecki2009}. Two numbers measure it for two qubits.
\begin{itemize}[nosep,leftmargin=*]
\item The \emph{negativity} $\mathcal N$ is the magnitude of the sum of the negative eigenvalues of the partial transpose $\rho^{T_B}$
(the transpose taken on $B$ alone)~\cite{vw}. For two qubits $\mathcal N>0$ exactly when the state is entangled~\cite{peres1996,horodecki1996}.
\item The \emph{concurrence} of a pure two-qubit state $a\ket{00}+b\ket{01}+c\ket{10}+d\ket{11}$ is $C=2|ad-bc|$, from 0 for a product to 1
for a maximally entangled state~\cite{wootters1998}. For pure two-qubit states $\mathcal N=C/2$.
\end{itemize}
An \emph{entanglement witness} is an observable $W$ with $\Tr(W\sigma)\ge0$ for every separable $\sigma$; a measured $\Tr(W\rho)<0$
certifies entanglement~\cite{guhne2009}. A given witness misses every entangled state for which $\Tr(W\rho)\ge0$.

\emph{LOCC.} In local operations and classical communication each party acts on its own system and the parties exchange only measurement
outcomes. LOCC maps take separable states to separable states, and no entanglement measure grows under them~\cite{vw,horodecki2009}.
Creating entanglement between two systems therefore needs a carrier that is itself quantum.

\emph{No-cloning and no-broadcasting.} No process copies two pure states that are neither equal nor orthogonal~\cite{wz,dieks}. A family
of states can be \emph{broadcast}, so that each of two readers ends up with the original, only if its members commute pairwise~\cite{barnum}.
A quantity that several readers hold in full is therefore diagonal in one basis, which is the core of quantum Darwinism~\cite{zurek}.

\emph{Worked example: a dephased mediator.} Qubits $A$ and $B$ start in $\ket+=(\ket0+\ket1)/\sqrt2$ and a mediator qubit $M$ in $\ket0$.
They interact only through $M$: a CNOT from $A$ to $M$, a controlled-$Z$ between $M$ and $B$, and the same CNOT again.
\begin{itemize}[nosep,leftmargin=*]
\item \emph{$M$ coherent.} The first CNOT gives $(\ket{0_A0_M}+\ket{1_A1_M})/\sqrt2$. The controlled-$Z$ turns $B$ into
$\ket-=(\ket0-\ket1)/\sqrt2$ in the second branch, and the second CNOT returns $M$ to $\ket0$. What remains is
$(\ket{0}_A\ket{+}_B+\ket{1}_A\ket{-}_B)/\sqrt2$, a maximally entangled state: $C=1$, $\mathcal N=1/2$.
\item \emph{$M$ dephased} in $\{\ket0,\ket1\}$ after the first CNOT. The state becomes the mixture
$\tfrac12(\ket{00}\bra{00}+\ket{11}\bra{11})_{AM}\otimes\ket+\bra+_B$, and the same two steps leave
$\tfrac12(\ket0\bra0_A\otimes\ket+\bra+_B+\ket1\bra1_A\otimes\ket-\bra-_B)$. This is separable: $\mathcal N=0$.
\end{itemize}
The dephased mediator carried one classical bit from $A$ to $B$; that is LOCC, and it correlates the two qubits without entangling them.
\end{primer}

\begin{primer}{Quantum information}{box:qi}
\emph{States and processes.} A quantum state is a \emph{density matrix} $\rho$: Hermitian, positive ($\rho\succeq0$) and of unit trace.
It is pure when $\rho=\ket\psi\bra\psi$ and mixed otherwise. The state of one part $A$ of a system $AB$ is the partial trace
$\rho_A=\Tr_B\rho_{AB}$~\cite{nielsen2010}. Every physical process is a \emph{completely positive, trace-preserving} (CPTP) map,
$\rho\mapsto\sum_kK_k\rho K_k^\dagger$ with $\sum_kK_k^\dagger K_k=\mathbb 1$; ``completely'' means that the map stays positive when it
acts on one part of a larger system. A measurement with outcomes $k$ returns $k$ with probability $p_k=\Tr(K_k\rho K_k^\dagger)$ and leaves
the state $K_k\rho K_k^\dagger/p_k$.

\emph{Dephasing is reading.} Dephasing in a basis $\{\ket k\}$, $\rho\mapsto\sum_k\ket k\bra k\rho\ket k\bra k$, erases the off-diagonal
elements (the coherences) and keeps the diagonal. It is what a measurement in that basis does to a system when the outcome is written
somewhere else. A system that is dephased in one fixed basis after every interaction behaves as a classical variable.

\emph{Entanglement.} A state of $AB$ is \emph{separable} if it is a mixture of products, $\rho_{AB}=\sum_ip_i\,\rho_A^i\otimes\rho_B^i$,
and \emph{entangled} otherwise~\cite{horodecki2009}. Two numbers measure it for two qubits.
\begin{itemize}[nosep,leftmargin=*]
\item The \emph{negativity} $\mathcal N$ is the magnitude of the sum of the negative eigenvalues of the partial transpose $\rho^{T_B}$
(the transpose taken on $B$ alone)~\cite{vw}. For two qubits $\mathcal N>0$ exactly when the state is entangled~\cite{peres1996,horodecki1996}.
\item The \emph{concurrence} of a pure two-qubit state $a\ket{00}+b\ket{01}+c\ket{10}+d\ket{11}$ is $C=2|ad-bc|$, from 0 for a product to 1
for a maximally entangled state~\cite{wootters1998}. For pure two-qubit states $\mathcal N=C/2$.
\end{itemize}
An \emph{entanglement witness} is an observable $W$ with $\Tr(W\sigma)\ge0$ for every separable $\sigma$; a measured $\Tr(W\rho)<0$
certifies entanglement~\cite{guhne2009}. A given witness misses every entangled state for which $\Tr(W\rho)\ge0$.

\emph{LOCC.} In local operations and classical communication each party acts on its own system and the parties exchange only measurement
outcomes. LOCC maps take separable states to separable states, and no entanglement measure grows under them~\cite{vw,horodecki2009}.
Creating entanglement between two systems therefore needs a carrier that is itself quantum.

\emph{No-cloning and no-broadcasting.} No process copies two pure states that are neither equal nor orthogonal~\cite{wz,dieks}. A family
of states can be \emph{broadcast}, so that each of two readers ends up with the original, only if its members commute pairwise~\cite{barnum}.
A quantity that several readers hold in full is therefore diagonal in one basis, which is the core of quantum Darwinism~\cite{zurek}.

\emph{Worked example: a dephased mediator.} Qubits $A$ and $B$ start in $\ket+=(\ket0+\ket1)/\sqrt2$ and a mediator qubit $M$ in $\ket0$.
They interact only through $M$: a CNOT from $A$ to $M$, a controlled-$Z$ between $M$ and $B$, and the same CNOT again.
\begin{itemize}[nosep,leftmargin=*]
\item \emph{$M$ coherent.} The first CNOT gives $(\ket{0_A0_M}+\ket{1_A1_M})/\sqrt2$. The controlled-$Z$ turns $B$ into
$\ket-=(\ket0-\ket1)/\sqrt2$ in the second branch, and the second CNOT returns $M$ to $\ket0$. What remains is
$(\ket{0}_A\ket{+}_B+\ket{1}_A\ket{-}_B)/\sqrt2$, a maximally entangled state: $C=1$, $\mathcal N=1/2$.
\item \emph{$M$ dephased} in $\{\ket0,\ket1\}$ after the first CNOT. The state becomes the mixture
$\tfrac12(\ket{00}\bra{00}+\ket{11}\bra{11})_{AM}\otimes\ket+\bra+_B$, and the same two steps leave
$\tfrac12(\ket0\bra0_A\otimes\ket+\bra+_B+\ket1\bra1_A\otimes\ket-\bra-_B)$. This is separable: $\mathcal N=0$.
\end{itemize}
The dephased mediator carried one classical bit from $A$ to $B$; that is LOCC, and it correlates the two qubits without entangling them.
\end{primer}

\begin{primer}{Quantum information}{box:qi}
\emph{States and processes.} A quantum state is a \emph{density matrix} $\rho$: Hermitian, positive ($\rho\succeq0$) and of unit trace.
It is pure when $\rho=\ket\psi\bra\psi$ and mixed otherwise. The state of one part $A$ of a system $AB$ is the partial trace
$\rho_A=\Tr_B\rho_{AB}$~\cite{nielsen2010}. Every physical process is a \emph{completely positive, trace-preserving} (CPTP) map,
$\rho\mapsto\sum_kK_k\rho K_k^\dagger$ with $\sum_kK_k^\dagger K_k=\mathbb 1$; ``completely'' means that the map stays positive when it
acts on one part of a larger system. A measurement with outcomes $k$ returns $k$ with probability $p_k=\Tr(K_k\rho K_k^\dagger)$ and leaves
the state $K_k\rho K_k^\dagger/p_k$.

\emph{Dephasing is reading.} Dephasing in a basis $\{\ket k\}$, $\rho\mapsto\sum_k\ket k\bra k\rho\ket k\bra k$, erases the off-diagonal
elements (the coherences) and keeps the diagonal. It is what a measurement in that basis does to a system when the outcome is written
somewhere else. A system that is dephased in one fixed basis after every interaction behaves as a classical variable.

\emph{Entanglement.} A state of $AB$ is \emph{separable} if it is a mixture of products, $\rho_{AB}=\sum_ip_i\,\rho_A^i\otimes\rho_B^i$,
and \emph{entangled} otherwise~\cite{horodecki2009}. Two numbers measure it for two qubits.
\begin{itemize}[nosep,leftmargin=*]
\item The \emph{negativity} $\mathcal N$ is the magnitude of the sum of the negative eigenvalues of the partial transpose $\rho^{T_B}$
(the transpose taken on $B$ alone)~\cite{vw}. For two qubits $\mathcal N>0$ exactly when the state is entangled~\cite{peres1996,horodecki1996}.
\item The \emph{concurrence} of a pure two-qubit state $a\ket{00}+b\ket{01}+c\ket{10}+d\ket{11}$ is $C=2|ad-bc|$, from 0 for a product to 1
for a maximally entangled state~\cite{wootters1998}. For pure two-qubit states $\mathcal N=C/2$.
\end{itemize}
An \emph{entanglement witness} is an observable $W$ with $\Tr(W\sigma)\ge0$ for every separable $\sigma$; a measured $\Tr(W\rho)<0$
certifies entanglement~\cite{guhne2009}. A given witness misses every entangled state for which $\Tr(W\rho)\ge0$.

\emph{LOCC.} In local operations and classical communication each party acts on its own system and the parties exchange only measurement
outcomes. LOCC maps take separable states to separable states, and no entanglement measure grows under them~\cite{vw,horodecki2009}.
Creating entanglement between two systems therefore needs a carrier that is itself quantum.

\emph{No-cloning and no-broadcasting.} No process copies two pure states that are neither equal nor orthogonal~\cite{wz,dieks}. A family
of states can be \emph{broadcast}, so that each of two readers ends up with the original, only if its members commute pairwise~\cite{barnum}.
A quantity that several readers hold in full is therefore diagonal in one basis, which is the core of quantum Darwinism~\cite{zurek}.

\emph{Worked example: a dephased mediator.} Qubits $A$ and $B$ start in $\ket+=(\ket0+\ket1)/\sqrt2$ and a mediator qubit $M$ in $\ket0$.
They interact only through $M$: a CNOT from $A$ to $M$, a controlled-$Z$ between $M$ and $B$, and the same CNOT again.
\begin{itemize}[nosep,leftmargin=*]
\item \emph{$M$ coherent.} The first CNOT gives $(\ket{0_A0_M}+\ket{1_A1_M})/\sqrt2$. The controlled-$Z$ turns $B$ into
$\ket-=(\ket0-\ket1)/\sqrt2$ in the second branch, and the second CNOT returns $M$ to $\ket0$. What remains is
$(\ket{0}_A\ket{+}_B+\ket{1}_A\ket{-}_B)/\sqrt2$, a maximally entangled state: $C=1$, $\mathcal N=1/2$.
\item \emph{$M$ dephased} in $\{\ket0,\ket1\}$ after the first CNOT. The state becomes the mixture
$\tfrac12(\ket{00}\bra{00}+\ket{11}\bra{11})_{AM}\otimes\ket+\bra+_B$, and the same two steps leave
$\tfrac12(\ket0\bra0_A\otimes\ket+\bra+_B+\ket1\bra1_A\otimes\ket-\bra-_B)$. This is separable: $\mathcal N=0$.
\end{itemize}
The dephased mediator carried one classical bit from $A$ to $B$; that is LOCC, and it correlates the two qubits without entangling them.
\end{primer}
}
\newcommand{\primerGR}{% GUID v6, Primer P3 (label box:gr): general relativity and the tetrad. Background only; it states no result of the paper.
% Place it before "The frame" (marker %PRIMER:gr in part1.tex): % GUID v6, Primer P3 (label box:gr): general relativity and the tetrad. Background only; it states no result of the paper.
% Place it before "The frame" (marker %PRIMER:gr in part1.tex): % GUID v6, Primer P3 (label box:gr): general relativity and the tetrad. Background only; it states no result of the paper.
% Place it before "The frame" (marker %PRIMER:gr in part1.tex): \input{primers/p3-gr}
\begin{primer}{General relativity and the tetrad}{box:gr}
\setlength{\parskip}{3pt plus 1pt}%
\emph{The metric.} General relativity describes gravity as the geometry of spacetime. A metric $g_{\mu\nu}$ gives the interval
$\dd s^2=g_{\mu\nu}\dd x^\mu\dd x^\nu$, with signature $(-,+,+,+)$, and Einstein's equation
$G_{\mu\nu}+\Lambda g_{\mu\nu}=(8\pi G/c^4)\,T_{\mu\nu}$ ties its curvature to the energy-momentum tensor $T_{\mu\nu}$ of matter and
radiation~\cite{wald1984}.

\emph{The tetrad.} A tetrad, or frame, is a set of four orthonormal one-forms $e^a=e^a{}_\mu\dd x^\mu$ ($a=0,\dots,3$) with
$g_{\mu\nu}=\eta_{ab}e^a{}_\mu e^b{}_\nu$ and $\eta=\operatorname{diag}(-1,1,1,1)$. At each point it can be rotated and boosted,
${e^a\to\Lambda^a{}_b(x)\,e^b}$, without changing $g_{\mu\nu}$: this is the \emph{local Lorentz freedom}. Of the tetrad's 16 components, 6 are
that freedom and 10 are the metric. General relativity can be written with the tetrad as its only field~\cite{aldrovandi}. A single unit
timelike vector field added to the metric is different: it picks out a rest frame at every point. That is Einstein--aether
theory~\cite{aether}, and its preferred-frame effects are tightly bounded, for example by the equal speeds of gravitational waves and
light~\cite{gw170817}.

\emph{Free fall.} Matter couples \emph{minimally} when it feels gravity only through $g_{\mu\nu}$, with partial derivatives replaced by
covariant ones. Then $\nabla^\mu T_{\mu\nu}=0$, and a small body follows a geodesic,
$\ddot x^\mu+\Gamma^\mu{}_{\alpha\beta}\dot x^\alpha\dot x^\beta=0$ (dots: derivatives by proper time), whatever it is made of. That is the
weak equivalence principle. Experiments compare the accelerations $a_1$, $a_2$ of two materials through the E\"otv\"os ratio
${\eta=2|a_1-a_2|/|a_1+a_2|}$~\cite{microscope}; with indices, $\eta$ is the Minkowski metric.

\emph{PPN parameters.} In weak fields, metric theories are compared through the parametrised post-Newtonian (PPN) expansion~\cite{will}.
With $U>0$ the Newtonian potential,
\[
g_{00}=-1+\frac{2U}{c^2}-2\beta\frac{U^2}{c^4},\qquad g_{ij}=\Big(1+2\gamma\frac{U}{c^2}\Big)\delta_{ij}.
\]
$\gamma$ says how much space curvature a unit mass produces; light bending and the Shapiro delay scale as $(1+\gamma)/2$. $\beta$ measures
the nonlinearity of gravity and enters the perihelion advance with $\gamma$. General relativity has $\gamma=\beta=1$.

\emph{The lapse.} Split into space and time, spacetime has a \emph{lapse} $N$: the proper time that a static clock records per unit of
coordinate time. In a weak potential $\Phi=-U$ it is $N=\sqrt{1+2\Phi/c^2}$, outside a mass $M$ it is $N=\sqrt{1-2GM/rc^2}$, and it
vanishes at a horizon.

\emph{Black holes.} A stationary black hole is fixed by three numbers: its mass $M$, angular momentum $J$ and charge
$Q$~\cite{israel,kerr}. The rotating hole (Kerr) has its horizon at $r_+=(GM/c^2)\big(1+\sqrt{1-a_*^2}\big)$, with spin
$a_*=cJ/GM^2\le1$~\cite{kerr}. Its \emph{surface gravity} $\kappa$ is the acceleration, as measured from far away, needed to hold a body
at the horizon; for a hole without spin $\kappa=c^4/4GM$. Quantum fields make the horizon radiate at $k_BT_H=\hbar\kappa/2\pi c$~\cite{hawking75}.
The \emph{area theorem}: with positive energy, the horizon area of classical general relativity never decreases~\cite{hawking1971}. It is
the second law of black-hole mechanics~\cite{bch}, with entropy $S=k_Bc^3A/4G\hbar$~\cite{bekenstein73}.
\end{primer}

\begin{primer}{General relativity and the tetrad}{box:gr}
\setlength{\parskip}{3pt plus 1pt}%
\emph{The metric.} General relativity describes gravity as the geometry of spacetime. A metric $g_{\mu\nu}$ gives the interval
$\dd s^2=g_{\mu\nu}\dd x^\mu\dd x^\nu$, with signature $(-,+,+,+)$, and Einstein's equation
$G_{\mu\nu}+\Lambda g_{\mu\nu}=(8\pi G/c^4)\,T_{\mu\nu}$ ties its curvature to the energy-momentum tensor $T_{\mu\nu}$ of matter and
radiation~\cite{wald1984}.

\emph{The tetrad.} A tetrad, or frame, is a set of four orthonormal one-forms $e^a=e^a{}_\mu\dd x^\mu$ ($a=0,\dots,3$) with
$g_{\mu\nu}=\eta_{ab}e^a{}_\mu e^b{}_\nu$ and $\eta=\operatorname{diag}(-1,1,1,1)$. At each point it can be rotated and boosted,
${e^a\to\Lambda^a{}_b(x)\,e^b}$, without changing $g_{\mu\nu}$: this is the \emph{local Lorentz freedom}. Of the tetrad's 16 components, 6 are
that freedom and 10 are the metric. General relativity can be written with the tetrad as its only field~\cite{aldrovandi}. A single unit
timelike vector field added to the metric is different: it picks out a rest frame at every point. That is Einstein--aether
theory~\cite{aether}, and its preferred-frame effects are tightly bounded, for example by the equal speeds of gravitational waves and
light~\cite{gw170817}.

\emph{Free fall.} Matter couples \emph{minimally} when it feels gravity only through $g_{\mu\nu}$, with partial derivatives replaced by
covariant ones. Then $\nabla^\mu T_{\mu\nu}=0$, and a small body follows a geodesic,
$\ddot x^\mu+\Gamma^\mu{}_{\alpha\beta}\dot x^\alpha\dot x^\beta=0$ (dots: derivatives by proper time), whatever it is made of. That is the
weak equivalence principle. Experiments compare the accelerations $a_1$, $a_2$ of two materials through the E\"otv\"os ratio
${\eta=2|a_1-a_2|/|a_1+a_2|}$~\cite{microscope}; with indices, $\eta$ is the Minkowski metric.

\emph{PPN parameters.} In weak fields, metric theories are compared through the parametrised post-Newtonian (PPN) expansion~\cite{will}.
With $U>0$ the Newtonian potential,
\[
g_{00}=-1+\frac{2U}{c^2}-2\beta\frac{U^2}{c^4},\qquad g_{ij}=\Big(1+2\gamma\frac{U}{c^2}\Big)\delta_{ij}.
\]
$\gamma$ says how much space curvature a unit mass produces; light bending and the Shapiro delay scale as $(1+\gamma)/2$. $\beta$ measures
the nonlinearity of gravity and enters the perihelion advance with $\gamma$. General relativity has $\gamma=\beta=1$.

\emph{The lapse.} Split into space and time, spacetime has a \emph{lapse} $N$: the proper time that a static clock records per unit of
coordinate time. In a weak potential $\Phi=-U$ it is $N=\sqrt{1+2\Phi/c^2}$, outside a mass $M$ it is $N=\sqrt{1-2GM/rc^2}$, and it
vanishes at a horizon.

\emph{Black holes.} A stationary black hole is fixed by three numbers: its mass $M$, angular momentum $J$ and charge
$Q$~\cite{israel,kerr}. The rotating hole (Kerr) has its horizon at $r_+=(GM/c^2)\big(1+\sqrt{1-a_*^2}\big)$, with spin
$a_*=cJ/GM^2\le1$~\cite{kerr}. Its \emph{surface gravity} $\kappa$ is the acceleration, as measured from far away, needed to hold a body
at the horizon; for a hole without spin $\kappa=c^4/4GM$. Quantum fields make the horizon radiate at $k_BT_H=\hbar\kappa/2\pi c$~\cite{hawking75}.
The \emph{area theorem}: with positive energy, the horizon area of classical general relativity never decreases~\cite{hawking1971}. It is
the second law of black-hole mechanics~\cite{bch}, with entropy $S=k_Bc^3A/4G\hbar$~\cite{bekenstein73}.
\end{primer}

\begin{primer}{General relativity and the tetrad}{box:gr}
\setlength{\parskip}{3pt plus 1pt}%
\emph{The metric.} General relativity describes gravity as the geometry of spacetime. A metric $g_{\mu\nu}$ gives the interval
$\dd s^2=g_{\mu\nu}\dd x^\mu\dd x^\nu$, with signature $(-,+,+,+)$, and Einstein's equation
$G_{\mu\nu}+\Lambda g_{\mu\nu}=(8\pi G/c^4)\,T_{\mu\nu}$ ties its curvature to the energy-momentum tensor $T_{\mu\nu}$ of matter and
radiation~\cite{wald1984}.

\emph{The tetrad.} A tetrad, or frame, is a set of four orthonormal one-forms $e^a=e^a{}_\mu\dd x^\mu$ ($a=0,\dots,3$) with
$g_{\mu\nu}=\eta_{ab}e^a{}_\mu e^b{}_\nu$ and $\eta=\operatorname{diag}(-1,1,1,1)$. At each point it can be rotated and boosted,
${e^a\to\Lambda^a{}_b(x)\,e^b}$, without changing $g_{\mu\nu}$: this is the \emph{local Lorentz freedom}. Of the tetrad's 16 components, 6 are
that freedom and 10 are the metric. General relativity can be written with the tetrad as its only field~\cite{aldrovandi}. A single unit
timelike vector field added to the metric is different: it picks out a rest frame at every point. That is Einstein--aether
theory~\cite{aether}, and its preferred-frame effects are tightly bounded, for example by the equal speeds of gravitational waves and
light~\cite{gw170817}.

\emph{Free fall.} Matter couples \emph{minimally} when it feels gravity only through $g_{\mu\nu}$, with partial derivatives replaced by
covariant ones. Then $\nabla^\mu T_{\mu\nu}=0$, and a small body follows a geodesic,
$\ddot x^\mu+\Gamma^\mu{}_{\alpha\beta}\dot x^\alpha\dot x^\beta=0$ (dots: derivatives by proper time), whatever it is made of. That is the
weak equivalence principle. Experiments compare the accelerations $a_1$, $a_2$ of two materials through the E\"otv\"os ratio
${\eta=2|a_1-a_2|/|a_1+a_2|}$~\cite{microscope}; with indices, $\eta$ is the Minkowski metric.

\emph{PPN parameters.} In weak fields, metric theories are compared through the parametrised post-Newtonian (PPN) expansion~\cite{will}.
With $U>0$ the Newtonian potential,
\[
g_{00}=-1+\frac{2U}{c^2}-2\beta\frac{U^2}{c^4},\qquad g_{ij}=\Big(1+2\gamma\frac{U}{c^2}\Big)\delta_{ij}.
\]
$\gamma$ says how much space curvature a unit mass produces; light bending and the Shapiro delay scale as $(1+\gamma)/2$. $\beta$ measures
the nonlinearity of gravity and enters the perihelion advance with $\gamma$. General relativity has $\gamma=\beta=1$.

\emph{The lapse.} Split into space and time, spacetime has a \emph{lapse} $N$: the proper time that a static clock records per unit of
coordinate time. In a weak potential $\Phi=-U$ it is $N=\sqrt{1+2\Phi/c^2}$, outside a mass $M$ it is $N=\sqrt{1-2GM/rc^2}$, and it
vanishes at a horizon.

\emph{Black holes.} A stationary black hole is fixed by three numbers: its mass $M$, angular momentum $J$ and charge
$Q$~\cite{israel,kerr}. The rotating hole (Kerr) has its horizon at $r_+=(GM/c^2)\big(1+\sqrt{1-a_*^2}\big)$, with spin
$a_*=cJ/GM^2\le1$~\cite{kerr}. Its \emph{surface gravity} $\kappa$ is the acceleration, as measured from far away, needed to hold a body
at the horizon; for a hole without spin $\kappa=c^4/4GM$. Quantum fields make the horizon radiate at $k_BT_H=\hbar\kappa/2\pi c$~\cite{hawking75}.
The \emph{area theorem}: with positive energy, the horizon area of classical general relativity never decreases~\cite{hawking1971}. It is
the second law of black-hole mechanics~\cite{bch}, with entropy $S=k_Bc^3A/4G\hbar$~\cite{bekenstein73}.
\end{primer}
}
\newcommand{\primerCQ}{% GUID v6, Primer P4 (label box:cq): classical-quantum dynamics and collapse. Background only; it states no result of the paper.
% Place it before "Consistent classical-quantum dynamics" (marker %PRIMER:cq in part1.tex): % GUID v6, Primer P4 (label box:cq): classical-quantum dynamics and collapse. Background only; it states no result of the paper.
% Place it before "Consistent classical-quantum dynamics" (marker %PRIMER:cq in part1.tex): % GUID v6, Primer P4 (label box:cq): classical-quantum dynamics and collapse. Background only; it states no result of the paper.
% Place it before "Consistent classical-quantum dynamics" (marker %PRIMER:cq in part1.tex): \input{primers/p4-cq}
\begin{primer}{Classical-quantum dynamics and collapse}{box:cq}
\setlength{\parskip}{3pt plus 1pt}%
\emph{The classical-quantum (CQ) master equation.} A consistent coupling keeps the joint description linear, positive and
normalised, as probability theory demands. The joint state is a family $\varrho(z)$ of positive operators, one for each value
$z$ of the classical variables (here the frame), normalised as $\int\dd z\,\Tr\varrho(z)=1$~\cite{oppenheim}. For one classical variable
$z$ driven by one quantum observable $L$ (for gravity, the mass density), continuous Markovian dynamics has the form~\cite{osswd}
\[
\partial_t\varrho=-\frac i\hbar[H,\varrho]+D_0\Big(L\varrho L-\tfrac12\{L^2,\varrho\}\Big)-\frac{D_1}2\,\partial_z\{L,\varrho\}
+\frac{D_2}2\,\partial_z^2\varrho ,
\]
with three coefficients beyond quantum mechanics:
\begin{itemize}[nosep,leftmargin=*]
\item $D_0$, \emph{decoherence}: the quantum system loses coherence between states that source the classical field differently (for
gravity, different mass distributions);
\item $D_1$, \emph{back-reaction}: a drift of $z$ driven by the quantum state, that is, how matter pushes the frame;
\item $D_2$, \emph{diffusion}: random kicks of $z$, that is, noise in the frame.
\end{itemize}
Complete positivity requires $D_1^2\le4D_0D_2$ here, and $D_0>0$ wherever $D_1\ne0$; Oppenheim et al.\ give the general matrix form of
this trade-off~\cite[Eq.~(23)]{osswd}, which is Eq.~\eqref{eq:tradeoff} of this paper. In words: \emph{a push implies noise}. If quantum matter pushes a
classical field at all, the matter decoheres and the field diffuses, and the product of the two is bounded below by the square of the push,
$D_0D_2\ge D_1^2/4$.

\emph{Measurement and feedback.} An older family of consistent models makes gravity a channel. Each mass's position, or its mass density,
is weakly and continuously measured, and a classical force computed from the record acts back on the other masses (Kafri, Taylor and
Milburn, KTM~\cite{ktm}; Tilloy and Di\'osi~\cite{tilloydiosi}). Newton's law then holds on average, the measurement decoheres
superpositions, and the feedback adds noise. The masses talk only through a classical record, so the channel is LOCC
(Primer~\ref{box:qi}) and cannot entangle them. A weaker measurement decoheres less but gives a noisier record and so more heating; KTM find
both effects at least of the order of the gravitational coupling. The least decoherence of two masses over all measurement strengths is
what this paper calls the \emph{KTM minimum}, $\omega/2$, where $\omega$ is the rate at which a coherent field would build up the
entangling phase.

\emph{Gravity-related collapse.} Di\'osi and Penrose proposed that a superposition of two mass distributions $\rho_1$, $\rho_2$ decays at
the rate $\lambda=E_\Delta/\hbar$, where
\[
E_\Delta=\frac G2\int\!\!\int\dd^3x\,\dd^3y\,\frac{[\rho_1(\mathbf x)-\rho_2(\mathbf x)][\rho_1(\mathbf y)-\rho_2(\mathbf y)]}
{|\mathbf x-\mathbf y|}
\]
is the gravitational self-energy of their difference~\cite{diosi87,penrose96,bassi2013}. For a body of mass $M$ and size $R$ displaced by
more than its size, $\lambda\approx GM^2/\hbar R$: the \emph{Di\'osi--Penrose (DP) rate}. For point-like nuclei $E_\Delta$ diverges, so the
mass density must be smeared over a length $R_0$. The same noise makes charged matter radiate. An underground germanium detector sees no
such radiation, which excludes smearing at nuclear size and requires $R_0$ of order $10^{-10}$\,m or more~\cite{donadi,figurato}. In the
Newtonian limit of CQ gravity the DP kernel is one admissible choice of the decoherence term~\cite{osswd}.
\end{primer}

\begin{primer}{Classical-quantum dynamics and collapse}{box:cq}
\setlength{\parskip}{3pt plus 1pt}%
\emph{The classical-quantum (CQ) master equation.} A consistent coupling keeps the joint description linear, positive and
normalised, as probability theory demands. The joint state is a family $\varrho(z)$ of positive operators, one for each value
$z$ of the classical variables (here the frame), normalised as $\int\dd z\,\Tr\varrho(z)=1$~\cite{oppenheim}. For one classical variable
$z$ driven by one quantum observable $L$ (for gravity, the mass density), continuous Markovian dynamics has the form~\cite{osswd}
\[
\partial_t\varrho=-\frac i\hbar[H,\varrho]+D_0\Big(L\varrho L-\tfrac12\{L^2,\varrho\}\Big)-\frac{D_1}2\,\partial_z\{L,\varrho\}
+\frac{D_2}2\,\partial_z^2\varrho ,
\]
with three coefficients beyond quantum mechanics:
\begin{itemize}[nosep,leftmargin=*]
\item $D_0$, \emph{decoherence}: the quantum system loses coherence between states that source the classical field differently (for
gravity, different mass distributions);
\item $D_1$, \emph{back-reaction}: a drift of $z$ driven by the quantum state, that is, how matter pushes the frame;
\item $D_2$, \emph{diffusion}: random kicks of $z$, that is, noise in the frame.
\end{itemize}
Complete positivity requires $D_1^2\le4D_0D_2$ here, and $D_0>0$ wherever $D_1\ne0$; Oppenheim et al.\ give the general matrix form of
this trade-off~\cite[Eq.~(23)]{osswd}, which is Eq.~\eqref{eq:tradeoff} of this paper. In words: \emph{a push implies noise}. If quantum matter pushes a
classical field at all, the matter decoheres and the field diffuses, and the product of the two is bounded below by the square of the push,
$D_0D_2\ge D_1^2/4$.

\emph{Measurement and feedback.} An older family of consistent models makes gravity a channel. Each mass's position, or its mass density,
is weakly and continuously measured, and a classical force computed from the record acts back on the other masses (Kafri, Taylor and
Milburn, KTM~\cite{ktm}; Tilloy and Di\'osi~\cite{tilloydiosi}). Newton's law then holds on average, the measurement decoheres
superpositions, and the feedback adds noise. The masses talk only through a classical record, so the channel is LOCC
(Primer~\ref{box:qi}) and cannot entangle them. A weaker measurement decoheres less but gives a noisier record and so more heating; KTM find
both effects at least of the order of the gravitational coupling. The least decoherence of two masses over all measurement strengths is
what this paper calls the \emph{KTM minimum}, $\omega/2$, where $\omega$ is the rate at which a coherent field would build up the
entangling phase.

\emph{Gravity-related collapse.} Di\'osi and Penrose proposed that a superposition of two mass distributions $\rho_1$, $\rho_2$ decays at
the rate $\lambda=E_\Delta/\hbar$, where
\[
E_\Delta=\frac G2\int\!\!\int\dd^3x\,\dd^3y\,\frac{[\rho_1(\mathbf x)-\rho_2(\mathbf x)][\rho_1(\mathbf y)-\rho_2(\mathbf y)]}
{|\mathbf x-\mathbf y|}
\]
is the gravitational self-energy of their difference~\cite{diosi87,penrose96,bassi2013}. For a body of mass $M$ and size $R$ displaced by
more than its size, $\lambda\approx GM^2/\hbar R$: the \emph{Di\'osi--Penrose (DP) rate}. For point-like nuclei $E_\Delta$ diverges, so the
mass density must be smeared over a length $R_0$. The same noise makes charged matter radiate. An underground germanium detector sees no
such radiation, which excludes smearing at nuclear size and requires $R_0$ of order $10^{-10}$\,m or more~\cite{donadi,figurato}. In the
Newtonian limit of CQ gravity the DP kernel is one admissible choice of the decoherence term~\cite{osswd}.
\end{primer}

\begin{primer}{Classical-quantum dynamics and collapse}{box:cq}
\setlength{\parskip}{3pt plus 1pt}%
\emph{The classical-quantum (CQ) master equation.} A consistent coupling keeps the joint description linear, positive and
normalised, as probability theory demands. The joint state is a family $\varrho(z)$ of positive operators, one for each value
$z$ of the classical variables (here the frame), normalised as $\int\dd z\,\Tr\varrho(z)=1$~\cite{oppenheim}. For one classical variable
$z$ driven by one quantum observable $L$ (for gravity, the mass density), continuous Markovian dynamics has the form~\cite{osswd}
\[
\partial_t\varrho=-\frac i\hbar[H,\varrho]+D_0\Big(L\varrho L-\tfrac12\{L^2,\varrho\}\Big)-\frac{D_1}2\,\partial_z\{L,\varrho\}
+\frac{D_2}2\,\partial_z^2\varrho ,
\]
with three coefficients beyond quantum mechanics:
\begin{itemize}[nosep,leftmargin=*]
\item $D_0$, \emph{decoherence}: the quantum system loses coherence between states that source the classical field differently (for
gravity, different mass distributions);
\item $D_1$, \emph{back-reaction}: a drift of $z$ driven by the quantum state, that is, how matter pushes the frame;
\item $D_2$, \emph{diffusion}: random kicks of $z$, that is, noise in the frame.
\end{itemize}
Complete positivity requires $D_1^2\le4D_0D_2$ here, and $D_0>0$ wherever $D_1\ne0$; Oppenheim et al.\ give the general matrix form of
this trade-off~\cite[Eq.~(23)]{osswd}, which is Eq.~\eqref{eq:tradeoff} of this paper. In words: \emph{a push implies noise}. If quantum matter pushes a
classical field at all, the matter decoheres and the field diffuses, and the product of the two is bounded below by the square of the push,
$D_0D_2\ge D_1^2/4$.

\emph{Measurement and feedback.} An older family of consistent models makes gravity a channel. Each mass's position, or its mass density,
is weakly and continuously measured, and a classical force computed from the record acts back on the other masses (Kafri, Taylor and
Milburn, KTM~\cite{ktm}; Tilloy and Di\'osi~\cite{tilloydiosi}). Newton's law then holds on average, the measurement decoheres
superpositions, and the feedback adds noise. The masses talk only through a classical record, so the channel is LOCC
(Primer~\ref{box:qi}) and cannot entangle them. A weaker measurement decoheres less but gives a noisier record and so more heating; KTM find
both effects at least of the order of the gravitational coupling. The least decoherence of two masses over all measurement strengths is
what this paper calls the \emph{KTM minimum}, $\omega/2$, where $\omega$ is the rate at which a coherent field would build up the
entangling phase.

\emph{Gravity-related collapse.} Di\'osi and Penrose proposed that a superposition of two mass distributions $\rho_1$, $\rho_2$ decays at
the rate $\lambda=E_\Delta/\hbar$, where
\[
E_\Delta=\frac G2\int\!\!\int\dd^3x\,\dd^3y\,\frac{[\rho_1(\mathbf x)-\rho_2(\mathbf x)][\rho_1(\mathbf y)-\rho_2(\mathbf y)]}
{|\mathbf x-\mathbf y|}
\]
is the gravitational self-energy of their difference~\cite{diosi87,penrose96,bassi2013}. For a body of mass $M$ and size $R$ displaced by
more than its size, $\lambda\approx GM^2/\hbar R$: the \emph{Di\'osi--Penrose (DP) rate}. For point-like nuclei $E_\Delta$ diverges, so the
mass density must be smeared over a length $R_0$. The same noise makes charged matter radiate. An underground germanium detector sees no
such radiation, which excludes smearing at nuclear size and requires $R_0$ of order $10^{-10}$\,m or more~\cite{donadi,figurato}. In the
Newtonian limit of CQ gravity the DP kernel is one admissible choice of the decoherence term~\cite{osswd}.
\end{primer}
}
\newcommand{\primerHolo}{% GUID v6, Primer P5 (label box:holo): holography and quantum error correction. Background only; it states no result of the paper.
% Place it at the start of Section sec:holo (marker %PRIMER:holo in part1.tex): % GUID v6, Primer P5 (label box:holo): holography and quantum error correction. Background only; it states no result of the paper.
% Place it at the start of Section sec:holo (marker %PRIMER:holo in part1.tex): % GUID v6, Primer P5 (label box:holo): holography and quantum error correction. Background only; it states no result of the paper.
% Place it at the start of Section sec:holo (marker %PRIMER:holo in part1.tex): \input{primers/p5-holo}
\begin{primer}{Holography and quantum error correction}{box:holo}
\emph{Stabiliser codes.} A quantum error-correcting code stores $k$ logical qubits in $n$ physical ones so that some errors can be
undone. In a \emph{stabiliser code} the code space is the joint $+1$ eigenspace of $n-k$ commuting products of Pauli matrices, the
\emph{stabilisers}~\cite{gottesman1997}. Measuring them gives the \emph{syndrome}. An error that anticommutes with a stabiliser flips its
sign and so shows in the syndrome. The stabilisers commute with the logical operators, so the syndrome reveals nothing about the stored
content. A code of distance $d$, written $[[n,k,d]]$, corrects any $\lfloor(d-1)/2\rfloor$ unknown errors, or any $d-1$
\emph{erasures} (losses at known places).

\emph{The five-qubit code.} The smallest code that corrects every single-qubit error is $[[5,1,3]]$, with the stabiliser $XZZXI$ and its
cyclic shifts~\cite{laflamme,bennett1996}. It corrects any two erasures, so its logical qubit can be recovered from any three of the five
qubits and from no two: a $((3,5))$ quantum threshold scheme~\cite{cgl}.

\emph{Perfect tensors and the HaPPY code.} A tensor with $2m$ legs is \emph{perfect} if every split of its legs into $m$ and $m$ is, up to
normalisation, a unitary map. The encoding map of the five-qubit code, with its logical leg as a sixth, is such a tensor. Pastawski,
Yoshida, Harlow and Preskill tile the hyperbolic disc with pentagons, place one perfect tensor on each, contract neighbouring legs, and
leave one leg per tile in the bulk and the outermost legs on the boundary~\cite{happy}. The network is an isometry from the bulk to the
boundary, a toy model of the bulk-boundary correspondence of AdS/CFT~\cite{maldacena}.

\emph{Ryu--Takayanagi as a minimal cut.} In AdS/CFT the entropy of a boundary region $A$ is $S(A)=\mathrm{Area}(\gamma_A)/4G\hbar$ (with
$c=k_B=1$), where $\gamma_A$ is the bulk surface of least area anchored on the edge of $A$~\cite{rt}. In a network of perfect tensors the
entropy of $A$, in bits, is for many regions the number of qubit legs cut by the smallest cut that separates $A$ from the rest of the
boundary~\cite{happy}.

\emph{Reconstruction.} The bulk region between $A$ and $\gamma_A$ is the \emph{entanglement wedge} of $A$. Every bulk operator inside it can
be represented by an operator acting on $A$ alone~\cite{adh,dong2016}: the bulk is encoded redundantly, like a logical qubit. A code has
\emph{complementary recovery} when what cannot be reconstructed on $A$ can be reconstructed on the rest of the boundary.

\emph{The area operator.} Harlow showed that every code with complementary recovery obeys a quantum Ryu--Takayanagi formula
(Eq.~\eqref{eq:harlow} in Section~\ref{sec:holo}). Its area term is the expectation value of an \emph{area operator} that lies in the centre
of the algebra reconstructable on $A$: it commutes with everything that either side can do~\cite{harlow}. A \emph{fixed-area state} is an
eigenstate of this operator, and such states have a flat entanglement spectrum~\cite{dhm}. A general state of quantum gravity is a
superposition of areas.
\end{primer}

\begin{primer}{Holography and quantum error correction}{box:holo}
\emph{Stabiliser codes.} A quantum error-correcting code stores $k$ logical qubits in $n$ physical ones so that some errors can be
undone. In a \emph{stabiliser code} the code space is the joint $+1$ eigenspace of $n-k$ commuting products of Pauli matrices, the
\emph{stabilisers}~\cite{gottesman1997}. Measuring them gives the \emph{syndrome}. An error that anticommutes with a stabiliser flips its
sign and so shows in the syndrome. The stabilisers commute with the logical operators, so the syndrome reveals nothing about the stored
content. A code of distance $d$, written $[[n,k,d]]$, corrects any $\lfloor(d-1)/2\rfloor$ unknown errors, or any $d-1$
\emph{erasures} (losses at known places).

\emph{The five-qubit code.} The smallest code that corrects every single-qubit error is $[[5,1,3]]$, with the stabiliser $XZZXI$ and its
cyclic shifts~\cite{laflamme,bennett1996}. It corrects any two erasures, so its logical qubit can be recovered from any three of the five
qubits and from no two: a $((3,5))$ quantum threshold scheme~\cite{cgl}.

\emph{Perfect tensors and the HaPPY code.} A tensor with $2m$ legs is \emph{perfect} if every split of its legs into $m$ and $m$ is, up to
normalisation, a unitary map. The encoding map of the five-qubit code, with its logical leg as a sixth, is such a tensor. Pastawski,
Yoshida, Harlow and Preskill tile the hyperbolic disc with pentagons, place one perfect tensor on each, contract neighbouring legs, and
leave one leg per tile in the bulk and the outermost legs on the boundary~\cite{happy}. The network is an isometry from the bulk to the
boundary, a toy model of the bulk-boundary correspondence of AdS/CFT~\cite{maldacena}.

\emph{Ryu--Takayanagi as a minimal cut.} In AdS/CFT the entropy of a boundary region $A$ is $S(A)=\mathrm{Area}(\gamma_A)/4G\hbar$ (with
$c=k_B=1$), where $\gamma_A$ is the bulk surface of least area anchored on the edge of $A$~\cite{rt}. In a network of perfect tensors the
entropy of $A$, in bits, is for many regions the number of qubit legs cut by the smallest cut that separates $A$ from the rest of the
boundary~\cite{happy}.

\emph{Reconstruction.} The bulk region between $A$ and $\gamma_A$ is the \emph{entanglement wedge} of $A$. Every bulk operator inside it can
be represented by an operator acting on $A$ alone~\cite{adh,dong2016}: the bulk is encoded redundantly, like a logical qubit. A code has
\emph{complementary recovery} when what cannot be reconstructed on $A$ can be reconstructed on the rest of the boundary.

\emph{The area operator.} Harlow showed that every code with complementary recovery obeys a quantum Ryu--Takayanagi formula
(Eq.~\eqref{eq:harlow} in Section~\ref{sec:holo}). Its area term is the expectation value of an \emph{area operator} that lies in the centre
of the algebra reconstructable on $A$: it commutes with everything that either side can do~\cite{harlow}. A \emph{fixed-area state} is an
eigenstate of this operator, and such states have a flat entanglement spectrum~\cite{dhm}. A general state of quantum gravity is a
superposition of areas.
\end{primer}

\begin{primer}{Holography and quantum error correction}{box:holo}
\emph{Stabiliser codes.} A quantum error-correcting code stores $k$ logical qubits in $n$ physical ones so that some errors can be
undone. In a \emph{stabiliser code} the code space is the joint $+1$ eigenspace of $n-k$ commuting products of Pauli matrices, the
\emph{stabilisers}~\cite{gottesman1997}. Measuring them gives the \emph{syndrome}. An error that anticommutes with a stabiliser flips its
sign and so shows in the syndrome. The stabilisers commute with the logical operators, so the syndrome reveals nothing about the stored
content. A code of distance $d$, written $[[n,k,d]]$, corrects any $\lfloor(d-1)/2\rfloor$ unknown errors, or any $d-1$
\emph{erasures} (losses at known places).

\emph{The five-qubit code.} The smallest code that corrects every single-qubit error is $[[5,1,3]]$, with the stabiliser $XZZXI$ and its
cyclic shifts~\cite{laflamme,bennett1996}. It corrects any two erasures, so its logical qubit can be recovered from any three of the five
qubits and from no two: a $((3,5))$ quantum threshold scheme~\cite{cgl}.

\emph{Perfect tensors and the HaPPY code.} A tensor with $2m$ legs is \emph{perfect} if every split of its legs into $m$ and $m$ is, up to
normalisation, a unitary map. The encoding map of the five-qubit code, with its logical leg as a sixth, is such a tensor. Pastawski,
Yoshida, Harlow and Preskill tile the hyperbolic disc with pentagons, place one perfect tensor on each, contract neighbouring legs, and
leave one leg per tile in the bulk and the outermost legs on the boundary~\cite{happy}. The network is an isometry from the bulk to the
boundary, a toy model of the bulk-boundary correspondence of AdS/CFT~\cite{maldacena}.

\emph{Ryu--Takayanagi as a minimal cut.} In AdS/CFT the entropy of a boundary region $A$ is $S(A)=\mathrm{Area}(\gamma_A)/4G\hbar$ (with
$c=k_B=1$), where $\gamma_A$ is the bulk surface of least area anchored on the edge of $A$~\cite{rt}. In a network of perfect tensors the
entropy of $A$, in bits, is for many regions the number of qubit legs cut by the smallest cut that separates $A$ from the rest of the
boundary~\cite{happy}.

\emph{Reconstruction.} The bulk region between $A$ and $\gamma_A$ is the \emph{entanglement wedge} of $A$. Every bulk operator inside it can
be represented by an operator acting on $A$ alone~\cite{adh,dong2016}: the bulk is encoded redundantly, like a logical qubit. A code has
\emph{complementary recovery} when what cannot be reconstructed on $A$ can be reconstructed on the rest of the boundary.

\emph{The area operator.} Harlow showed that every code with complementary recovery obeys a quantum Ryu--Takayanagi formula
(Eq.~\eqref{eq:harlow} in Section~\ref{sec:holo}). Its area term is the expectation value of an \emph{area operator} that lies in the centre
of the algebra reconstructable on $A$: it commutes with everything that either side can do~\cite{harlow}. A \emph{fixed-area state} is an
eigenstate of this operator, and such states have a flat entanglement spectrum~\cite{dhm}. A general state of quantum gravity is a
superposition of areas.
\end{primer}
}
\newcommand{\primerCosmo}{% GUID v6, Primer P6 (label box:cosmo): cosmology. Background only; it states no result of the paper.
% Place it at the start of Section sec:cosmo (marker %PRIMER:cosmo in part1.tex): % GUID v6, Primer P6 (label box:cosmo): cosmology. Background only; it states no result of the paper.
% Place it at the start of Section sec:cosmo (marker %PRIMER:cosmo in part1.tex): % GUID v6, Primer P6 (label box:cosmo): cosmology. Background only; it states no result of the paper.
% Place it at the start of Section sec:cosmo (marker %PRIMER:cosmo in part1.tex): \input{primers/p6-cosmo}
\begin{primer}{Cosmology}{box:cosmo}
\emph{The expanding universe.} On large scales the universe is homogeneous and isotropic. Distances grow with a scale factor $a(t)$, light
from the past is redshifted by $1+z=a_0/a$, and the expansion rate $H=\dot a/a$ obeys, for a flat universe,
$H^2=8\pi G\rho/3+\Lambda c^2/3$. Each component's share is $\Omega_i=\rho_i/\rho_c$ with the critical density $\rho_c=3H^2/8\pi G$; the
value today is $H_0$.

\emph{$\Lambda$CDM.} The standard model of cosmology adds to the known particles a cold, collisionless \emph{dark matter} and a
cosmological constant $\Lambda$. Six parameters fix it: the densities of baryons and dark matter ($\Omega_bh^2$, $\Omega_ch^2$), the angular
size of the sound horizon, the optical depth to reionisation, and the amplitude $A_s$ and tilt $n_s$ of the primordial
fluctuations~\cite{planck18}.

\emph{Inflation.} An early epoch of nearly exponential expansion, $a\propto e^{Ht}$, flattens and smooths the universe and stretches
quantum fluctuations into the seeds of galaxies~\cite{guth}. It is driven by a scalar field, the \emph{inflaton}, with potential $V$, and
its length is counted in \emph{e-folds}, $N=\ln(a_{\rm end}/a_{\rm start})$. In \emph{slow roll} the field creeps down a flat potential;
with the reduced Planck mass $M_{\rm pl}$, the parameters $\epsilon=\tfrac12M_{\rm pl}^2(V'/V)^2$ and $\eta_V=M_{\rm pl}^2V''/V$ are small. The
fluctuations then have a nearly scale-invariant spectrum with tilt $n_s\approx1-6\epsilon+2\eta_V$, and gravitational waves with a
tensor-to-scalar ratio $r\approx16\epsilon$~\cite{planck2020x}.

\emph{Reheating and oscillons.} After inflation the inflaton oscillates about the minimum of its potential and decays into ordinary
particles, which \emph{reheats} the universe~\cite{kls}. In a broad class of models it first fragments into \emph{oscillons}, long-lived
localised lumps of the field~\cite{amin}.

\emph{The probes.} The cosmic microwave background (CMB) was released at recombination, $z\approx1100$. Its acoustic peaks are the imprint of
sound waves in the early plasma: their spacing fixes the sound horizon and their heights the densities~\cite{planck18}. The same sound
horizon, about 150\,Mpc today, is imprinted on the clustering of galaxies as \emph{baryon acoustic oscillations} (BAO), a standard ruler
for distance against redshift~\cite{desi2}.

\emph{Dark energy.} A cosmological constant has the equation of state $w=p/\rho c^2=-1$ at all times. DESI fits the
Chevallier--Polarski--Linder (CPL) form $w(a)=w_0+w_a(1-a)$~\cite{chevallier2001,linder2003}. A \emph{thawing} dark energy has $w_0>-1$
and $w_a<0$: it was close to $-1$ early and weakens today.

\emph{Data names and tensions.} P-ACT combines the CMB spectra of Planck and of the Atacama Cosmology Telescope (ACT, data release 6);
P-ACT-LB adds CMB lensing (L) and DESI BAO (B)~\cite{louis2025}. The \emph{$H_0$ tension} is the disagreement between the Hubble constant
inferred from the CMB with $\Lambda$CDM, near 67\,km\,s$^{-1}$\,Mpc$^{-1}$, and the value from a distance ladder of Cepheids and supernovae,
near 73~\cite{planck18,riess2022}. A ladder calibrated on the tip of the red-giant branch lies in between~\cite{freedman2025}.
\end{primer}

\begin{primer}{Cosmology}{box:cosmo}
\emph{The expanding universe.} On large scales the universe is homogeneous and isotropic. Distances grow with a scale factor $a(t)$, light
from the past is redshifted by $1+z=a_0/a$, and the expansion rate $H=\dot a/a$ obeys, for a flat universe,
$H^2=8\pi G\rho/3+\Lambda c^2/3$. Each component's share is $\Omega_i=\rho_i/\rho_c$ with the critical density $\rho_c=3H^2/8\pi G$; the
value today is $H_0$.

\emph{$\Lambda$CDM.} The standard model of cosmology adds to the known particles a cold, collisionless \emph{dark matter} and a
cosmological constant $\Lambda$. Six parameters fix it: the densities of baryons and dark matter ($\Omega_bh^2$, $\Omega_ch^2$), the angular
size of the sound horizon, the optical depth to reionisation, and the amplitude $A_s$ and tilt $n_s$ of the primordial
fluctuations~\cite{planck18}.

\emph{Inflation.} An early epoch of nearly exponential expansion, $a\propto e^{Ht}$, flattens and smooths the universe and stretches
quantum fluctuations into the seeds of galaxies~\cite{guth}. It is driven by a scalar field, the \emph{inflaton}, with potential $V$, and
its length is counted in \emph{e-folds}, $N=\ln(a_{\rm end}/a_{\rm start})$. In \emph{slow roll} the field creeps down a flat potential;
with the reduced Planck mass $M_{\rm pl}$, the parameters $\epsilon=\tfrac12M_{\rm pl}^2(V'/V)^2$ and $\eta_V=M_{\rm pl}^2V''/V$ are small. The
fluctuations then have a nearly scale-invariant spectrum with tilt $n_s\approx1-6\epsilon+2\eta_V$, and gravitational waves with a
tensor-to-scalar ratio $r\approx16\epsilon$~\cite{planck2020x}.

\emph{Reheating and oscillons.} After inflation the inflaton oscillates about the minimum of its potential and decays into ordinary
particles, which \emph{reheats} the universe~\cite{kls}. In a broad class of models it first fragments into \emph{oscillons}, long-lived
localised lumps of the field~\cite{amin}.

\emph{The probes.} The cosmic microwave background (CMB) was released at recombination, $z\approx1100$. Its acoustic peaks are the imprint of
sound waves in the early plasma: their spacing fixes the sound horizon and their heights the densities~\cite{planck18}. The same sound
horizon, about 150\,Mpc today, is imprinted on the clustering of galaxies as \emph{baryon acoustic oscillations} (BAO), a standard ruler
for distance against redshift~\cite{desi2}.

\emph{Dark energy.} A cosmological constant has the equation of state $w=p/\rho c^2=-1$ at all times. DESI fits the
Chevallier--Polarski--Linder (CPL) form $w(a)=w_0+w_a(1-a)$~\cite{chevallier2001,linder2003}. A \emph{thawing} dark energy has $w_0>-1$
and $w_a<0$: it was close to $-1$ early and weakens today.

\emph{Data names and tensions.} P-ACT combines the CMB spectra of Planck and of the Atacama Cosmology Telescope (ACT, data release 6);
P-ACT-LB adds CMB lensing (L) and DESI BAO (B)~\cite{louis2025}. The \emph{$H_0$ tension} is the disagreement between the Hubble constant
inferred from the CMB with $\Lambda$CDM, near 67\,km\,s$^{-1}$\,Mpc$^{-1}$, and the value from a distance ladder of Cepheids and supernovae,
near 73~\cite{planck18,riess2022}. A ladder calibrated on the tip of the red-giant branch lies in between~\cite{freedman2025}.
\end{primer}

\begin{primer}{Cosmology}{box:cosmo}
\emph{The expanding universe.} On large scales the universe is homogeneous and isotropic. Distances grow with a scale factor $a(t)$, light
from the past is redshifted by $1+z=a_0/a$, and the expansion rate $H=\dot a/a$ obeys, for a flat universe,
$H^2=8\pi G\rho/3+\Lambda c^2/3$. Each component's share is $\Omega_i=\rho_i/\rho_c$ with the critical density $\rho_c=3H^2/8\pi G$; the
value today is $H_0$.

\emph{$\Lambda$CDM.} The standard model of cosmology adds to the known particles a cold, collisionless \emph{dark matter} and a
cosmological constant $\Lambda$. Six parameters fix it: the densities of baryons and dark matter ($\Omega_bh^2$, $\Omega_ch^2$), the angular
size of the sound horizon, the optical depth to reionisation, and the amplitude $A_s$ and tilt $n_s$ of the primordial
fluctuations~\cite{planck18}.

\emph{Inflation.} An early epoch of nearly exponential expansion, $a\propto e^{Ht}$, flattens and smooths the universe and stretches
quantum fluctuations into the seeds of galaxies~\cite{guth}. It is driven by a scalar field, the \emph{inflaton}, with potential $V$, and
its length is counted in \emph{e-folds}, $N=\ln(a_{\rm end}/a_{\rm start})$. In \emph{slow roll} the field creeps down a flat potential;
with the reduced Planck mass $M_{\rm pl}$, the parameters $\epsilon=\tfrac12M_{\rm pl}^2(V'/V)^2$ and $\eta_V=M_{\rm pl}^2V''/V$ are small. The
fluctuations then have a nearly scale-invariant spectrum with tilt $n_s\approx1-6\epsilon+2\eta_V$, and gravitational waves with a
tensor-to-scalar ratio $r\approx16\epsilon$~\cite{planck2020x}.

\emph{Reheating and oscillons.} After inflation the inflaton oscillates about the minimum of its potential and decays into ordinary
particles, which \emph{reheats} the universe~\cite{kls}. In a broad class of models it first fragments into \emph{oscillons}, long-lived
localised lumps of the field~\cite{amin}.

\emph{The probes.} The cosmic microwave background (CMB) was released at recombination, $z\approx1100$. Its acoustic peaks are the imprint of
sound waves in the early plasma: their spacing fixes the sound horizon and their heights the densities~\cite{planck18}. The same sound
horizon, about 150\,Mpc today, is imprinted on the clustering of galaxies as \emph{baryon acoustic oscillations} (BAO), a standard ruler
for distance against redshift~\cite{desi2}.

\emph{Dark energy.} A cosmological constant has the equation of state $w=p/\rho c^2=-1$ at all times. DESI fits the
Chevallier--Polarski--Linder (CPL) form $w(a)=w_0+w_a(1-a)$~\cite{chevallier2001,linder2003}. A \emph{thawing} dark energy has $w_0>-1$
and $w_a<0$: it was close to $-1$ early and weakens today.

\emph{Data names and tensions.} P-ACT combines the CMB spectra of Planck and of the Atacama Cosmology Telescope (ACT, data release 6);
P-ACT-LB adds CMB lensing (L) and DESI BAO (B)~\cite{louis2025}. The \emph{$H_0$ tension} is the disagreement between the Hubble constant
inferred from the CMB with $\Lambda$CDM, near 67\,km\,s$^{-1}$\,Mpc$^{-1}$, and the value from a distance ladder of Cepheids and supernovae,
near 73~\cite{planck18,riess2022}. A ladder calibrated on the tip of the red-giant branch lies in between~\cite{freedman2025}.
\end{primer}
}
\newcommand{\primerGalaxy}{% GUID v6, Primer P7 (label box:galaxy): galaxy dynamics. Background only; it states no result of the paper.
% Place it at the start of Section sec:sparc (or at Section sec:darkmatter if the lead prefers): % GUID v6, Primer P7 (label box:galaxy): galaxy dynamics. Background only; it states no result of the paper.
% Place it at the start of Section sec:sparc (or at Section sec:darkmatter if the lead prefers): % GUID v6, Primer P7 (label box:galaxy): galaxy dynamics. Background only; it states no result of the paper.
% Place it at the start of Section sec:sparc (or at Section sec:darkmatter if the lead prefers): \input{primers/p7-galaxy}
\begin{primer}{Galaxy dynamics}{box:galaxy}
\setlength{\parskip}{3pt plus 1pt}%
\emph{Rotation curves.} Gas and stars in a disc galaxy move on nearly circular orbits. Their speed $V$ against radius $R$ is the
\emph{rotation curve}, and the centripetal acceleration is $g_{\rm obs}=V^2/R$. Far out, most curves stay flat where the visible mass alone
would make them fall as $R^{-1/2}$: the missing-mass problem.

\emph{Baryonic mass models.} The visible, or \emph{baryonic}, mass has three parts, each with its own contribution to $V^2$, computed from
its observed surface density: the gas (atomic hydrogen with its helium), the stellar disc and the central bulge. Starlight becomes stellar
mass through a mass-to-light ratio $\Upsilon$; at a wavelength of 3.6\,\textmu m it is near 0.5 for discs and 0.7 for
bulges~\cite{mcgaugh2016}. The sum of the three, divided by $R$, is the baryonic acceleration $g_{\rm bar}$.

\emph{Two empirical laws.} The \emph{baryonic Tully--Fisher relation} (BTFR): a galaxy's baryonic mass grows as the fourth power of its flat
rotation speed, $M_{\rm bar}\propto V_{\rm flat}^4$, with small scatter~\cite{lelli2016a}. The \emph{radial acceleration relation} (RAR): at
every radius of every galaxy, $g_{\rm obs}$ is one function of $g_{\rm bar}$. It is close to $g_{\rm bar}$ at high acceleration and close to
$\sqrt{g_{\rm bar}\,g_\dagger}$ below a scale $g_\dagger\approx10^{-10}$\,m\,s$^{-2}$~\cite{mcgaugh2016,lelli2017}.

\emph{MOND.} In Milgrom's modified Newtonian dynamics, $g_{\rm obs}=\nu(g_{\rm bar}/a_0)\,g_{\rm bar}$, where the new constant
$a_0$ is about $1.2\times10^{-10}$\,m\,s$^{-2}$~\cite{milgrom}. The \emph{interpolating function} $\nu(y)$ tends to 1 at high acceleration
($y\gg1$) and to $y^{-1/2}$ at low; the deep regime gives $V^4=GMa_0$, the BTFR. Common choices are the ``simple'' function
$\nu(y)=\tfrac12+(\tfrac14+1/y)^{1/2}$ and the ``standard'' one, and the RAR's fitting function is another~\cite{famaey2012}.

\emph{Dark-matter halos.} In cold dark matter each galaxy sits in a halo. Simulated halos follow the NFW profile
\[
\rho(r)=\frac{\rho_s}{(r/r_s)(1+r/r_s)^2}
\]
of Navarro, Frenk and White~\cite{navarro1996}. Its mass $M_{200}$ is counted inside the radius $r_{200}$ that encloses 200 times the
critical density, and its concentration $c=r_{200}/r_s$ falls slowly with mass~\cite{dutton2014}. At the centre the density rises as
$r^{-1}$, a \emph{cusp}. Many dwarf and low-surface-brightness galaxies prefer a \emph{core} of nearly constant density, as in the Burkert
profile $\rho=\rho_0r_0^3/[(r+r_0)(r^2+r_0^2)]$~\cite{burkert1995}; this mismatch is the cusp-core problem~\cite{deblok2010}. The
\emph{stellar-to-halo relation}, found by matching the abundances of galaxies and halos, gives the stellar mass a halo of given mass hosts;
the ratio peaks at a few per cent near $10^{12}$ solar masses~\cite{moster2013}.

\emph{SPARC.} The Spitzer Photometry and Accurate Rotation Curves sample holds 175 disc galaxies with 3.6\,\textmu m photometry and
rotation curves from atomic-hydrogen and H$\alpha$ data, with mass models for gas, disc and bulge~\cite{lelli2016}.
\end{primer}

\begin{primer}{Galaxy dynamics}{box:galaxy}
\setlength{\parskip}{3pt plus 1pt}%
\emph{Rotation curves.} Gas and stars in a disc galaxy move on nearly circular orbits. Their speed $V$ against radius $R$ is the
\emph{rotation curve}, and the centripetal acceleration is $g_{\rm obs}=V^2/R$. Far out, most curves stay flat where the visible mass alone
would make them fall as $R^{-1/2}$: the missing-mass problem.

\emph{Baryonic mass models.} The visible, or \emph{baryonic}, mass has three parts, each with its own contribution to $V^2$, computed from
its observed surface density: the gas (atomic hydrogen with its helium), the stellar disc and the central bulge. Starlight becomes stellar
mass through a mass-to-light ratio $\Upsilon$; at a wavelength of 3.6\,\textmu m it is near 0.5 for discs and 0.7 for
bulges~\cite{mcgaugh2016}. The sum of the three, divided by $R$, is the baryonic acceleration $g_{\rm bar}$.

\emph{Two empirical laws.} The \emph{baryonic Tully--Fisher relation} (BTFR): a galaxy's baryonic mass grows as the fourth power of its flat
rotation speed, $M_{\rm bar}\propto V_{\rm flat}^4$, with small scatter~\cite{lelli2016a}. The \emph{radial acceleration relation} (RAR): at
every radius of every galaxy, $g_{\rm obs}$ is one function of $g_{\rm bar}$. It is close to $g_{\rm bar}$ at high acceleration and close to
$\sqrt{g_{\rm bar}\,g_\dagger}$ below a scale $g_\dagger\approx10^{-10}$\,m\,s$^{-2}$~\cite{mcgaugh2016,lelli2017}.

\emph{MOND.} In Milgrom's modified Newtonian dynamics, $g_{\rm obs}=\nu(g_{\rm bar}/a_0)\,g_{\rm bar}$, where the new constant
$a_0$ is about $1.2\times10^{-10}$\,m\,s$^{-2}$~\cite{milgrom}. The \emph{interpolating function} $\nu(y)$ tends to 1 at high acceleration
($y\gg1$) and to $y^{-1/2}$ at low; the deep regime gives $V^4=GMa_0$, the BTFR. Common choices are the ``simple'' function
$\nu(y)=\tfrac12+(\tfrac14+1/y)^{1/2}$ and the ``standard'' one, and the RAR's fitting function is another~\cite{famaey2012}.

\emph{Dark-matter halos.} In cold dark matter each galaxy sits in a halo. Simulated halos follow the NFW profile
\[
\rho(r)=\frac{\rho_s}{(r/r_s)(1+r/r_s)^2}
\]
of Navarro, Frenk and White~\cite{navarro1996}. Its mass $M_{200}$ is counted inside the radius $r_{200}$ that encloses 200 times the
critical density, and its concentration $c=r_{200}/r_s$ falls slowly with mass~\cite{dutton2014}. At the centre the density rises as
$r^{-1}$, a \emph{cusp}. Many dwarf and low-surface-brightness galaxies prefer a \emph{core} of nearly constant density, as in the Burkert
profile $\rho=\rho_0r_0^3/[(r+r_0)(r^2+r_0^2)]$~\cite{burkert1995}; this mismatch is the cusp-core problem~\cite{deblok2010}. The
\emph{stellar-to-halo relation}, found by matching the abundances of galaxies and halos, gives the stellar mass a halo of given mass hosts;
the ratio peaks at a few per cent near $10^{12}$ solar masses~\cite{moster2013}.

\emph{SPARC.} The Spitzer Photometry and Accurate Rotation Curves sample holds 175 disc galaxies with 3.6\,\textmu m photometry and
rotation curves from atomic-hydrogen and H$\alpha$ data, with mass models for gas, disc and bulge~\cite{lelli2016}.
\end{primer}

\begin{primer}{Galaxy dynamics}{box:galaxy}
\setlength{\parskip}{3pt plus 1pt}%
\emph{Rotation curves.} Gas and stars in a disc galaxy move on nearly circular orbits. Their speed $V$ against radius $R$ is the
\emph{rotation curve}, and the centripetal acceleration is $g_{\rm obs}=V^2/R$. Far out, most curves stay flat where the visible mass alone
would make them fall as $R^{-1/2}$: the missing-mass problem.

\emph{Baryonic mass models.} The visible, or \emph{baryonic}, mass has three parts, each with its own contribution to $V^2$, computed from
its observed surface density: the gas (atomic hydrogen with its helium), the stellar disc and the central bulge. Starlight becomes stellar
mass through a mass-to-light ratio $\Upsilon$; at a wavelength of 3.6\,\textmu m it is near 0.5 for discs and 0.7 for
bulges~\cite{mcgaugh2016}. The sum of the three, divided by $R$, is the baryonic acceleration $g_{\rm bar}$.

\emph{Two empirical laws.} The \emph{baryonic Tully--Fisher relation} (BTFR): a galaxy's baryonic mass grows as the fourth power of its flat
rotation speed, $M_{\rm bar}\propto V_{\rm flat}^4$, with small scatter~\cite{lelli2016a}. The \emph{radial acceleration relation} (RAR): at
every radius of every galaxy, $g_{\rm obs}$ is one function of $g_{\rm bar}$. It is close to $g_{\rm bar}$ at high acceleration and close to
$\sqrt{g_{\rm bar}\,g_\dagger}$ below a scale $g_\dagger\approx10^{-10}$\,m\,s$^{-2}$~\cite{mcgaugh2016,lelli2017}.

\emph{MOND.} In Milgrom's modified Newtonian dynamics, $g_{\rm obs}=\nu(g_{\rm bar}/a_0)\,g_{\rm bar}$, where the new constant
$a_0$ is about $1.2\times10^{-10}$\,m\,s$^{-2}$~\cite{milgrom}. The \emph{interpolating function} $\nu(y)$ tends to 1 at high acceleration
($y\gg1$) and to $y^{-1/2}$ at low; the deep regime gives $V^4=GMa_0$, the BTFR. Common choices are the ``simple'' function
$\nu(y)=\tfrac12+(\tfrac14+1/y)^{1/2}$ and the ``standard'' one, and the RAR's fitting function is another~\cite{famaey2012}.

\emph{Dark-matter halos.} In cold dark matter each galaxy sits in a halo. Simulated halos follow the NFW profile
\[
\rho(r)=\frac{\rho_s}{(r/r_s)(1+r/r_s)^2}
\]
of Navarro, Frenk and White~\cite{navarro1996}. Its mass $M_{200}$ is counted inside the radius $r_{200}$ that encloses 200 times the
critical density, and its concentration $c=r_{200}/r_s$ falls slowly with mass~\cite{dutton2014}. At the centre the density rises as
$r^{-1}$, a \emph{cusp}. Many dwarf and low-surface-brightness galaxies prefer a \emph{core} of nearly constant density, as in the Burkert
profile $\rho=\rho_0r_0^3/[(r+r_0)(r^2+r_0^2)]$~\cite{burkert1995}; this mismatch is the cusp-core problem~\cite{deblok2010}. The
\emph{stellar-to-halo relation}, found by matching the abundances of galaxies and halos, gives the stellar mass a halo of given mass hosts;
the ratio peaks at a few per cent near $10^{12}$ solar masses~\cite{moster2013}.

\emph{SPARC.} The Spitzer Photometry and Accurate Rotation Curves sample holds 175 disc galaxies with 3.6\,\textmu m photometry and
rotation curves from atomic-hydrogen and H$\alpha$ data, with mass models for gas, disc and bulge~\cite{lelli2016}.
\end{primer}
}
\newcommand{\primerStats}{% GUID v6, Primer P8 (label box:stats): the statistics behind the scoring rules of Section sec:scoring. Background only; it states
% no result of the paper and does not repeat the rules themselves. Place it in Section sec:scoring: % GUID v6, Primer P8 (label box:stats): the statistics behind the scoring rules of Section sec:scoring. Background only; it states
% no result of the paper and does not repeat the rules themselves. Place it in Section sec:scoring: % GUID v6, Primer P8 (label box:stats): the statistics behind the scoring rules of Section sec:scoring. Background only; it states
% no result of the paper and does not repeat the rules themselves. Place it in Section sec:scoring: \input{primers/p8-stats}
\begin{primer}{Statistics}{box:stats}
\emph{Pulls and 95\%.} A measurement $x\pm\sigma$ usually means a Gaussian error with standard deviation $\sigma$. If a prediction
$x_G\pm\sigma_G$ is right and the two errors are independent, the difference $x-x_G$ is Gaussian with variance $\sigma^2+\sigma_G^2$, so the
\emph{pull} $p=(x-x_G)/\sqrt{\sigma^2+\sigma_G^2}$ is a standard normal number. It falls outside $\pm1.96$ with probability 5\% (two-sided),
so $|p|<1.96$ is a test at 95\%, and a published 95\% half-width $b$ corresponds to $\sigma=b/1.96$. \emph{Worked pull:} a model predicts
$10.0\pm0.3$ and an experiment finds $11.0\pm0.4$. Then $p=1.0/\sqrt{0.4^2+0.3^2}=1.0/0.5=2.0$: a FAIL at 95\%, although a right model
gives $|p|\ge2.0$ in 4.6\% of independent tests. For $n$ independent tests, about $0.05\,n$ FAILs are expected by chance, and the sum of the
squared pulls follows a $\chi^2$ distribution with $n$ degrees of freedom.

\emph{BIC and $\Delta$BIC.} A model with $k$ fitted parameters that reaches the maximum likelihood $\mathcal L_{\max}$ on $N$ data points has
$\mathrm{BIC}=-2\ln\mathcal L_{\max}+k\ln N$~\cite{schwarz1978}. The penalty $k\ln N$ charges each extra parameter. For large $N$ and
independent points, the difference $\Delta$BIC of two models approximates $-2$ times the logarithm of their Bayes factor, the ratio of the
probabilities of the data under the two models. A $\Delta$BIC of 2 to 6 is then positive evidence, 6 to 10 strong and above 10 very
strong~\cite{kassraftery1995}. The approximation assumes $N$ independent points. If the points come in groups that share an error, such
as the distance of one galaxy, the effective number of independent data is closer to the number of groups. A point-wise $\Delta$BIC then
overstates the evidence, roughly in proportion to the number of points per group.

\emph{The bootstrap.} To learn how much a fitted value would vary in a new sample, one draws many new samples of the same size from the
data themselves, with replacement, and refits each~\cite{efron1979}. The spread of the refitted values estimates the uncertainty; the 16th
to 84th percentiles bracket a range comparable to $\pm1\sigma$. The units resampled must be the independent ones, for rotation curves whole
galaxies, not single points.

\emph{Many tests.} With $m$ tests at 5\% each, the chance of at least one false positive, the \emph{family-wise error rate}, grows towards
$1-0.95^m$. Bonferroni's remedy tests each at $0.05/m$. Holm's method orders the $p$-values, $p_{(1)}\le\dots\le p_{(m)}$, and rejects
$p_{(i)}$ while $p_{(i)}\le0.05/(m-i+1)$; it keeps the family-wise error rate at 5\% under any dependence and never rejects less than
Bonferroni~\cite{holm1979}. Fisher's method combines $n$ \emph{independent} $p$-values: under the null each $p_i$ is uniform, so
$-2\ln p_i$ follows $\chi^2$ with 2 degrees of freedom, and $X^2=-2\sum_i\ln p_i$ follows $\chi^2$ with $2n$~\cite{fisher1932}. It responds
strongly to a few very small $p_i$.

\emph{Look-elsewhere and preregistration.} A search over many masses, scales or candidate values finds a small \emph{local} $p$-value
somewhere by chance. The \emph{global} $p$-value, the chance of so small a value anywhere in the search, is larger, roughly the local one
times a trials factor~\cite{grossvitells2010}. A test is confirmatory only if its hypothesis, data, statistic and decision rule were
fixed before the outcome was seen; \emph{preregistration} writes them down in a dated record beforehand, and every choice made afterwards
is then visible as exploratory~\cite{nosek2018}.
\end{primer}

\begin{primer}{Statistics}{box:stats}
\emph{Pulls and 95\%.} A measurement $x\pm\sigma$ usually means a Gaussian error with standard deviation $\sigma$. If a prediction
$x_G\pm\sigma_G$ is right and the two errors are independent, the difference $x-x_G$ is Gaussian with variance $\sigma^2+\sigma_G^2$, so the
\emph{pull} $p=(x-x_G)/\sqrt{\sigma^2+\sigma_G^2}$ is a standard normal number. It falls outside $\pm1.96$ with probability 5\% (two-sided),
so $|p|<1.96$ is a test at 95\%, and a published 95\% half-width $b$ corresponds to $\sigma=b/1.96$. \emph{Worked pull:} a model predicts
$10.0\pm0.3$ and an experiment finds $11.0\pm0.4$. Then $p=1.0/\sqrt{0.4^2+0.3^2}=1.0/0.5=2.0$: a FAIL at 95\%, although a right model
gives $|p|\ge2.0$ in 4.6\% of independent tests. For $n$ independent tests, about $0.05\,n$ FAILs are expected by chance, and the sum of the
squared pulls follows a $\chi^2$ distribution with $n$ degrees of freedom.

\emph{BIC and $\Delta$BIC.} A model with $k$ fitted parameters that reaches the maximum likelihood $\mathcal L_{\max}$ on $N$ data points has
$\mathrm{BIC}=-2\ln\mathcal L_{\max}+k\ln N$~\cite{schwarz1978}. The penalty $k\ln N$ charges each extra parameter. For large $N$ and
independent points, the difference $\Delta$BIC of two models approximates $-2$ times the logarithm of their Bayes factor, the ratio of the
probabilities of the data under the two models. A $\Delta$BIC of 2 to 6 is then positive evidence, 6 to 10 strong and above 10 very
strong~\cite{kassraftery1995}. The approximation assumes $N$ independent points. If the points come in groups that share an error, such
as the distance of one galaxy, the effective number of independent data is closer to the number of groups. A point-wise $\Delta$BIC then
overstates the evidence, roughly in proportion to the number of points per group.

\emph{The bootstrap.} To learn how much a fitted value would vary in a new sample, one draws many new samples of the same size from the
data themselves, with replacement, and refits each~\cite{efron1979}. The spread of the refitted values estimates the uncertainty; the 16th
to 84th percentiles bracket a range comparable to $\pm1\sigma$. The units resampled must be the independent ones, for rotation curves whole
galaxies, not single points.

\emph{Many tests.} With $m$ tests at 5\% each, the chance of at least one false positive, the \emph{family-wise error rate}, grows towards
$1-0.95^m$. Bonferroni's remedy tests each at $0.05/m$. Holm's method orders the $p$-values, $p_{(1)}\le\dots\le p_{(m)}$, and rejects
$p_{(i)}$ while $p_{(i)}\le0.05/(m-i+1)$; it keeps the family-wise error rate at 5\% under any dependence and never rejects less than
Bonferroni~\cite{holm1979}. Fisher's method combines $n$ \emph{independent} $p$-values: under the null each $p_i$ is uniform, so
$-2\ln p_i$ follows $\chi^2$ with 2 degrees of freedom, and $X^2=-2\sum_i\ln p_i$ follows $\chi^2$ with $2n$~\cite{fisher1932}. It responds
strongly to a few very small $p_i$.

\emph{Look-elsewhere and preregistration.} A search over many masses, scales or candidate values finds a small \emph{local} $p$-value
somewhere by chance. The \emph{global} $p$-value, the chance of so small a value anywhere in the search, is larger, roughly the local one
times a trials factor~\cite{grossvitells2010}. A test is confirmatory only if its hypothesis, data, statistic and decision rule were
fixed before the outcome was seen; \emph{preregistration} writes them down in a dated record beforehand, and every choice made afterwards
is then visible as exploratory~\cite{nosek2018}.
\end{primer}

\begin{primer}{Statistics}{box:stats}
\emph{Pulls and 95\%.} A measurement $x\pm\sigma$ usually means a Gaussian error with standard deviation $\sigma$. If a prediction
$x_G\pm\sigma_G$ is right and the two errors are independent, the difference $x-x_G$ is Gaussian with variance $\sigma^2+\sigma_G^2$, so the
\emph{pull} $p=(x-x_G)/\sqrt{\sigma^2+\sigma_G^2}$ is a standard normal number. It falls outside $\pm1.96$ with probability 5\% (two-sided),
so $|p|<1.96$ is a test at 95\%, and a published 95\% half-width $b$ corresponds to $\sigma=b/1.96$. \emph{Worked pull:} a model predicts
$10.0\pm0.3$ and an experiment finds $11.0\pm0.4$. Then $p=1.0/\sqrt{0.4^2+0.3^2}=1.0/0.5=2.0$: a FAIL at 95\%, although a right model
gives $|p|\ge2.0$ in 4.6\% of independent tests. For $n$ independent tests, about $0.05\,n$ FAILs are expected by chance, and the sum of the
squared pulls follows a $\chi^2$ distribution with $n$ degrees of freedom.

\emph{BIC and $\Delta$BIC.} A model with $k$ fitted parameters that reaches the maximum likelihood $\mathcal L_{\max}$ on $N$ data points has
$\mathrm{BIC}=-2\ln\mathcal L_{\max}+k\ln N$~\cite{schwarz1978}. The penalty $k\ln N$ charges each extra parameter. For large $N$ and
independent points, the difference $\Delta$BIC of two models approximates $-2$ times the logarithm of their Bayes factor, the ratio of the
probabilities of the data under the two models. A $\Delta$BIC of 2 to 6 is then positive evidence, 6 to 10 strong and above 10 very
strong~\cite{kassraftery1995}. The approximation assumes $N$ independent points. If the points come in groups that share an error, such
as the distance of one galaxy, the effective number of independent data is closer to the number of groups. A point-wise $\Delta$BIC then
overstates the evidence, roughly in proportion to the number of points per group.

\emph{The bootstrap.} To learn how much a fitted value would vary in a new sample, one draws many new samples of the same size from the
data themselves, with replacement, and refits each~\cite{efron1979}. The spread of the refitted values estimates the uncertainty; the 16th
to 84th percentiles bracket a range comparable to $\pm1\sigma$. The units resampled must be the independent ones, for rotation curves whole
galaxies, not single points.

\emph{Many tests.} With $m$ tests at 5\% each, the chance of at least one false positive, the \emph{family-wise error rate}, grows towards
$1-0.95^m$. Bonferroni's remedy tests each at $0.05/m$. Holm's method orders the $p$-values, $p_{(1)}\le\dots\le p_{(m)}$, and rejects
$p_{(i)}$ while $p_{(i)}\le0.05/(m-i+1)$; it keeps the family-wise error rate at 5\% under any dependence and never rejects less than
Bonferroni~\cite{holm1979}. Fisher's method combines $n$ \emph{independent} $p$-values: under the null each $p_i$ is uniform, so
$-2\ln p_i$ follows $\chi^2$ with 2 degrees of freedom, and $X^2=-2\sum_i\ln p_i$ follows $\chi^2$ with $2n$~\cite{fisher1932}. It responds
strongly to a few very small $p_i$.

\emph{Look-elsewhere and preregistration.} A search over many masses, scales or candidate values finds a small \emph{local} $p$-value
somewhere by chance. The \emph{global} $p$-value, the chance of so small a value anywhere in the search, is larger, roughly the local one
times a trials factor~\cite{grossvitells2010}. A test is confirmatory only if its hypothesis, data, statistic and decision rule were
fixed before the outcome was seen; \emph{preregistration} writes them down in a dated record beforehand, and every choice made afterwards
is then visible as exploratory~\cite{nosek2018}.
\end{primer}
}
\newcommand{\primerParticle}{% GUID v6, Primer P9 (label box:particle): particle physics. Background only; it states no result of the paper.
% Place it before "The spin ladder" (marker %PRIMER:particle in part1.tex): % GUID v6, Primer P9 (label box:particle): particle physics. Background only; it states no result of the paper.
% Place it before "The spin ladder" (marker %PRIMER:particle in part1.tex): % GUID v6, Primer P9 (label box:particle): particle physics. Background only; it states no result of the paper.
% Place it before "The spin ladder" (marker %PRIMER:particle in part1.tex): \input{primers/p9-particle}
\begin{primer}{Particle physics}{box:particle}
\setlength{\parskip}{3pt plus 1pt}%
\emph{The Higgs mechanism.} The Standard Model is a gauge theory of $SU(3)\times SU(2)\times U(1)$: colour for the strong force, weak isospin
and hypercharge for the electroweak force~\cite{glashow,weinberg67,salam}. The Higgs field is a doublet whose potential has its minimum
away from zero, at a vacuum value $v=246$\,GeV~\cite{higgs64,be64}. In that vacuum the W and Z bosons acquire masses proportional to $v$,
while one combination of the electroweak fields, the photon, stays massless. Quarks and charged leptons get their masses from
\emph{Yukawa couplings} $y_f$ to the doublet, $m_f=y_fv/\sqrt2$. The Higgs boson therefore couples to each fermion with strength $m_f/v$
and to the W and Z with $2m_V^2/v$.

\emph{The $\kappa$ framework.} Measurements of Higgs production and decay are summarised by \emph{coupling modifiers} $\kappa_i$, which
scale each Standard Model coupling, so that a rate involving particle $i$ scales as $\kappa_i^2$; the Standard Model has every
$\kappa_i=1$~\cite{david2012}.

\emph{Schemes and running.} A coupling or mass in quantum field theory is defined only together with a \emph{renormalisation scheme} and
an energy scale $\mu$. In the $\overline{\text{MS}}$ scheme only the divergent part of each loop is subtracted, and parameters run with
$\mu$. The \emph{on-shell} scheme instead fixes masses at the poles of the propagators and $\alpha$ at zero momentum transfer. The same
quantity can differ between the two by more than its measurement error, so every comparison states its scheme~\cite{erler2017,navas2024}.
The dependence on $\mu$ obeys the renormalisation-group equations (RGEs). At one loop
$\dd\alpha_i^{-1}/\dd\ln\mu=-b_i/2\pi$, with $(b_1,b_2,b_3)=(41/10,-19/6,-7)$ in the Standard Model (hypercharge normalised as in grand
unification). The strong coupling therefore weakens at high energy, and the hypercharge coupling grows. Precision work runs at two or
three loops~\cite{buttazzo2013}.

\emph{Dirac or Majorana.} A \emph{Dirac} mass joins a left-handed neutrino to a right-handed partner and conserves lepton number $L$. A
\emph{Majorana} mass joins the neutrino to its own antiparticle and changes $L$ by two. The lowest-dimension operator of Standard Model
fields that gives neutrinos mass is Weinberg's dimension-five operator, $(c/\Lambda)\,(LH)(LH)$, with $L$ the lepton doublet; in the Higgs
vacuum it gives a Majorana mass of order $cv^2/\Lambda$, small when the new scale $\Lambda$ is large~\cite{weinberg1979}. Majorana masses
can therefore also come from the Higgs vacuum.

\emph{Neutrinoless double-beta decay} ($0\nu\beta\beta$) turns two neutrons of a nucleus into two protons and two electrons with no
neutrinos. It changes $L$ by two and can occur only if neutrinos have a Majorana mass. Through light-neutrino exchange its rate is
proportional to $m_{\beta\beta}^2$, with the effective mass $m_{\beta\beta}=|\sum_iU_{ei}^2m_i|$ set by the mixing matrix $U$ and the
masses $m_i$~\cite{navas2024}. Xenon and germanium experiments search for it~\cite{kamlandzen2025,legend2025}.

\emph{Grand unification.} A grand unified theory (GUT) embeds the three gauge groups in one simple group, the smallest being
$SU(5)$~\cite{georgiglashow}, with a single coupling above a scale $M_X$. With the Standard Model's particles alone, the three running
couplings do not meet at one point~\cite{amaldi}. A GUT has two generic consequences. Its heavy gauge bosons let protons decay, with a
lifetime of order $M_X^4/(\alpha_G^2m_p^5)$; Super-Kamiokande's bound on $p\to e^+\pi^0$ excludes the minimal $SU(5)$
model~\cite{superk}. And its breaking leaves magnetic monopoles of mass near $M_X/\alpha_G$, about one per causal volume, which would
overfill the universe unless inflation dilutes them~\cite{preskill,guth}.
\end{primer}

\begin{primer}{Particle physics}{box:particle}
\setlength{\parskip}{3pt plus 1pt}%
\emph{The Higgs mechanism.} The Standard Model is a gauge theory of $SU(3)\times SU(2)\times U(1)$: colour for the strong force, weak isospin
and hypercharge for the electroweak force~\cite{glashow,weinberg67,salam}. The Higgs field is a doublet whose potential has its minimum
away from zero, at a vacuum value $v=246$\,GeV~\cite{higgs64,be64}. In that vacuum the W and Z bosons acquire masses proportional to $v$,
while one combination of the electroweak fields, the photon, stays massless. Quarks and charged leptons get their masses from
\emph{Yukawa couplings} $y_f$ to the doublet, $m_f=y_fv/\sqrt2$. The Higgs boson therefore couples to each fermion with strength $m_f/v$
and to the W and Z with $2m_V^2/v$.

\emph{The $\kappa$ framework.} Measurements of Higgs production and decay are summarised by \emph{coupling modifiers} $\kappa_i$, which
scale each Standard Model coupling, so that a rate involving particle $i$ scales as $\kappa_i^2$; the Standard Model has every
$\kappa_i=1$~\cite{david2012}.

\emph{Schemes and running.} A coupling or mass in quantum field theory is defined only together with a \emph{renormalisation scheme} and
an energy scale $\mu$. In the $\overline{\text{MS}}$ scheme only the divergent part of each loop is subtracted, and parameters run with
$\mu$. The \emph{on-shell} scheme instead fixes masses at the poles of the propagators and $\alpha$ at zero momentum transfer. The same
quantity can differ between the two by more than its measurement error, so every comparison states its scheme~\cite{erler2017,navas2024}.
The dependence on $\mu$ obeys the renormalisation-group equations (RGEs). At one loop
$\dd\alpha_i^{-1}/\dd\ln\mu=-b_i/2\pi$, with $(b_1,b_2,b_3)=(41/10,-19/6,-7)$ in the Standard Model (hypercharge normalised as in grand
unification). The strong coupling therefore weakens at high energy, and the hypercharge coupling grows. Precision work runs at two or
three loops~\cite{buttazzo2013}.

\emph{Dirac or Majorana.} A \emph{Dirac} mass joins a left-handed neutrino to a right-handed partner and conserves lepton number $L$. A
\emph{Majorana} mass joins the neutrino to its own antiparticle and changes $L$ by two. The lowest-dimension operator of Standard Model
fields that gives neutrinos mass is Weinberg's dimension-five operator, $(c/\Lambda)\,(LH)(LH)$, with $L$ the lepton doublet; in the Higgs
vacuum it gives a Majorana mass of order $cv^2/\Lambda$, small when the new scale $\Lambda$ is large~\cite{weinberg1979}. Majorana masses
can therefore also come from the Higgs vacuum.

\emph{Neutrinoless double-beta decay} ($0\nu\beta\beta$) turns two neutrons of a nucleus into two protons and two electrons with no
neutrinos. It changes $L$ by two and can occur only if neutrinos have a Majorana mass. Through light-neutrino exchange its rate is
proportional to $m_{\beta\beta}^2$, with the effective mass $m_{\beta\beta}=|\sum_iU_{ei}^2m_i|$ set by the mixing matrix $U$ and the
masses $m_i$~\cite{navas2024}. Xenon and germanium experiments search for it~\cite{kamlandzen2025,legend2025}.

\emph{Grand unification.} A grand unified theory (GUT) embeds the three gauge groups in one simple group, the smallest being
$SU(5)$~\cite{georgiglashow}, with a single coupling above a scale $M_X$. With the Standard Model's particles alone, the three running
couplings do not meet at one point~\cite{amaldi}. A GUT has two generic consequences. Its heavy gauge bosons let protons decay, with a
lifetime of order $M_X^4/(\alpha_G^2m_p^5)$; Super-Kamiokande's bound on $p\to e^+\pi^0$ excludes the minimal $SU(5)$
model~\cite{superk}. And its breaking leaves magnetic monopoles of mass near $M_X/\alpha_G$, about one per causal volume, which would
overfill the universe unless inflation dilutes them~\cite{preskill,guth}.
\end{primer}

\begin{primer}{Particle physics}{box:particle}
\setlength{\parskip}{3pt plus 1pt}%
\emph{The Higgs mechanism.} The Standard Model is a gauge theory of $SU(3)\times SU(2)\times U(1)$: colour for the strong force, weak isospin
and hypercharge for the electroweak force~\cite{glashow,weinberg67,salam}. The Higgs field is a doublet whose potential has its minimum
away from zero, at a vacuum value $v=246$\,GeV~\cite{higgs64,be64}. In that vacuum the W and Z bosons acquire masses proportional to $v$,
while one combination of the electroweak fields, the photon, stays massless. Quarks and charged leptons get their masses from
\emph{Yukawa couplings} $y_f$ to the doublet, $m_f=y_fv/\sqrt2$. The Higgs boson therefore couples to each fermion with strength $m_f/v$
and to the W and Z with $2m_V^2/v$.

\emph{The $\kappa$ framework.} Measurements of Higgs production and decay are summarised by \emph{coupling modifiers} $\kappa_i$, which
scale each Standard Model coupling, so that a rate involving particle $i$ scales as $\kappa_i^2$; the Standard Model has every
$\kappa_i=1$~\cite{david2012}.

\emph{Schemes and running.} A coupling or mass in quantum field theory is defined only together with a \emph{renormalisation scheme} and
an energy scale $\mu$. In the $\overline{\text{MS}}$ scheme only the divergent part of each loop is subtracted, and parameters run with
$\mu$. The \emph{on-shell} scheme instead fixes masses at the poles of the propagators and $\alpha$ at zero momentum transfer. The same
quantity can differ between the two by more than its measurement error, so every comparison states its scheme~\cite{erler2017,navas2024}.
The dependence on $\mu$ obeys the renormalisation-group equations (RGEs). At one loop
$\dd\alpha_i^{-1}/\dd\ln\mu=-b_i/2\pi$, with $(b_1,b_2,b_3)=(41/10,-19/6,-7)$ in the Standard Model (hypercharge normalised as in grand
unification). The strong coupling therefore weakens at high energy, and the hypercharge coupling grows. Precision work runs at two or
three loops~\cite{buttazzo2013}.

\emph{Dirac or Majorana.} A \emph{Dirac} mass joins a left-handed neutrino to a right-handed partner and conserves lepton number $L$. A
\emph{Majorana} mass joins the neutrino to its own antiparticle and changes $L$ by two. The lowest-dimension operator of Standard Model
fields that gives neutrinos mass is Weinberg's dimension-five operator, $(c/\Lambda)\,(LH)(LH)$, with $L$ the lepton doublet; in the Higgs
vacuum it gives a Majorana mass of order $cv^2/\Lambda$, small when the new scale $\Lambda$ is large~\cite{weinberg1979}. Majorana masses
can therefore also come from the Higgs vacuum.

\emph{Neutrinoless double-beta decay} ($0\nu\beta\beta$) turns two neutrons of a nucleus into two protons and two electrons with no
neutrinos. It changes $L$ by two and can occur only if neutrinos have a Majorana mass. Through light-neutrino exchange its rate is
proportional to $m_{\beta\beta}^2$, with the effective mass $m_{\beta\beta}=|\sum_iU_{ei}^2m_i|$ set by the mixing matrix $U$ and the
masses $m_i$~\cite{navas2024}. Xenon and germanium experiments search for it~\cite{kamlandzen2025,legend2025}.

\emph{Grand unification.} A grand unified theory (GUT) embeds the three gauge groups in one simple group, the smallest being
$SU(5)$~\cite{georgiglashow}, with a single coupling above a scale $M_X$. With the Standard Model's particles alone, the three running
couplings do not meet at one point~\cite{amaldi}. A GUT has two generic consequences. Its heavy gauge bosons let protons decay, with a
lifetime of order $M_X^4/(\alpha_G^2m_p^5)$; Super-Kamiokande's bound on $p\to e^+\pi^0$ excludes the minimal $SU(5)$
model~\cite{superk}. And its breaking leaves magnetic monopoles of mass near $M_X/\alpha_G$, about one per causal volume, which would
overfill the universe unless inflation dilutes them~\cite{preskill,guth}.
\end{primer}
}
\newcommand{\primerHadron}{% GUID v6, Primer P10 (label box:hadron): hadron spectroscopy. Background only; it states no result of the paper.
% Place it before "Hadrons as trees of dualities" (marker %PRIMER:hadron in part1.tex): % GUID v6, Primer P10 (label box:hadron): hadron spectroscopy. Background only; it states no result of the paper.
% Place it before "Hadrons as trees of dualities" (marker %PRIMER:hadron in part1.tex): % GUID v6, Primer P10 (label box:hadron): hadron spectroscopy. Background only; it states no result of the paper.
% Place it before "Hadrons as trees of dualities" (marker %PRIMER:hadron in part1.tex): \input{primers/p10-hadron}
\begin{primer}{Hadron spectroscopy}{box:hadron}
\emph{Colour.} Quarks carry the colour charge of $SU(3)$: a quark is a triplet $\mathbf3$, an antiquark an antitriplet $\mathbf{\bar3}$.
Only colour \emph{singlets} exist as free particles, the hadrons~\cite{navas2024}. The products
$\mathbf3\otimes\mathbf{\bar3}=\mathbf1\oplus\mathbf8$ and $\mathbf3\otimes\mathbf3=\mathbf{\bar3}\oplus\mathbf6$ give the ordinary
hadrons: a meson $q\bar q$ in the singlet, and a baryon $qqq$, whose first two quarks couple to $\mathbf{\bar3}$ and then with the third to
a singlet. One-gluon exchange attracts in the singlet $q\bar q$ and, half as strongly, in the $\mathbf{\bar3}$ $qq$, and repels in the
$\mathbf8$ and the $\mathbf6$.

\emph{Diquarks.} A \emph{diquark} is a pair of quarks in the attractive $\mathbf{\bar3}$. The ``good'' diquark, with spin 0 and a flavour-
antisymmetric pair such as $[ud]$, binds most; the ``bad'' one has spin 1~\cite{jaffe2005}.

\emph{Compact or molecular.} A hadron with more than three quarks, an \emph{exotic}, can be \emph{compact}, for example a diquark and an
antidiquark bound by colour forces, $[qq][\bar q\bar q]$~\cite{maiani2005}. Or it can be a \emph{hadronic molecule}: two colour-singlet
hadrons held together by the residual strong force, like the deuteron~\cite{guo2018}. Molecules sit just below the threshold
$m_1+m_2$ of their two constituents.

\emph{The colour-spin interaction.} In the constituent-quark model a hadron's mass is the sum of effective quark masses plus a hyperfine
term proportional to $(\boldsymbol\lambda_i\cdot\boldsymbol\lambda_j)(\boldsymbol\sigma_i\cdot\boldsymbol\sigma_j)/m_im_j$ for each pair,
where the $\boldsymbol\lambda$ are colour and the $\boldsymbol\sigma$ spin matrices. This chromomagnetic interaction of De R\'ujula, Georgi
and Glashow explains, for example, why the $\Delta$ is heavier than the nucleon~\cite{derujula1975}. Karliner and Rosner (KR) refined it
with separate effective quark masses in mesons and in baryons, the baryon masses for every constituent of a compact multiquark, and binding
terms for heavy pairs with $B(QQ)=B(Q\bar Q)/2$; with it they predicted the doubly charmed baryon and doubly heavy
tetraquarks~\cite{karliner2014,karliner2017}.

\emph{Thresholds and cusps.} Where a new two-hadron channel opens, at $M=m_1+m_2$, the scattering amplitude has a branch point. A
\emph{cusp} can then appear in a mass spectrum without any state, and a state that couples to the channel is pulled towards the threshold.
Many exotic candidates lie within a few MeV of a threshold~\cite{guo2018}.

\emph{Names.} The Particle Data Group names hadrons by their quantum numbers, not by their structure~\cite{navas2024}; the
$X(3872)$, for example, is listed as $\chi_{c1}(3872)$. LHCb's convention for exotics~\cite{gershon2022} uses $T$ for four-quark and $P$
for five-quark states, with subscripts for heavy flavour, $\psi$ standing for hidden charm $c\bar c$. So $T_{cc}^+$ is $cc\bar u\bar d$,
seen just below the $D^{*+}D^0$ threshold~\cite{lhcb2022h}, and the hidden-charm pentaquark $P_c(4312)^+$~\cite{lhcb2019h} becomes
$P_\psi^N(4312)^+$.

\emph{String junctions.} In a baryon the three colour flux tubes meet at a point, a \emph{string junction}~\cite{rossiveneziano1977}. A
compact tetraquark has two junctions and a compact pentaquark three, and some models add a constant energy per
junction~\cite{karliner2017a}.
\end{primer}

\begin{primer}{Hadron spectroscopy}{box:hadron}
\emph{Colour.} Quarks carry the colour charge of $SU(3)$: a quark is a triplet $\mathbf3$, an antiquark an antitriplet $\mathbf{\bar3}$.
Only colour \emph{singlets} exist as free particles, the hadrons~\cite{navas2024}. The products
$\mathbf3\otimes\mathbf{\bar3}=\mathbf1\oplus\mathbf8$ and $\mathbf3\otimes\mathbf3=\mathbf{\bar3}\oplus\mathbf6$ give the ordinary
hadrons: a meson $q\bar q$ in the singlet, and a baryon $qqq$, whose first two quarks couple to $\mathbf{\bar3}$ and then with the third to
a singlet. One-gluon exchange attracts in the singlet $q\bar q$ and, half as strongly, in the $\mathbf{\bar3}$ $qq$, and repels in the
$\mathbf8$ and the $\mathbf6$.

\emph{Diquarks.} A \emph{diquark} is a pair of quarks in the attractive $\mathbf{\bar3}$. The ``good'' diquark, with spin 0 and a flavour-
antisymmetric pair such as $[ud]$, binds most; the ``bad'' one has spin 1~\cite{jaffe2005}.

\emph{Compact or molecular.} A hadron with more than three quarks, an \emph{exotic}, can be \emph{compact}, for example a diquark and an
antidiquark bound by colour forces, $[qq][\bar q\bar q]$~\cite{maiani2005}. Or it can be a \emph{hadronic molecule}: two colour-singlet
hadrons held together by the residual strong force, like the deuteron~\cite{guo2018}. Molecules sit just below the threshold
$m_1+m_2$ of their two constituents.

\emph{The colour-spin interaction.} In the constituent-quark model a hadron's mass is the sum of effective quark masses plus a hyperfine
term proportional to $(\boldsymbol\lambda_i\cdot\boldsymbol\lambda_j)(\boldsymbol\sigma_i\cdot\boldsymbol\sigma_j)/m_im_j$ for each pair,
where the $\boldsymbol\lambda$ are colour and the $\boldsymbol\sigma$ spin matrices. This chromomagnetic interaction of De R\'ujula, Georgi
and Glashow explains, for example, why the $\Delta$ is heavier than the nucleon~\cite{derujula1975}. Karliner and Rosner (KR) refined it
with separate effective quark masses in mesons and in baryons, the baryon masses for every constituent of a compact multiquark, and binding
terms for heavy pairs with $B(QQ)=B(Q\bar Q)/2$; with it they predicted the doubly charmed baryon and doubly heavy
tetraquarks~\cite{karliner2014,karliner2017}.

\emph{Thresholds and cusps.} Where a new two-hadron channel opens, at $M=m_1+m_2$, the scattering amplitude has a branch point. A
\emph{cusp} can then appear in a mass spectrum without any state, and a state that couples to the channel is pulled towards the threshold.
Many exotic candidates lie within a few MeV of a threshold~\cite{guo2018}.

\emph{Names.} The Particle Data Group names hadrons by their quantum numbers, not by their structure~\cite{navas2024}; the
$X(3872)$, for example, is listed as $\chi_{c1}(3872)$. LHCb's convention for exotics~\cite{gershon2022} uses $T$ for four-quark and $P$
for five-quark states, with subscripts for heavy flavour, $\psi$ standing for hidden charm $c\bar c$. So $T_{cc}^+$ is $cc\bar u\bar d$,
seen just below the $D^{*+}D^0$ threshold~\cite{lhcb2022h}, and the hidden-charm pentaquark $P_c(4312)^+$~\cite{lhcb2019h} becomes
$P_\psi^N(4312)^+$.

\emph{String junctions.} In a baryon the three colour flux tubes meet at a point, a \emph{string junction}~\cite{rossiveneziano1977}. A
compact tetraquark has two junctions and a compact pentaquark three, and some models add a constant energy per
junction~\cite{karliner2017a}.
\end{primer}

\begin{primer}{Hadron spectroscopy}{box:hadron}
\emph{Colour.} Quarks carry the colour charge of $SU(3)$: a quark is a triplet $\mathbf3$, an antiquark an antitriplet $\mathbf{\bar3}$.
Only colour \emph{singlets} exist as free particles, the hadrons~\cite{navas2024}. The products
$\mathbf3\otimes\mathbf{\bar3}=\mathbf1\oplus\mathbf8$ and $\mathbf3\otimes\mathbf3=\mathbf{\bar3}\oplus\mathbf6$ give the ordinary
hadrons: a meson $q\bar q$ in the singlet, and a baryon $qqq$, whose first two quarks couple to $\mathbf{\bar3}$ and then with the third to
a singlet. One-gluon exchange attracts in the singlet $q\bar q$ and, half as strongly, in the $\mathbf{\bar3}$ $qq$, and repels in the
$\mathbf8$ and the $\mathbf6$.

\emph{Diquarks.} A \emph{diquark} is a pair of quarks in the attractive $\mathbf{\bar3}$. The ``good'' diquark, with spin 0 and a flavour-
antisymmetric pair such as $[ud]$, binds most; the ``bad'' one has spin 1~\cite{jaffe2005}.

\emph{Compact or molecular.} A hadron with more than three quarks, an \emph{exotic}, can be \emph{compact}, for example a diquark and an
antidiquark bound by colour forces, $[qq][\bar q\bar q]$~\cite{maiani2005}. Or it can be a \emph{hadronic molecule}: two colour-singlet
hadrons held together by the residual strong force, like the deuteron~\cite{guo2018}. Molecules sit just below the threshold
$m_1+m_2$ of their two constituents.

\emph{The colour-spin interaction.} In the constituent-quark model a hadron's mass is the sum of effective quark masses plus a hyperfine
term proportional to $(\boldsymbol\lambda_i\cdot\boldsymbol\lambda_j)(\boldsymbol\sigma_i\cdot\boldsymbol\sigma_j)/m_im_j$ for each pair,
where the $\boldsymbol\lambda$ are colour and the $\boldsymbol\sigma$ spin matrices. This chromomagnetic interaction of De R\'ujula, Georgi
and Glashow explains, for example, why the $\Delta$ is heavier than the nucleon~\cite{derujula1975}. Karliner and Rosner (KR) refined it
with separate effective quark masses in mesons and in baryons, the baryon masses for every constituent of a compact multiquark, and binding
terms for heavy pairs with $B(QQ)=B(Q\bar Q)/2$; with it they predicted the doubly charmed baryon and doubly heavy
tetraquarks~\cite{karliner2014,karliner2017}.

\emph{Thresholds and cusps.} Where a new two-hadron channel opens, at $M=m_1+m_2$, the scattering amplitude has a branch point. A
\emph{cusp} can then appear in a mass spectrum without any state, and a state that couples to the channel is pulled towards the threshold.
Many exotic candidates lie within a few MeV of a threshold~\cite{guo2018}.

\emph{Names.} The Particle Data Group names hadrons by their quantum numbers, not by their structure~\cite{navas2024}; the
$X(3872)$, for example, is listed as $\chi_{c1}(3872)$. LHCb's convention for exotics~\cite{gershon2022} uses $T$ for four-quark and $P$
for five-quark states, with subscripts for heavy flavour, $\psi$ standing for hidden charm $c\bar c$. So $T_{cc}^+$ is $cc\bar u\bar d$,
seen just below the $D^{*+}D^0$ threshold~\cite{lhcb2022h}, and the hidden-charm pentaquark $P_c(4312)^+$~\cite{lhcb2019h} becomes
$P_\psi^N(4312)^+$.

\emph{String junctions.} In a baryon the three colour flux tubes meet at a point, a \emph{string junction}~\cite{rossiveneziano1977}. A
compact tetraquark has two junctions and a compact pentaquark three, and some models add a constant energy per
junction~\cite{karliner2017a}.
\end{primer}
}
